%=======================================================================
%  Chapter 1.  MATHEMATICAL PRELIMINARIES  (complete, with solved problems)
%  A from-scratch, motivated account of Steigmann & Shirani,
%  "Principles of Continuum Mechanics", Chapter 1 (1.1.1-1.3), matching
%  the book's section numbering and order.  Every step is shown.
%
%  Compile:  pdflatex continuum_ch1.tex   (run twice, for references)
%=======================================================================
\documentclass[17pt,a4paper]{extreport}

\usepackage[margin=1.8cm]{geometry}
\usepackage{amsmath,amssymb,amsthm}
\usepackage{bm}
\usepackage{accents}
% book style: boldface is rendered by an under-tilde
\DeclareRobustCommand{\vek}[1]{\underaccent{\tilde}{#1}}
\usepackage{xcolor}
\usepackage{enumitem}
\usepackage{tikz}
\usetikzlibrary{arrows.meta,calc,positioning,angles,quotes,patterns}
\tikzset{vec/.style={-{Stealth[length=2.3mm]},thick},
         thinvec/.style={-{Stealth[length=1.8mm]},semithick},
         guide/.style={gray!65,densely dashed,thin}}
\usepackage[colorlinks=true,linkcolor=blue!55!black,citecolor=blue!55!black,
            hypertexnames=false,pdfencoding=unicode]{hyperref}

%=============================== notation ==============================
\newcommand{\Rn}{\mathbb{R}}
\newcommand{\V}{\mathcal{V}}
\newcommand{\Vn}{\mathcal{V}_{n}}
\newcommand{\En}{\mathbb{E}^{n}}
\newcommand{\Et}{\mathbb{E}^{3}}
\newcommand{\hEt}{\hat{\mathbb{E}}^{3}}
\newcommand{\Lin}{\operatorname{Lin}}
\newcommand{\Sym}{\operatorname{Sym}}
\newcommand{\Skw}{\operatorname{Skw}}
\newcommand{\Sph}{\operatorname{Sph}}
\newcommand{\Dev}{\operatorname{Dev}}
\newcommand{\Proj}[1]{\mathbb{P}_{(#1)}}
\newcommand{\Refl}[1]{\mathbb{R}_{(#1)}}
\newcommand{\tr}{\operatorname{tr}}
\newcommand{\grad}{\nabla}
\newcommand{\dvg}{\operatorname{div}}
\newcommand{\crl}{\operatorname{curl}}
\newcommand{\I}{\vek{I}}
\newcommand{\Ob}{\vek{O}}
\newcommand{\ba}{\vek{a}}
\newcommand{\bb}{\vek{b}}
\newcommand{\bc}{\vek{c}}
\newcommand{\bd}{\vek{d}}
\newcommand{\bh}{\vek{h}}
\newcommand{\bp}{\vek{p}}
\newcommand{\bq}{\vek{q}}
\newcommand{\bu}{\vek{u}}
\newcommand{\bv}{\vek{v}}
\newcommand{\bw}{\vek{w}}
\newcommand{\bx}{\vek{x}}
\newcommand{\by}{\vek{y}}
\newcommand{\bz}{\vek{z}}
\newcommand{\bo}{\vek{o}}
\newcommand{\be}{\vek{e}}
\newcommand{\bmv}{\vek{m}}
\newcommand{\bn}{\vek{n}}
\newcommand{\btt}{\vek{t}}
\newcommand{\bk}{\vek{k}}
\newcommand{\bze}{\vek{0}}
\newcommand{\bA}{\vek{A}}
\newcommand{\bB}{\vek{B}}
\newcommand{\bC}{\vek{C}}
\newcommand{\bD}{\vek{D}}
\newcommand{\bE}{\vek{E}}
\newcommand{\bL}{\vek{L}}
\newcommand{\bP}{\vek{P}}
\newcommand{\bQ}{\vek{Q}}
\newcommand{\bR}{\vek{R}}
\newcommand{\bS}{\vek{S}}
\newcommand{\bT}{\vek{T}}
\newcommand{\bV}{\vek{V}}
\newcommand{\bW}{\vek{W}}
\newcommand{\hbE}{\hat{\vek{E}}}
\newcommand{\bOm}{\vek{\Omega}}
\newcommand{\bom}{\vek{\omega}}
\newcommand{\bF}{\vek{F}}
\newcommand{\bN}{\vek{N}}
\newcommand{\bff}{\vek{f}}
\newcommand{\Vt}{\mathbb{E}^{3}}
\newcommand{\ber}{\vek{e}_{r}}
\renewcommand{\beth}{\vek{e}_{\theta}}
\newcommand{\beph}{\vek{e}_{\varphi}}
\newcommand{\berho}{\vek{e}_{\rho}}
\newcommand{\dd}{\delta}
\newcommand{\eps}{\varepsilon}
\newcommand{\dl}{\mathrm{d}}
\newcommand{\tbox}[3]{[\,#1,#2,#3\,]}
\newcommand{\ip}[2]{#1\cdot#2}
\newcommand{\rest}[1]{\big|_{#1}}

% theorem headers in small capitals rather than bold
\newtheoremstyle{scplain}{\topsep}{\topsep}{\itshape}{0pt}
  {\scshape}{.}{5pt plus 1pt minus 1pt}{}
\newtheoremstyle{scdefn}{\topsep}{\topsep}{\normalfont}{0pt}
  {\scshape}{.}{5pt plus 1pt minus 1pt}{}

\theoremstyle{scplain}
\newtheorem{theorem}{Theorem}[chapter]
\newtheorem{proposition}[theorem]{Proposition}
\newtheorem{lemma}[theorem]{Lemma}
\newtheorem{corollary}[theorem]{Corollary}
\theoremstyle{scdefn}
\newtheorem{definition}[theorem]{Definition}
\newtheorem{remark}[theorem]{Remark}
\newtheorem{example}[theorem]{Example}

\theoremstyle{scplain}
\newtheorem{problem}{Problem}
\newenvironment{solution}
 {\par\smallskip\noindent\textit{Solution.}\ }
 {\hfill$\square$\par\medskip}

\renewcommand{\qedsymbol}{$\blacksquare$}
\allowdisplaybreaks

%=======================================================================

%------------------- clickable equation references ---------------------
% \eqr{4.11} prints "(4.11)" exactly as before and turns it into a link
% when equation (4.11) is present in this document.  When it is not -- a
% standalone chapter referring to an equation of another chapter -- it
% degrades silently to plain text instead of raising an undefined
% reference.  Must sit after hyperref.
\makeatletter
% a display carrying \tag{2.48--2.50} can hold only one \label, so the
% remaining members of the range are routed to it by this alias table
\newcommand{\eqalias}[2]{\expandafter\gdef\csname eqal@#1\endcsname{#2}}
\newcommand{\eqr@do}[1]{%
  \begingroup
    \@ifundefined{eqal@#1}{\def\eqr@t{#1}}%
                          {\edef\eqr@t{\csname eqal@#1\endcsname}}%
    \@ifundefined{r@eq:\eqr@t}%
      {\textup{(#1)}}%
      {\hyperref[eq:\eqr@t]{\textup{(#1)}}}%
  \endgroup}
\DeclareRobustCommand{\eqr}[1]{%
  \ifmmode\text{\eqr@do{#1}}\else\eqr@do{#1}\fi}
\makeatother
% generated: extra members of a range tag point at the tag itself
\eqalias{1.104}{1.103}
\eqalias{1.105}{1.103}
\eqalias{1.106}{1.103}
\eqalias{1.114}{1.113}
\eqalias{1.127}{1.126}
\eqalias{2.6}{2.5}
\eqalias{2.9}{2.8}
\eqalias{2.19}{2.18}
\eqalias{2.23}{2.22}
\eqalias{2.26}{2.25}
\eqalias{2.44}{2.43}
\eqalias{2.46}{2.45}
\eqalias{2.49}{2.48}
\eqalias{2.50}{2.48}
\eqalias{2.53}{2.52}
\eqalias{2.54}{2.52}
\eqalias{2.72}{2.71}
\eqalias{2.73}{2.71}
\eqalias{2.78}{2.77}
\eqalias{2.81}{2.80}
\eqalias{2.87}{2.86}
\eqalias{2.88}{2.86}
\eqalias{2.90}{2.89}
\eqalias{2.94}{2.93}
\eqalias{2.97}{2.96}
\eqalias{2.99}{2.98}
\eqalias{2.103}{2.102}
\eqalias{2.108}{2.107}
\eqalias{3.10}{3.9}


\begin{document}

\begin{titlepage}
\centering
\vspace*{3cm}
{\Large\bfseries Mathematical Preliminaries\par}
\vspace{4mm}
{\large Vector and tensor algebra, fields, integral theorems,\\
and full solutions to the chapter problems\par}
\vspace{8mm}
{\itshape A from-scratch, definition--theorem--proof companion to}\\[2mm]
{Steigmann and Shirani, \emph{Principles of Continuum Mechanics},
Chapter~1}\\[1mm]
{(\S\,1.1.1 through \S\,1.3, in the book's order and numbering)}
\vfill
{\small Each subsection opens with plain expository ground, in the spirit
of the text, before any formal statement. Nothing in the chapter is
skipped; every algebraic step, including those the book leaves to the
reader, is carried out. The chapter problems are all solved in \S1.3.\par}
\vspace{1cm}
\end{titlepage}

\tableofcontents

\chapter{Mathematical Preliminaries}

A working knowledge of vector and tensor algebra and analysis is a
prerequisite for continuum mechanics. This chapter assembles the results
we will need. The style is deliberately slow: we begin each topic with the
naive question that forces the next idea, then make it precise and prove
it. Throughout, $\Vn$ is an $n$ dimensional real vector space, $\En$ the
same carrying an inner product, and $\Et$ the three dimensional case;
$\{\be_i\}$ is a fixed orthonormal basis, indices run over $\{1,2,3\}$
unless stated otherwise, and a repeated index in a term is summed. The
symbol $\dd_{ij}$ is Kronecker's, $\eps_{ijk}$ the permutation symbol, and
$\Lin$ the space of second order tensors. Appendix~\ref{app:dict} maps
every numbered book equation to the place it is produced here.

%=======================================================================
\section{Vector and tensor theory}\label{sec:vtt}
%=======================================================================

\subsection{Vector spaces}\label{ss:vs}

Displacements in space, velocities, and forces share exactly two features:
any two of them may be added tip to tail, and any one may be scaled by a
real number. Nothing else about ``arrows'' is used in the algebra that
follows, so instead of drawing pictures we simply write down those two
operations and the rules they obey. The reward is leverage: a theorem
proved from the rules holds at once for displacements, forces, and later
for tensors and fields. The list of rules is the definition of a vector
space.

\begin{definition}[Vector space]\label{def:vs}
A real vector space is a set $\V$ with maps $+:\V\times\V\to\V$ and
$\cdot:\Rn\times\V\to\V$ such that, for all $\ba,\bb,\bc\in\V$ and
$\alpha,\beta\in\Rn$,
\[
\begin{array}{ll}
\text{(a)}\ \alpha\ba+\beta\bb\in\V, &
\text{(b)}\ (\ba+\bb)+\bc=\ba+(\bb+\bc), \\[1mm]
\text{(c)}\ \ba+\bb=\bb+\ba, &
\text{(d)}\ \exists\,\bze:\ \ba+\bze=\ba, \\[1mm]
\text{(e)}\ \exists\,(-\ba):\ \ba+(-\ba)=\bze, &
\text{(f)}\ (\alpha\beta)\ba=\alpha(\beta\ba), \\[1mm]
\text{(g)}\ (\alpha+\beta)\ba=\alpha\ba+\beta\ba, &
\text{(h)}\ \alpha(\ba+\bb)=\alpha\ba+\alpha\bb, \\[1mm]
\text{(i)}\ 1\ba=\ba, & \text{(j)}\ 0\ba=\bze.
\end{array}
\]
Vectors are set in bold; real numbers (scalars) in light face.
\end{definition}

We want to pin any vector down by a finite list of numbers. For that we
need a set of reference vectors that neither wastes effort duplicating each
other nor leaves gaps. ``No duplication'' is linear independence; ``no
gaps'' together with independence is a basis.

\begin{definition}[Linear dependence]\label{def:lindep}
$\{\bv_{1},\dots,\bv_{p}\}\subset\V$ is linearly independent iff
$\alpha_{i}\bv_{i}=\bze\Rightarrow\alpha_{1}=\cdots=\alpha_{p}=0$, and
linearly dependent otherwise.
\end{definition}

\begin{definition}[Dimension, basis]\label{def:dim}
$\V$ is finite dimensional of dimension $n$, written $\Vn$, iff every
linearly independent set has at most $n$ elements and some such set has
exactly $n$. Any linearly independent set of $n$ elements is a basis.
\end{definition}

\begin{theorem}[Coordinates exist and are unique]\label{thm:coord}
Let $\{\bu_{1},\dots,\bu_{n}\}$ be a basis of $\Vn$. Then every
$\bv\in\Vn$ has a representation $\bv=\alpha_{i}\bu_{i}$ with unique
coefficients $\alpha_{i}$.
\end{theorem}

\begin{proof}
\emph{Existence.} $\{\bu_{1},\dots,\bu_{n},\bv\}$ has $n+1>n$ elements,
hence is dependent: $\lambda\bv+\lambda_{i}\bu_{i}=\bze$ with the
coefficients not all zero. If $\lambda=0$ then $\lambda_{i}\bu_{i}=\bze$
forces every $\lambda_{i}=0$, a contradiction; so $\lambda\neq0$ and
$\bv=(-\lambda_{i}/\lambda)\bu_{i}$. \emph{Uniqueness.}
$\alpha_{i}\bu_{i}=\beta_{i}\bu_{i}\Rightarrow(\alpha_{i}-\beta_{i})\bu_{i}
=\bze\Rightarrow\alpha_{i}=\beta_{i}$ by independence.
\end{proof}

\medskip
The axioms can add and scale, but they cannot yet say how \emph{long} a
vector is or what \emph{angle} two of them make: those words never appear
in (a)--(j). To measure we add one operation that eats two vectors and
returns a number, is symmetric, is linear in each slot, and is strictly
positive on nonzero vectors. From it flow length, angle, orthogonality,
and all of Euclidean geometry.

\begin{definition}[Inner product, norm]\label{def:ip}
$\En$ is $\Vn$ with a map $\cdot:\En\times\En\to\Rn$ such that, for all
$\bu,\bv,\bw$ and $\alpha\in\Rn$,
\[
\text{(a)}\ \ip\bu\bv=\ip\bv\bu,\quad
\text{(b)}\ \ip\bu{(\bv+\bw)}=\ip\bu\bv+\ip\bu\bw,\quad
\text{(c)}\ \ip{(\alpha\bu)}\bv=\alpha(\ip\bu\bv),
\]
\[
\text{(d)}\ \ip\bu\bu\ge0,\qquad \ip\bu\bu=0\iff\bu=\bze.
\]
The norm is $|\bu|=(\ip\bu\bu)^{1/2}$; $\bu$ is a unit vector iff
$|\bu|=1$; $\bu\perp\bv$ iff $\ip\bu\bv=0$.
\end{definition}

Before we may define the angle by $\cos\theta=\ip\bu\bv/(|\bu||\bv|)$, the
right side must lie in $[-1,1]$, or the arccosine is meaningless. That
bound is not a new assumption: positivity already forces it, because a
particular squared length can never be negative.

\begin{theorem}[Cauchy--Schwarz]\label{thm:cs}
$|\ip\bu\bv|\le|\bu|\,|\bv|$ for all $\bu,\bv\in\En$, with equality iff
$\bu,\bv$ are parallel.
\end{theorem}

\begin{proof}
If $\bv=\bze$ both sides vanish. Otherwise positivity (d) applied to
$\bu-\lambda\bv$ gives, for every $\lambda\in\Rn$,
\[
0\le|\bu-\lambda\bv|^{2}=|\bu|^{2}-2\lambda\,\ip\bu\bv+\lambda^{2}|\bv|^{2},
\]
using bilinearity and symmetry. The right side is least at
$\lambda=\ip\bu\bv/|\bv|^{2}$; substituting,
$0\le|\bu|^{2}-(\ip\bu\bv)^{2}/|\bv|^{2}$, i.e.
$(\ip\bu\bv)^{2}\le|\bu|^{2}|\bv|^{2}$. Equality forces
$\bu-\lambda\bv=\bze$ by (d), i.e. $\bu\parallel\bv$.
\end{proof}

\begin{corollary}[Triangle inequality]\label{cor:tri}
$|\bu+\bv|\le|\bu|+|\bv|$.
\end{corollary}

\begin{proof}
$|\bu+\bv|^{2}=|\bu|^{2}+2\ip\bu\bv+|\bv|^{2}
\le|\bu|^{2}+2|\bu||\bv|+|\bv|^{2}=(|\bu|+|\bv|)^{2}$ by
Theorem~\ref{thm:cs}; take square roots.
\end{proof}

\begin{definition}[Angle]\label{def:angle}
For $\bu,\bv\neq\bze$, $\theta=\arccos(\ip\bu\bv/(|\bu||\bv|))\in[0,\pi]$,
well defined by Theorem~\ref{thm:cs}.
\end{definition}

A general basis reproduces vectors but scrambles the metric: components are
not projections, and $\ip\bu\bv$ is not $u_iv_i$. A basis of mutually
orthogonal unit vectors removes every cross term, so a component becomes a
projection and the inner product becomes the plain sum of products. This is
why all computation below uses an orthonormal basis.

\begin{definition}[Orthonormal basis]\label{def:onb}
$\{\be_{1},\dots,\be_{n}\}$ is orthonormal iff
$\ip{\be_{i}}{\be_{j}}=\dd_{ij}$.
\end{definition}

\begin{proposition}\label{prop:onlindep}
Every orthonormal set is linearly independent.
\end{proposition}

\begin{proof}
$\alpha_{i}\be_{i}=\bze\Rightarrow
0=\ip{(\alpha_{i}\be_{i})}{\be_{j}}=\alpha_{i}\dd_{ij}=\alpha_{j}$.
\end{proof}

\begin{proposition}[Components and the dot product]\label{prop:comp}
Relative to an orthonormal basis $\{\be_{i}\}$, every $\bv$ satisfies
$\bv=v_{i}\be_{i}$ with $v_{i}=\ip\bv{\be_{i}}$, and for
$\bu=u_{i}\be_{i}$, $\bv=v_{i}\be_{i}$, $\ip\bu\bv=u_{i}v_{i}$.
\end{proposition}

\begin{proof}
By Theorem~\ref{thm:coord}, $\bv=v_{i}\be_{i}$ uniquely; then
$\ip\bv{\be_{j}}=v_{i}\dd_{ij}=v_{j}$, and
$\ip\bu\bv=u_{i}v_{j}\dd_{ij}=u_{i}v_{i}$.
\end{proof}

\begin{remark}[the summation convention, stated once]\label{rem:sum}
In $v_{i}\be_{i}$ the index $i$ occurs twice and is summed; it is a dummy
and may be relabelled, $v_{i}\be_{i}=v_{j}\be_{j}$. An index occurring once
per term is free and must match across every term, as $j$ in
$\ip\bv{\be_{j}}=v_{j}$. A free index must never share a label with a dummy
in the same term, and a symbol appearing three or more times in a term is
undefined. Intuition: the dummy index is a hidden $\textstyle\sum$; the
free index is the ``address'' of the component the equation is really
about.
\end{remark}

\begin{remark}[projection, and why $\bze$ is orthogonal to all]\label{rem:proj}
$\bv=(\ip\bv{\be_{k}})\be_{k}$ exhibits $v_{i}=\ip\bv{\be_{i}}$ as the
orthogonal projection of $\bv$ on $\be_{i}$: geometrically, drop a
perpendicular from the tip of $\bv$ onto the $\be_i$ axis. Axiom (d) with
$\bu=\bze$ gives $\ip\bze\bv=0$ for all $\bv$: the zero vector, having no
direction, is perpendicular to everything.
\end{remark}

\begin{proposition}[Gram--Schmidt]\label{prop:gs}
Every basis $\{\bu_{1},\dots,\bu_{n}\}$ yields an orthonormal basis
$\{\be_{i}\}$ with
\[\operatorname{span}\{\be_{1},\dots,\be_{k}\}
=\operatorname{span}\{\bu_{1},\dots,\bu_{k}\}\]
for each $k$; hence there are
infinitely many orthonormal bases when $n\ge1$.
\end{proposition}

\begin{proof}
Recursively $\bw_{1}=\bu_{1}$ and
$\bw_{k}=\bu_{k}-\sum_{i<k}(\ip{\bu_{k}}{\be_{i}})\be_{i}$,
$\be_{k}=\bw_{k}/|\bw_{k}|$. Independence of the $\bu_i$ gives
$\bw_{k}\neq\bze$, and for $j<k$,
$\ip{\bw_{k}}{\be_{j}}=\ip{\bu_{k}}{\be_{j}}-\ip{\bu_{k}}{\be_{j}}=0$.
Infinitude: for any orthogonal $\bQ$ (Definition~\ref{def:orth}),
$\{\bQ\be_{i}\}$ is again orthonormal, and such $\bQ$ form a continuum.
\end{proof}

Many physical quantities are linear scalar functions of a vector: the work
of a fixed force along a variable displacement, a single component, a
directional derivative. Any such function is fixed by its values on the
basis, and reassembling those values yields one vector whose dot product
reproduces the function. Isolating this fact now pays off later, when the
increment of a field becomes its gradient.

\begin{theorem}[Riesz representation]\label{thm:riesz}
$f:\En\to\Rn$ is linear iff there is a unique $\ba\in\En$ with
$f(\bv)=\ip\ba\bv$ for all $\bv$; explicitly $\ba=f(\be_{i})\be_{i}$.
\end{theorem}

\begin{proof}
($\Leftarrow$) $\bv\mapsto\ip\ba\bv$ is linear by (b),(c).
($\Rightarrow$) By linearity and Proposition~\ref{prop:comp},
$f(\bv)=f(v_{i}\be_{i})=v_{i}f(\be_{i})=a_{i}v_{i}=\ip\ba\bv$ with
$a_{i}:=f(\be_{i})$. Uniqueness: $\ip{\ba}{\bv}=\ip{\ba'}{\bv}\,\forall\bv
\Rightarrow\ip{(\ba-\ba')}{\bv}=0$; take $\bv=\ba-\ba'$ and use (d).
\end{proof}

\begin{example}[the Riesz construction in action]\label{ex:riesz}
Suppose $f:\Et\to\Rn$ is linear and we know its values on the standard
basis: $f(\be_{1})=3$, $f(\be_{2})=-1$, $f(\be_{3})=2$. Set
$a_{i}:=f(\be_{i})$, so $\ba=3\be_{1}-\be_{2}+2\be_{3}$. Then for any
$\bv=v_{1}\be_{1}+v_{2}\be_{2}+v_{3}\be_{3}$,
\[
f(\bv)=f(v_{i}\be_{i})=v_{i}f(\be_{i})=3v_{1}-v_{2}+2v_{3}=\ip\ba\bv.
\]
The abstract function $f$ is now the concrete operation ``dot with
$(3,-1,2)$''. Every linear scalar function on $\Et$ hides a single vector;
finding that vector is the entire content of the Riesz theorem.
\end{example}

\begin{remark}[why Riesz is the engine of continuum calculus]\label{rem:rieszengine}
The pattern ``linear functional $\Rightarrow$ dot product with a fixed
vector'' will reappear every time we define a derivative. The gradient of
a scalar field (\S\ref{ss:grad}) is the Riesz representer of the linear
part of the increment. For a vector field the increment is linear and
vector valued, so its representer is a \emph{tensor} (\S\ref{ss:gradv}).
Recognising this template once makes every later gradient obvious.
\end{remark}

\subsection{Change-of-basis formulas}\label{ss:cob}

A vector is a geometric object, but its components $v_i=\ip\bv{\be_i}$ are
read off a chosen basis. Turn the basis and the same $\bv$ shows different
numbers. Mechanics must track that change exactly, so that physical
statements can be made basis independent. The bookkeeping is a single
matrix of dot products between the old and new axes, and it turns out to be
orthogonal.

\begin{definition}[Direction cosines]\label{def:dircos}
For two orthonormal bases $\{\be_{i}\}$, $\{\be'_{i}\}$, set
$a_{ij}=\ip{\be_{i}}{\be'_{j}}$.
\end{definition}

\begin{theorem}[Transformation law and orthogonality]\label{thm:cob}
With $v_{i}=\ip\bv{\be_{i}}$, $v'_{i}=\ip\bv{\be'_{i}}$,
\[
v'_{j}=a_{ij}v_{i},\qquad v_{i}=a_{ij}v'_{j},\qquad
a_{ik}a_{jk}=\dd_{ij}=a_{ki}a_{kj}.
\]
\end{theorem}

\begin{proof}
Expanding each old axis in the new basis,
$\be_{i}=\ip{\be_{i}}{\be'_{j}}\be'_{j}=a_{ij}\be'_{j}$, so
$v'_{j}\be'_{j}=\bv=v_{i}\be_{i}=v_{i}a_{ij}\be'_{j}$, whence
$v'_{j}=a_{ij}v_{i}$ by independence of $\{\be'_{j}\}$; swapping the two
bases gives the inverse. Orthogonality:
$a_{ik}a_{jk}=\ip{\be_{i}}{\be'_{k}}\ip{\be_{j}}{\be'_{k}}
=\ip{\be_{i}}{(\ip{\be_{j}}{\be'_{k}}\be'_{k})}=\ip{\be_{i}}{\be_{j}}
=\dd_{ij}$, and $a_{ki}a_{kj}=\dd_{ij}$ likewise.
\end{proof}

\begin{remark}[what ``orthogonal matrix'' means here]\label{rem:cobintuition}
$a_{ik}a_{jk}=\dd_{ij}$ says the rows of $(a_{ij})$ are orthonormal; it is
the algebraic echo of the geometric fact that both frames are orthonormal.
Nothing is rotated in space, only the labels; the vector itself is
untouched. This is the passive view of a change of basis, to be contrasted
later with the active view, where an orthogonal \emph{tensor} genuinely
moves vectors (\S\ref{ss:ops}).
\end{remark}

\subsection{Vector products in $\Et$}\label{ss:cross}

Two vectors span a parallelogram. Its \emph{area}, and the \emph{normal}
direction to its plane, are geometric data the inner product cannot supply.
In exactly three dimensions a plane has a unique normal line, so ``area
times normal'' can be packed into one vector, the cross product; feeding a
third vector then measures the signed \emph{volume} of the spanned
parallelepiped, the box product. These are the tools behind area, volume,
orientation, determinants, and the axial vector of a rotation.

\begin{definition}[Cross and box products]\label{def:cross}
$\times:\Et\times\Et\to\Et$ satisfies, for all $\bu,\bv,\bw$ and
$\alpha,\beta$,
\[
\text{(a)}\ \bu\times\bv=-\bv\times\bu,\quad
\text{(b)}\ |\bu\times\bv|=|\bu||\bv|\sin\theta,
\]
\[
\text{(c)}\ (\alpha\bu+\beta\bv)\times\bw
 =\alpha\,\bu\times\bw+\beta\,\bv\times\bw,
\]
\[
\text{(d)}\ \tbox\bu\bv\bw:=\ip\bu{(\bv\times\bw)}
=\ip\bv{(\bw\times\bu)}=\ip\bw{(\bu\times\bv)}.
\]
\end{definition}

\begin{proposition}\label{prop:boxprops}
$\bu\times\bu=\bze$; $\bu\times\bv\perp\bu,\bv$; and $\tbox\cdot\cdot\cdot$
is trilinear and alternating (it changes sign under any transposition of
its arguments).
\end{proposition}

\begin{proof}
(a) gives $\bu\times\bu=-\bu\times\bu=\bze$; then (d) gives
$\ip\bu{(\bu\times\bv)}=\tbox\bu\bu\bv=0$ and
$\ip\bv{(\bu\times\bv)}=\tbox\bv\bu\bv=0$. Trilinearity in slot one is (c)
via (d); the other slots follow from cyclicity in (d). Alternation: (a)
gives $\tbox\bu\bv\bw=-\ip\bu{\bw\times\bv}=-\tbox\bu\bw\bv$, and cyclicity
extends this to every transposition.
\end{proof}

\begin{definition}[Right-handed basis; permutation symbol]\label{def:perm}
An orthonormal basis of $\Et$ is right handed iff
$\be_{1}\times\be_{2}=\be_{3}$ and cyclically. Then
\[
\eps_{ijk}:=\tbox{\be_{i}}{\be_{j}}{\be_{k}}=
\begin{cases}+1,&(i,j,k)\ \text{even perm.\ of}\ (1,2,3),\\
-1,&(i,j,k)\ \text{odd perm.},\\ 0,&\text{repeated index.}\end{cases}
\]
\end{definition}

\begin{proposition}[Components of the products]\label{prop:crosscomp}
In a right-handed orthonormal basis,
$\be_{i}\times\be_{j}=\eps_{ijk}\be_{k}$,
$(\ba\times\bb)_{k}=\eps_{ijk}a_{i}b_{j}$, and
$\tbox\ba\bb\bc=\eps_{ijk}a_{i}b_{j}c_{k}$.
\end{proposition}

\begin{proof}
$\be_{i}\times\be_{j}=\ip{(\be_{i}\times\be_{j})}{\be_{k}}\be_{k}
=\eps_{ijk}\be_{k}$; then
$\ba\times\bb=a_{i}b_{j}\eps_{ijk}\be_{k}$ and
$\tbox\ba\bb\bc=\ip{(\ba\times\bb)}\bc=\eps_{ijk}a_{i}b_{j}c_{k}$.
\end{proof}

\begin{lemma}[Alternating trilinear forms are one dimensional]\label{lem:altform}
If $\omega:\Et\times\Et\times\Et\to\Rn$ is trilinear and alternating, then
$\omega(\ba,\bb,\bc)=\omega(\be_{1},\be_{2},\be_{3})\,\tbox\ba\bb\bc$.
\end{lemma}

\begin{proof}
\[
\begin{aligned}
\omega(\ba,\bb,\bc)&=a_{i}b_{j}c_{k}\,\omega(\be_{i},\be_{j},\be_{k})\\
&=a_{i}b_{j}c_{k}\,\eps_{ijk}\,\omega(\be_{1},\be_{2},\be_{3})\\
&=\omega(\be_{1},\be_{2},\be_{3})\,\tbox\ba\bb\bc,
\end{aligned}
\]
using
$\omega(\be_{i},\be_{j},\be_{k})=\eps_{ijk}\omega(\be_1,\be_2,\be_3)$ from
alternation.
\end{proof}

\begin{remark}[the workhorse]\label{rem:workhorse}
Lemma~\ref{lem:altform} is the one fact that makes determinants and
invariants painless: \emph{any} quantity linear in three vectors and sign
reversing under a swap is a fixed multiple of the box product. Each time we
invoke it, a page of index juggling collapses to a line. Intuitively, in
three dimensions there is essentially one way to measure oriented volume,
and everything alternating and trilinear must be proportional to it.
\end{remark}

\begin{definition}[Determinant]\label{def:det}
For $B=(b_{ij})\in\Rn^{3\times3}$, $\det B:=\eps_{ijk}b_{i1}b_{j2}b_{k3}$.
\end{definition}

\begin{proposition}\label{prop:detforms}
$\eps_{mnp}\det B=\eps_{ijk}b_{im}b_{jn}b_{kp}$, hence
$\eps_{ijk}\eps_{mnp}=\det(\dd_{\bullet\bullet})$, the $3\times3$
determinant of the Kronecker array with rows $i,j,k$ and
columns $m,n,p$.
\end{proposition}

\begin{proof}
$(\bx,\by,\bz)\mapsto\eps_{ijk}x_{i}y_{j}z_{k}$ is trilinear and
alternating in its three columns, so by Lemma~\ref{lem:altform} equals
$\tbox\bx\by\bz$ times its value $1$ at $(\be_1,\be_2,\be_3)$. Applied to
columns $(b_{\cdot m},b_{\cdot n},b_{\cdot p})$ this gives
$\eps_{ijk}b_{im}b_{jn}b_{kp}=\eps_{mnp}\det B$. Putting $b_{ij}=\dd_{ij}$
(so $\det B=1$) yields the second identity.
\end{proof}

\begin{theorem}[$\eps$--$\dd$ identity]\label{thm:epsdelta}
$\eps_{ijk}\eps_{mnk}=\dd_{im}\dd_{jn}-\dd_{in}\dd_{jm}$,
$\ \eps_{ijk}\eps_{mjk}=2\dd_{im}$, $\ \eps_{ijk}\eps_{ijk}=6$.
\end{theorem}

\begin{proof}
Expand the determinant of Proposition~\ref{prop:detforms}, set $p=k$, sum
($\dd_{kk}=3$):
$\eps_{ijk}\eps_{mnk}=\dd_{im}(3\dd_{jn}-\dd_{jn})-\dd_{in}(3\dd_{jm}
-\dd_{jm})+(\dd_{in}\dd_{jm}-\dd_{im}\dd_{jn})
=\dd_{im}\dd_{jn}-\dd_{in}\dd_{jm}$. Set $n=j$, sum:
$\eps_{ijk}\eps_{mjk}=3\dd_{im}-\dd_{im}=2\dd_{im}$. Set $m=i$, sum:
$\eps_{ijk}\eps_{ijk}=2\dd_{ii}=6$.
\end{proof}

\begin{corollary}[Determinant, contracted]\label{cor:det6}
\[\det B=\tfrac16\,\eps_{ijk}\eps_{mnp}\,b_{im}b_{jn}b_{kp}.\]
\end{corollary}

\begin{proof}
Contract $\eps_{ijk}b_{im}b_{jn}b_{kp}=\eps_{mnp}\det B$ with
$\eps_{mnp}$ and use $\eps_{mnp}\eps_{mnp}=6$.
\end{proof}

\subsection{Tensors}\label{ss:tensor}

Riesz handled linear rules returning a \emph{number}. Physics abounds in
linear rules returning a \emph{vector}: stress sends a surface normal to a
traction, the inertia tensor sends angular velocity to angular momentum,
the deformation gradient sends a material fibre to its deformed image. The
object encoding such a rule is a second order tensor. Its simplest instance,
$\ba\otimes\bb$, is the rule ``project onto $\bb$, then point along $\ba$''.

\begin{definition}[Tensor]\label{def:tensor}
A (second order) tensor is a linear map $\bA:\Et\to\Et$. Equality, scalar
multiple and sum are defined by $\bA=\bB\iff\bA\bv=\bB\bv\,\forall\bv$,
$(\alpha\bA)\bv=\alpha(\bA\bv)$, $(\bA+\bB)\bv=\bA\bv+\bB\bv$. The set $\Lin$
is a vector space; $\Ob\bv=\bze$ and $\I\bv=\bv$ for all $\bv$.
\end{definition}

\begin{definition}[Tensor product]\label{def:tp}
For $\ba,\bb\in\Et$, $(\ba\otimes\bb)\bv:=(\ip\bb\bv)\,\ba$.
\end{definition}

\begin{theorem}[Component representation]\label{thm:comprep}
$\{\be_{i}\otimes\be_{j}\}$ is a basis of $\Lin$, and
$\bA=A_{ij}\be_{i}\otimes\be_{j}$ with $A_{ij}=\ip{\be_{i}}{\bA\be_{j}}$;
hence $\dim\Lin=9$.
\end{theorem}

\begin{proof}
\emph{Spanning.} For $\bv=v_{i}\be_{i}$,
$\bA\bv=v_{i}\bA\be_{i}=v_{i}\ip{\be_{j}}{\bA\be_{i}}\be_{j}
=A_{ji}(\ip{\be_{i}}\bv)\be_{j}=(A_{ji}\be_{j}\otimes\be_{i})\bv$, so
$\bA=A_{ij}\be_{i}\otimes\be_{j}$ with $A_{ij}=\ip{\be_{i}}{\bA\be_{j}}$.
\emph{Independence.} $\alpha_{ij}\be_{i}\otimes\be_{j}=\Ob$ applied to
$\be_{k}$ gives $\alpha_{ik}\be_{i}=\bze$, so $\alpha_{ik}=0$.
\end{proof}

\begin{example}[reading off components of tensor products]\label{ex:tp}
\emph{(i) The tensor product $\ba\otimes\bb$.}
By Definition~\ref{def:tp},
$(\ba\otimes\bb)\be_{j}=(\ip\bb{\be_{j}})\ba=b_{j}\ba$. Dotting with
$\be_{i}$,
\[
(\ba\otimes\bb)_{ij}=\ip{\be_{i}}{(\ba\otimes\bb)\be_{j}}
=\ip{\be_{i}}{b_{j}\ba}=b_{j}\ip{\be_{i}}{\ba}=a_{i}b_{j}.
\]
Thus the component matrix is the outer product of the column $(a_i)$ with
the row $(b_j)$:
\[
\bigl((\ba\otimes\bb)_{ij}\bigr)=
\begin{pmatrix}a_{1}b_{1}&a_{1}b_{2}&a_{1}b_{3}\\
a_{2}b_{1}&a_{2}b_{2}&a_{2}b_{3}\\
a_{3}b_{1}&a_{3}b_{2}&a_{3}b_{3}\end{pmatrix}.
\]
Such a matrix has rank $\le1$: every row is a multiple of $(b_1,b_2,b_3)$.
Consequently $\det(\ba\otimes\bb)=0$ unless $\ba$ or $\bb$ is $\bze$.

\emph{(ii) The identity tensor.}
$\I\be_{j}=\be_{j}$, so $I_{ij}=\ip{\be_{i}}{\be_{j}}=\dd_{ij}$. In the
basis expansion,
$\I=\dd_{ij}\be_{i}\otimes\be_{j}=\be_{1}\otimes\be_{1}+\be_{2}\otimes\be_{2}
+\be_{3}\otimes\be_{3}=\be_{i}\otimes\be_{i}$,
with matrix $(I_{ij})=\operatorname{diag}(1,1,1)$.

\emph{(iii) The zero tensor.}
$\Ob\bv=\bze$ for all $\bv$, so $O_{ij}=\ip{\be_{i}}{\bze}=0$;
every entry of the matrix is zero.
\end{example}

\begin{remark}[outer versus inner product]\label{rem:outer}
Two operations eat a pair of vectors and return very different objects.
The \emph{inner product} $\ip\ba\bb=a_ib_i$ sums all matching pairs and
returns one number: it measures overlap and forgets direction. The
\emph{outer product} (tensor product) stores every pair, returning
nine numbers arranged in the matrix $(a_ib_j)$: it remembers both
directions. In $\Et$ one may think of $\ba\otimes\bb$ as the instruction
``measure how much $\bv$ aligns with $\bb$ (via $\ip\bb\bv$), then lay
that amount along $\ba$''. The trace connects the two:
$\tr(\ba\otimes\bb)=a_ib_j\dd_{ij}=a_ib_i=\ip\ba\bb$,
collapsing the nine entries back to one.
\end{remark}

\begin{example}[a diagonal tensor: the reflection]\label{ex:diag}
Suppose a linear map $\bA$ is specified by its action on the basis:
$\bA\be_{1}=-\be_{1}$, $\bA\be_{2}=\be_{2}$, $\bA\be_{3}=\be_{3}$.
We compute every component $A_{ij}=\ip{\be_{i}}{\bA\be_{j}}$ in turn.

\emph{Column $j=1$}: $\bA\be_{1}=-\be_{1}$, so
$A_{11}=\ip{\be_{1}}{-\be_{1}}=-1$,
$A_{21}=\ip{\be_{2}}{-\be_{1}}=0$,
$A_{31}=\ip{\be_{3}}{-\be_{1}}=0$.

\emph{Column $j=2$}: $\bA\be_{2}=\be_{2}$, so
$A_{12}=0$, $A_{22}=1$, $A_{32}=0$.

\emph{Column $j=3$}: $\bA\be_{3}=\be_{3}$, so
$A_{13}=0$, $A_{23}=0$, $A_{33}=1$.

Hence
\[
(A_{ij})=\begin{pmatrix}-1&0&0\\0&1&0\\0&0&1\end{pmatrix}
=\operatorname{diag}(-1,1,1),
\]
and the tensor product form is
$\bA=-\be_{1}\otimes\be_{1}+\be_{2}\otimes\be_{2}+\be_{3}\otimes\be_{3}$.
For any $\bv=v_i\be_i$,
$\bA\bv=-v_1\be_1+v_2\be_2+v_3\be_3$: the component along $\be_1$ is
negated while those along $\be_2,\be_3$ are unchanged. This is precisely
the \emph{mirror reflection} in the plane perpendicular to $\be_1$.
\end{example}

\begin{remark}[why the diagonal form matters]\label{rem:diag}
A diagonal tensor $\bA=\lambda_i\be_i\otimes\be_i$ (no sum) acts by
scaling each basis direction independently:
$\bA\bv=\lambda_1v_1\be_1+\lambda_2v_2\be_2+\lambda_3v_3\be_3$. The
$\lambda_i$ are the eigenvalues and the $\be_i$ the eigenvectors. In
continuum mechanics, when $\bA$ is the (symmetric) stretch tensor, the
eigenvectors are the principal directions and $|\lambda_i|$ the principal
stretches. The spectral theorem (Chapter~3) guarantees that every
symmetric tensor can be written in this form for a suitable orthonormal
basis.
\end{remark}

\begin{lemma}\label{lem:Aab}
$\bA(\ba\otimes\bb)=(\bA\ba)\otimes\bb$.
\end{lemma}

\begin{proof}
$[\bA(\ba\otimes\bb)]\bv=\bA[(\ip\bb\bv)\ba]=(\ip\bb\bv)\bA\ba
=[(\bA\ba)\otimes\bb]\bv$.
\end{proof}

\begin{theorem}[Transformation law for tensor components]\label{thm:tenstrans}
With $a_{ij}=\ip{\be_{i}}{\be'_{j}}$, $A'_{ij}=a_{ki}a_{lj}A_{kl}$.
\end{theorem}

\begin{proof}
$\bA\be'_{n}=(A_{kl}\be_{k}\otimes\be_{l})\be'_{n}=A_{kl}a_{ln}\be_{k}$,
so $A'_{jn}=\ip{\be'_{j}}{\bA\be'_{n}}=A_{kl}a_{ln}\ip{\be'_{j}}{\be_{k}}
=A_{kl}a_{kj}a_{ln}$.
\end{proof}

\begin{remark}[order $=$ number of cosine factors]\label{rem:order}
Scalars, vectors and tensors transform with the direction cosine matrix
zero, one and two times: $\varphi'=\varphi$, $v'_{i}=a_{ki}v_{k}$,
$A'_{ij}=a_{ki}a_{lj}A_{kl}$. The count of $a$'s is the tensor's
\emph{order}, and is often taken as the definition of a Cartesian tensor.
To illustrate: the trace $\tr\bA=A_{ii}$ carries zero free indices, hence
zero cosines and transforms as a scalar ($\tr\bA'=A'_{ii}=a_{ki}a_{li}A_{kl}
=\dd_{kl}A_{kl}=A_{kk}=\tr\bA$). The traction vector
$t_i=A_{ij}n_j$ carries one free index, hence one cosine factor ($t'_i
=a_{ki}t_k$), and transforms as a vector. The tensor itself carries two.
Invariants (\S\ref{ss:ops}) exist precisely because even combinations of
cosines can pair up into Kronecker deltas, cancelling all basis dependence.
\end{remark}

\subsection{Algebraic operations with second-order tensors}\label{ss:ops}

\paragraph*{Transpose and symmetry.}
In $\ip\bu{\bA\bv}$ one often needs to let $\bA$ act on the other factor.
The tensor legalising this, $\ip\bu{\bA\bv}=\ip{\bA^{t}\bu}\bv$, is the
transpose. Whether a tensor equals or negates its transpose splits $\Lin$
into the two parts that rule mechanics: strain is symmetric, spin is skew.

\begin{definition}[Transpose]\label{def:transpose}
$\bA^{t}$ is defined by $\ip\bu{\bA\bv}=\ip\bv{\bA^{t}\bu}$ for all
$\bu,\bv$.
\end{definition}

\begin{proposition}\label{prop:transpose}
$(\bA^{t})_{ij}=A_{ji}$, $(\bA^{t})^{t}=\bA$, $(\bA\bB)^{t}=\bB^{t}\bA^{t}$.
\end{proposition}

\begin{proof}
$(\bA^{t})_{ij}=\ip{\be_{i}}{\bA^{t}\be_{j}}=\ip{\be_{j}}{\bA\be_{i}}
=A_{ji}$. The rest follow from Definition~\ref{def:transpose} and
uniqueness: $\ip\bu{(\bA\bB)\bv}=\ip{\bA^{t}\bu}{\bB\bv}
=\ip{\bB^{t}\bA^{t}\bu}\bv$.
\end{proof}

\begin{definition}[Symmetric, skew]\label{def:symskew}
$\bA$ is symmetric iff $\bA^{t}=\bA$, skew iff $\bA^{t}=-\bA$.
\end{definition}

\begin{proposition}[Axial vector of a skew tensor]\label{prop:axial}
For skew $\bA$ there is a unique $\bw$ with $\bA\bv=\bw\times\bv$ for all
$\bv$; $A_{ij}=-\eps_{ijk}w_{k}$, $w_{k}=-\tfrac12\eps_{kij}A_{ij}$.
\end{proposition}

\begin{proof}
Skewness gives $A_{ii}=0$ and three free entries; the general skew array is
$A_{ij}=-\eps_{ijk}w_{k}$. Then
$(\bA\bv)_{i}=-\eps_{ijk}w_{k}v_{j}=\eps_{ikj}w_{k}v_{j}
=(\bw\times\bv)_{i}$. Inversion by Theorem~\ref{thm:epsdelta}:
$-\eps_{ijm}A_{ij}=\eps_{ijm}\eps_{ijk}w_{k}=2w_{m}$.
\end{proof}

\paragraph*{Product and inverse.}

\begin{definition}[Product, inverse]\label{def:prodinv}
$(\bA\bB)\bv:=\bA(\bB\bv)$. $\bA$ is invertible iff $\bA\bu=\bv$ has a
unique solution $\bu=\bA^{-1}\bv$ for every $\bv$.
\end{definition}

\begin{proposition}\label{prop:prodinv}
$(\bA\bB)_{ij}=A_{ik}B_{kj}$; in general $\bA\bB\neq\bB\bA$; and if $\bA$ is
invertible, $\bA^{-1}\bA=\bA\bA^{-1}=\I$.
\end{proposition}

\begin{proof}
$(\bA\bB)\be_{j}=\bA(B_{kj}\be_{k})=B_{kj}A_{ik}\be_{i}$, so
$(\bA\bB)_{ij}=A_{ik}B_{kj}$; and
$(\bA^{-1}\bA)\bu=\bA^{-1}\bv=\bu$, $(\bA\bA^{-1})\bv=\bA\bu=\bv$.
\end{proof}

\paragraph*{Orthogonal tensors.}
Rigid motions neither stretch nor shear a body: they preserve every length
and angle, hence every inner product. The tensors doing this are the
orthogonal ones, and they are exactly those whose transpose is their
inverse. Those with $\det=+1$ are rotations; $\det=-1$ adds a reflection.

\begin{definition}[Orthogonal]\label{def:orth}
$\bA$ is orthogonal iff $\ip{\bA\bu}{\bA\bv}=\ip\bu\bv$ for all $\bu,\bv$.
\end{definition}

\begin{proposition}\label{prop:orth}
$\bA$ orthogonal $\iff\bA^{t}\bA=\I\iff\bA\bA^{t}=\I$; then
$\bA^{-1}=\bA^{t}$ and $\det\bA=\pm1$.
\end{proposition}

\begin{proof}
$\ip{\bA\bu}{\bA\bv}=\ip\bu{\bA^{t}\bA\bv}$, so orthogonality is
$\ip\bu{(\bA^{t}\bA-\I)\bv}=0\,\forall\bu,\bv$; take
$\bu=(\bA^{t}\bA-\I)\bv$ to get $\bA^{t}\bA=\I$, then
$\bA^{-1}=\bA^{t}$ and $\bA\bA^{t}=\I$. Finally
$(\det\bA)^{2}=\det\bA^{t}\det\bA=\det\I=1$ (Theorem~\ref{thm:det}).
\end{proof}

\paragraph*{Principal invariants.}
The nine components of a tensor change with the basis; its physics does
not. So we seek scalar combinations surviving every change of basis.
Building them from the box product, itself geometric, makes invariance
automatic. Three exist, and they are the coefficients of the characteristic
polynomial (Problem~23), hence carry the eigenvalue data used throughout
deformation analysis.

\begin{definition}[Invariants]\label{def:inv3}
For all $\ba,\bb,\bc$,
\[
\begin{aligned}
I_{1}(\bA)\tbox\ba\bb\bc&=\tbox{\bA\ba}\bb\bc+\tbox\ba{\bA\bb}\bc+\tbox\ba\bb{\bA\bc},\\
I_{2}(\bA)\tbox\ba\bb\bc&=\tbox{\bA\ba}{\bA\bb}\bc+\tbox{\bA\ba}\bb{\bA\bc}+\tbox\ba{\bA\bb}{\bA\bc},\\
I_{3}(\bA)\tbox\ba\bb\bc&=\tbox{\bA\ba}{\bA\bb}{\bA\bc}.
\end{aligned}
\]
\end{definition}

\begin{proposition}[Well-posedness]\label{prop:invwd}
$I_{1},I_{2},I_{3}$ do not depend on $\ba,\bb,\bc$.
\end{proposition}

\begin{proof}
Each right side is trilinear and alternating in $(\ba,\bb,\bc)$; by
Lemma~\ref{lem:altform} it equals a fixed scalar times $\tbox\ba\bb\bc$.
\end{proof}

\begin{theorem}[Trace]\label{thm:trace}
$I_{1}$ is linear, $I_{1}(\bA)=\tr\bA:=A_{ii}$,
$\tr(\be_{i}\otimes\be_{j})=\dd_{ij}$, $\tr(\bu\otimes\bv)=\ip\bu\bv$.
\end{theorem}

\begin{proof}
Definition~\ref{def:inv3}$_1$ is linear in $\bA$. With $\tbox{\be_1}{\be_2}{\be_3}=1$,
\[
\tr(\be_{i}\otimes\be_{j})
=\dd_{j1}\eps_{i23}+\dd_{j2}\eps_{1i3}+\dd_{j3}\eps_{12i}=\dd_{ij},
\]
so $\tr\bA=A_{ij}\dd_{ij}=A_{ii}$ and $\tr(\bu\otimes\bv)=u_iv_j\dd_{ij}
=\ip\bu\bv$.
\end{proof}

\begin{theorem}[Determinant]\label{thm:det}
$I_{3}(\bA)=\det(A_{ij})$; $\bA$ invertible $\iff\det\bA\neq0$;
$\det(\bA\bB)=\det\bA\det\bB$; $\det\bA^{t}=\det\bA$.
\end{theorem}

\begin{proof}
$I_{3}(\bA)=\tbox{\bA\be_1}{\bA\be_2}{\bA\be_3}=\eps_{ijk}A_{i1}A_{j2}A_{k3}
=\det(A_{ij})$ by Proposition~\ref{prop:crosscomp} and
Definition~\ref{def:det}. Multiplicativity and transpose invariance follow
from Proposition~\ref{prop:detforms} applied to $(A_{ik}B_{kj})$ and to
$(A_{ji})$. Invertibility: $\det\bA\neq0\iff$ the columns $\bA\be_i$ are
independent $\iff\bA$ is a bijection.
\end{proof}

\begin{remark}[determinant $=$ oriented volume ratio]\label{rem:vol}
When $\tbox\ba\bb\bc>0$ it is the volume of the parallelepiped on
$\ba,\bb,\bc$, so Definition~\ref{def:inv3}$_3$ reads
$\det\bA=\tbox{\bA\ba}{\bA\bb}{\bA\bc}/\tbox\ba\bb\bc$: $\det\bA$ is the
factor by which $\bA$ scales oriented volume. Hence $\det>0$ means
orientation preserving, and later $\det$ of the deformation gradient is the
local volume ratio, with $\det=1$ signalling incompressibility.
\end{remark}

\begin{definition}[Cofactor]\label{def:cof}
$\bA^{*}$ is defined by $\bA^{*}(\ba\times\bb)=\bA\ba\times\bA\bb$.
\end{definition}

\begin{proposition}\label{prop:I2cof}
$I_{2}(\bA)=\tr\bA^{*}$; $I_{2}$ is quadratic and $I_{3}$ cubic in $\bA$.
\end{proposition}

\begin{proof}
Applying the trace formula to $\bA^{*}$ and
$\bA^{*}(\be_2\times\be_3)=\bA\be_2\times\bA\be_3$ etc.\ (see Problem~22a
for the full computation) gives $\tr\bA^{*}=I_{2}(\bA)$. Degrees are read
off Definition~\ref{def:inv3}$_{2,3}$.
\end{proof}

\begin{remark}[the cofactor maps area elements]%
\label{rem:cofgeom}
If $\ba\times\bb$ is the area vector of a parallelogram, then
$\bA^{*}(\ba\times\bb)=\bA\ba\times\bA\bb$ is the area vector of the
deformed parallelogram. Thus the determinant tells how $\bA$ scales
\emph{volumes} (Remark~\ref{rem:vol}), while the cofactor tells how
$\bA$ scales \emph{oriented areas}. In the notation of continuum
mechanics, the deformation gradient $\vek{F}$ satisfies
$n\,da=\vek{F}^{*}(N\,dA)$ (Nanson's formula): the cofactor converts
a reference area element into its current image. Because area is one
dimension lower than volume, the cofactor is quadratic in $\bA$ while
the determinant is cubic; this is exactly what the degrees in
Proposition~\ref{prop:I2cof} express. The identity
$\bA^{t}\bA^{*}=(\det\bA)\I$ (Problem~22c) ties the two together.
\end{remark}

\paragraph*{Inner product of tensors.}
$\Lin$ is nine dimensional; to do analysis on it, e.g.\ to call a stress
increment small, it needs a length. The natural one pairs components
entrywise and is generated intrinsically by trace and transpose.

\begin{definition}[Tensor inner product]\label{def:tip}
$\bA\cdot\bB:=\tr(\bA\bB^{t})$.
\end{definition}

\begin{theorem}\label{thm:tip}
$\bA\cdot\bB=A_{ij}B_{ij}=\bB\cdot\bA$; $|\bA|=(A_{ij}A_{ij})^{1/2}$ makes
$\Lin$ isometric to $\mathbb{E}^{9}$, so
$|\bA\cdot\bB|\le|\bA||\bB|$ and $|\bA+\bB|\le|\bA|+|\bB|$; and
$\bA\cdot(\bu\otimes\bv)=\ip\bu{\bA\bv}$.
\end{theorem}

\begin{proof}
$(\bA\bB^{t})_{ik}=A_{ij}B_{kj}$, so $\tr(\bA\bB^{t})=A_{ij}B_{ij}$,
symmetric. This is the Euclidean inner product of nine components, so the
inequalities transfer from Theorem~\ref{thm:cs} and
Corollary~\ref{cor:tri}. Finally, by Lemma~\ref{lem:Aab},
$\bA\cdot(\bu\otimes\bv)=\tr((\bA\bv)\otimes\bu)=\ip{\bA\bv}\bu
=\ip\bu{\bA\bv}$.
\end{proof}

\begin{definition}[Two orthogonal splittings]\label{def:splits}
$\bA=\Sym\bA+\Skw\bA$ with $\Sym\bA=\tfrac12(\bA+\bA^{t})$,
$\Skw\bA=\tfrac12(\bA-\bA^{t})$; and
$\bA=\Sph\bA+\Dev\bA$ with $\Sph\bA=\tfrac13(\tr\bA)\I$,
$\Dev\bA=\bA-\tfrac13(\tr\bA)\I$.
\end{definition}

\begin{proposition}\label{prop:splits}
Each summand is of the named type, the splittings are unique,
$\tr(\Dev\bA)=0$, and $\Sym\bA\cdot\Skw\bB=0$,
$\Sph\bA\cdot\Dev\bB=0$ for all $\bA,\bB$. (Details in Problems~14--15.)
\end{proposition}

\begin{proof}
$\tfrac12(\bA+\bA^t)$ is symmetric, $\tfrac12(\bA-\bA^t)$ skew, and their
sum is $\bA$; $\tr(\bA-\tfrac13(\tr\bA)\I)=\tr\bA-\tr\bA=0$. For symmetric
$\bS$ and skew $\bW$, $\bS\cdot\bW=S_{ij}W_{ij}=S_{ji}(-W_{ji})
=-\bS\cdot\bW=0$; and
$\Sph\bA\cdot\Dev\bB=\tfrac13(\tr\bA)\tr(\Dev\bB)=0$. Uniqueness is proved
in Problems~14a,~15a.
\end{proof}

%=======================================================================
\section{Scalar, vector and tensor fields; gradient, divergence and curl;
integral formulae}\label{sec:fields}
%=======================================================================

\subsection{Euclidean point spaces}\label{ss:pts}

The algebra above treats \emph{arrows}: they add, scale, and share a zero.
Mechanics needs instead the \emph{places} occupied by particles, and places
behave differently. Two places cannot be added: ``the doorway plus the
window'' names nothing. There is no zero place until someone plants a flag
and calls it the origin. What \emph{is} intrinsic is the displacement from
one place to another, and a displacement does not move when the flag moves.
So positions are not vectors, but their differences are. We therefore work
with a set of points whose subtraction lands in a vector space.

\begin{definition}[Euclidean point space]\label{def:eps}
A Euclidean point space is a set $\Et$ of points
$\mathsf{x},\mathsf{y},\dots$ with a subtraction
$(\mathsf{x},\mathsf{y})\mapsto\mathsf{y}-\mathsf{x}$ into a three
dimensional inner product space $\Et$, the translation space, such that
\begin{enumerate}[label=\textup{(P\arabic*)},leftmargin=2.6em]
\item \textup{(transport)} for every $\mathsf{x}$ and $\bv\in\Et$ there is
      a unique $\mathsf{y}$ with $\mathsf{y}-\mathsf{x}=\bv$, written
      $\mathsf{y}=\mathsf{x}+\bv$;
\item \textup{(Chasles)} $(\mathsf{y}-\mathsf{x})+(\mathsf{z}-\mathsf{y})
      =\mathsf{z}-\mathsf{x}$.
\end{enumerate}
\end{definition}

\begin{remark}[what each axiom buys]\label{rem:axioms}
(P1) says the translation space acts freely and transitively on points:
every displacement carries each point to a unique new one, and every pair
of points is joined by a unique displacement. This is precisely what lets
us later compare a field's value at $\mathsf{x}$ with its values at
neighbours $\mathsf{x}+s\bu$, the comparison a derivative is built from.
Setting $\mathsf{z}=\mathsf{x}$ in (P2), with $\mathsf{x}-\mathsf{x}=\bze$
(the unique solution of $\mathsf{y}-\mathsf{x}=\bze$ being
$\mathsf{y}=\mathsf{x}$), gives $\mathsf{x}-\mathsf{y}
=-(\mathsf{y}-\mathsf{x})$: reversing a displacement negates its arrow.
\end{remark}

The definition is origin free, which is clean but useless for computation.
Naming a reference place, the origin, is planting the flag; every point
then acquires one arrow with three components. The cost is dependence on
the flag, measured precisely below.

\begin{proposition}[Position vector, Cartesian coordinates]\label{prop:posvec}
Fix an origin $\mathsf{O}$. Then (a) every point $\mathsf{x}$ determines a
unique position vector $\bx:=\mathsf{x}-\mathsf{O}$, and
$\mathsf{x}\mapsto\bx$ is a bijection; (b) $\bx=x_i\be_i$ with
$x_i=\ip\bx{\be_i}$; (c) under $\mathsf{O}\to\mathsf{O}'$,
$\bx'=\bx-\bc$ with $\bc=\mathsf{O}'-\mathsf{O}$, and $x_i'=x_i-c_i$.
\end{proposition}

\begin{proof}
(a) $\bx\in\Et$ is well defined. Injective: $\mathsf{x}_1-\mathsf{O}
=\mathsf{x}_2-\mathsf{O}$ gives, by (P2),
$\bze=(\mathsf{x}_1-\mathsf{O})+(\mathsf{O}-\mathsf{x}_2)
=\mathsf{x}_1-\mathsf{x}_2$. Surjective: $\mathsf{x}=\mathsf{O}+\bv$ has
$\bx=\bv$. (b) $\ip\bx{\be_i}=x_j\dd_{ji}=x_i$. (c) By (P2),
$\mathsf{x}-\mathsf{O}'=(\mathsf{x}-\mathsf{O})-\bc$; dot with $\be_i$.
\end{proof}

\begin{remark}[the deliberate abuse $\bx(=\bx-\bo)$]\label{rem:abuse}
The book writes the position vector as a bare $\bx$, so the same letter
$\bx$ names both the point and its position vector. Part (a) makes this
legitimate \emph{after} an origin is fixed, since the correspondence is
then a bijection; part (c) is the warning label, that any formula with a
bare $\bx$ silently carries the chosen origin. Good continuum theory keeps
its final physical statements independent of that choice.
\end{remark}

\begin{definition}[Fields]\label{def:field}
A scalar field is $\varphi:\Et\to\Rn$, a vector field $\bv:\Et\to\Et$, a
tensor field $\bA:\Et\to\Lin$; their argument is the position $\bx$.
\end{definition}

\subsection{Gradients}\label{ss:grad}

In one variable a differentiable $f$ satisfies
$f(x_0+h)=f(x_0)+f'(x_0)h+(\text{small})$: near $x_0$ the graph is a
straight line, and $f'(x_0)$ its slope. The single idea of differential
calculus is that a smooth function is, up close, linear. We now transplant
it to a field, where the increment depends on the \emph{vector}
$\bh=\bx-\bx_0$, so the linear part is a linear scalar function of $\bh$.
Two things must be made precise: what ``small'' means for the remainder,
and how to replace the abstract linear part by one computable vector.

\begin{definition}[Landau symbol]\label{def:landau}
For $r$ defined near $0$, $r(\eps)=o(\eps)$ iff $r(\eps)/\eps\to0$ as
$\eps\to0$; equivalently, for every $\eta>0$ there is $\dd>0$ with
$|r(\eps)|\le\eta\eps$ for $0<\eps<\dd$.
\end{definition}

\begin{remark}[why ``faster than linear'' removes ambiguity]\label{rem:kills}
If a field had two linear parts $f_0,g_0$ in \eqr{1.96}, subtracting gives
$f_0(\bh)-g_0(\bh)=-o(|\bh|)$; putting $\bh=s\bu$ ($|\bu|=1$) and using
linearity, $s[f_0(\bu)-g_0(\bu)]=-o(s)$, so
$f_0(\bu)-g_0(\bu)=-o(s)/s\to0$. The left side is independent of $s$, hence
$0$ for every unit $\bu$, so $f_0=g_0$. Demanding the remainder die
\emph{faster} than $\bh$ is exactly what pins the linear part down.
\end{remark}

\begin{definition}[Differentiable scalar field]\label{def:diff}
$\varphi$ is differentiable at $\bx_0$ iff there is a linear
$f_0:\Et\to\Rn$ with
\begin{equation}
\varphi(\bx)=\varphi(\bx_0)+f_0(\bx-\bx_0)+o(|\bx-\bx_0|)\qquad(\bx\to\bx_0).
\tag{1.96}\label{eq:1.96}
\end{equation}
\end{definition}

\begin{theorem}[Existence, uniqueness, naming of the gradient]\label{thm:grad}
If $\varphi$ is differentiable at $\bx_0$, there is a unique
$\grad\varphi\rest{\bx_0}\in\Et$ with
$f_0(\bh)=\grad\varphi\rest{\bx_0}\cdot\bh$, hence
\begin{equation}
\varphi(\bx)=\varphi(\bx_0)+\grad\varphi\rest{\bx_0}\cdot(\bx-\bx_0)
+o(|\bx-\bx_0|). \tag{1.99}\label{eq:1.99}
\end{equation}
\end{theorem}

\begin{proof}
$f_0$ is a linear functional on $\Et$; by Riesz (Theorem~\ref{thm:riesz})
there is a unique $\bc$ with $f_0(\bh)=\ip\bc\bh$. Explicitly, with
$c_i:=f_0(\be_i)$ and $\bh=h_i\be_i$,
$f_0(\bh)=h_if_0(\be_i)=h_ic_i=\ip\bc\bh$. Name
$\bc=:\grad\varphi\rest{\bx_0}$ and substitute into \eqr{1.96}.
\end{proof}

\begin{remark}[the gradient is a slope, and it points uphill]\label{rem:slope}
$\grad\varphi\rest{\bx_0}$ is one vector attached to the place $\bx_0$. By
\eqr{1.99} the first order change under a step $\bh$ is
$\grad\varphi\cdot\bh$, largest when $\bh\parallel\grad\varphi$
(Cauchy--Schwarz). So the gradient points in the direction of steepest
increase and its length is that steepest rate; its level sets are the
directions of no first order change, orthogonal to $\grad\varphi$.
\end{remark}

\begin{remark}[the double construction reused throughout]\label{rem:double}
The argument has a fixed three step shape:
$\text{increment}\to\text{linear part}\to\text{Riesz representer}$. For a
scalar field the representer is a vector, $\grad\varphi$. Applied to a
\emph{vector} field the linear part is a linear map and Riesz upgrades to
``a linear map is a tensor'', producing $\grad\bv\in\Lin$
(\S\ref{ss:gradv}); applied to a flux it produces the divergence. Spotting
this template once explains every later derivative at a glance.
\end{remark}

\subsection{Directional derivatives}\label{ss:dir}

The differentiability limit \eqr{1.96} lets $\bx\to\bx_0$ from all directions
at once, a limit in three variables. The cheapest way to probe it is to fix
a direction and move only along it: set $\bv=\bx-\bx_0$, let $s=|\bv|$ so
that $\bv=s\bu$ with $|\bu|=1$, and watch
$f(s)=\varphi(\bx_0+s\bu)$ as an ordinary function of one real variable.
The coefficient of $s$ in its expansion is the rate of change of $\varphi$
along $\bu$. Let us compute it and see that the gradient already contains
the answer for every direction.

\begin{definition}[Directional derivative]\label{def:dir}
For a unit vector $\bu$ and fixed $\bx_0$, with $f(s)=\varphi(\bx_0+s\bu)$,
the derivative of $\varphi$ at $\bx_0$ along $\bu$ is $f'(0)$.
\end{definition}

\begin{theorem}[The gradient gives every directional derivative]\label{thm:dir}
If $\varphi$ is differentiable at $\bx_0$, then for every unit $\bu$
\begin{equation}
\Big(\tfrac{\dl f}{\dl s}\Big)\Big|_{s=0}=\bu\cdot\grad\varphi\rest{\bx_0}.
\tag{1.102}\label{eq:1.102}
\end{equation}
\end{theorem}

\begin{proof}
With $\bh=s\bu$, $|\bh|=s$ (as $|\bu|=1$, $s\ge0$), \eqr{1.99} gives
\begin{equation}
f(s)=\varphi(\bx_0)+\grad\varphi\rest{\bx_0}\cdot(s\bu)+o(s)
=\varphi(\bx_0)+s\,(\bu\cdot\grad\varphi\rest{\bx_0})+o(s). \tag{1.100}\label{eq:1.100}
\end{equation}
Subtract $f(0)$, divide by $s\neq0$:
$[f(s)-f(0)]/s=\bu\cdot\grad\varphi\rest{\bx_0}+o(s)/s$ (1.101); let
$s\to0$.
\end{proof}

Choosing the special directions $\be_i$, which change only one coordinate,
turns directional derivatives into ordinary partial derivatives and hands
us the components of the gradient.

\begin{corollary}[Partials as gradient components]\label{cor:partial}
With $\hat\varphi(x_1,x_2,x_3):=\varphi(x_i\be_i)$,
\begin{equation}
\be_i\cdot\grad\varphi\rest{\bx_0}
=\Big(\tfrac{\partial\hat\varphi}{\partial x_i}\Big)
\Big|_{(x_1^0,x_2^0,x_3^0)}=:\hat\varphi_{,i}\rest{\bx_0}. \tag{1.103--1.106}\label{eq:1.103}
\end{equation}
\end{corollary}

\begin{proof}
Take $\bu=\be_i$ in Theorem~\ref{thm:dir}; $\bx_0+s\be_i$ increments only
the $i$th coordinate, so
$\be_i\cdot\grad\varphi\rest{\bx_0}=\lim_{s\to0}\tfrac1s[
\hat\varphi(\dots,x_i^0+s,\dots)-\hat\varphi(\dots,x_i^0,\dots)]
=\partial\hat\varphi/\partial x_i$. As $\bx_0$ is arbitrary,
$\be_i\cdot\grad\varphi=\hat\varphi_{,i}$.
\end{proof}

\begin{theorem}[Cartesian gradient]\label{thm:cart}
$\grad\varphi=\hat\varphi_{,i}\be_i=\varphi_{,i}\be_i$.
\hfill\textup{(1.107--1.108)}
\end{theorem}

\begin{proof}
By orthonormal expansion (Proposition~\ref{prop:comp}),
$\bw=(\ip\bw{\be_i})\be_i$; take $\bw=\grad\varphi$ and use
$\ip{\grad\varphi}{\be_i}=\hat\varphi_{,i}$. Dropping the caret gives
(1.108).
\end{proof}

\begin{remark}[caret and comma]\label{rem:caret}
$\varphi(\bx)$ eats a point; $\hat\varphi(x_1,x_2,x_3)$ eats three real
numbers. Different domains make them different functions, yet they take
equal values when $\bx=x_i\be_i$. The comma abbreviates a Cartesian
partial, $\varphi_{,i}=\partial\varphi/\partial x_i$, from here on;
dropping the caret is the same controlled abuse as writing a bare $\bx$ for
the position (Remark~\ref{rem:abuse}).
\end{remark}

\subsection{Parameterized paths and the chain rule}\label{ss:chain}

Directional derivatives run along straight lines of fixed direction. A
particle's trajectory, or a coordinate line in curvilinear coordinates,
bends: its direction $\bx'(u)$ varies along the way. To differentiate a
field along such a curve we compose $\varphi$ with the path and use the one
variable derivative. The resulting rule, in its coordinate free form
$\dl\varphi=\grad\varphi\cdot\dl\bx$, is powerful enough to generate
$\grad\varphi$ in any coordinate system.

\begin{theorem}[Chain rule]\label{thm:chain}
For a smooth simple curve $\bx(u)$ and $\psi(u)=\varphi(\bx(u))$,
\begin{equation}
\psi'(u_0)=\grad\varphi\rest{\bx(u_0)}\cdot\bx'(u_0). \tag{1.112}\label{eq:1.112}
\end{equation}
\end{theorem}

\begin{proof}
With $\bx_0=\bx(u_0)$, $\bx_1=\bx(u_1)$, \eqr{1.99} gives
\begin{equation}
\psi(u_1)=\psi(u_0)+\grad\varphi\rest{\bx_0}\cdot(\bx_1-\bx_0)
+o(|\bx_1-\bx_0|). \tag{1.109}\label{eq:1.109}
\end{equation}
Smoothness of the path gives
\begin{equation}
\bx_1-\bx_0=\bx'(u_0)(u_1-u_0)+o(u_1-u_0), \tag{1.110}\label{eq:1.110}
\end{equation}
so $|\bx_1-\bx_0|=O(u_1-u_0)$ and $o(|\bx_1-\bx_0|)=o(u_1-u_0)$. Insert
\eqr{1.110} into \eqr{1.109}, subtract $\psi(u_0)$, divide by $u_1-u_0$:
\begin{equation}
\frac{\psi(u_1)-\psi(u_0)}{u_1-u_0}
=\grad\varphi\rest{\bx_0}\cdot\bx'(u_0)+\frac{o(u_1-u_0)}{u_1-u_0}, \tag{1.111}\label{eq:1.111}
\end{equation}
and let $u_1\to u_0$.
\end{proof}

\begin{corollary}[Coordinate free differential]\label{cor:dphi}
With $\dl\bx=\bx'(u)\,\dl u$,
\begin{equation}
\frac{\dl\varphi}{\dl u}=\grad\varphi\cdot\frac{\dl\bx}{\dl u},
\qquad \dl\varphi=\grad\varphi\cdot\dl\bx. \tag{1.113--1.114}\label{eq:1.113}
\end{equation}
This holds for every $\dl\bx$ and uses no coordinates.
\end{corollary}

\begin{remark}[the power of \eqr{1.114}]\label{rem:dphi}
Equation \eqr{1.114}, $\dl\varphi=\grad\varphi\cdot\dl\bx$, is coordinate free:
it never mentions a basis. Yet it can \emph{generate} the gradient formula
in \emph{any} coordinate system. The recipe is always the same: write
$\dl\varphi$ from the conventional chain rule, express the coordinate
increments as dot products with $\dl\bx$, read off $\grad\varphi$ by
matching with \eqr{1.114}. The next two examples execute this recipe in
Cartesian and cylindrical coordinates.
\end{remark}

\begin{example}[recovering the Cartesian formula]\label{ex:cartrecover}
\emph{Step 1 (conventional chain rule).} With
$\hat\varphi(x_1,x_2,x_3)=\varphi(x_i\be_i)$ and $\{\be_i\}$ fixed,
\[
\dl\hat\varphi=\frac{\partial\hat\varphi}{\partial x_1}\dl x_1
+\frac{\partial\hat\varphi}{\partial x_2}\dl x_2
+\frac{\partial\hat\varphi}{\partial x_3}\dl x_3
=\hat\varphi_{,i}\,\dl x_i.
\]
\emph{Step 2 (express $\dl x_i$ via $\dl\bx$).} Since the basis is fixed,
$\bx=x_i\be_i$ gives $\dl\bx=\dl x_i\,\be_i$. Dotting with $\be_i$
and using orthonormality, $\dl x_i=\be_i\cdot\dl\bx$. Insert:
\[
\dl\hat\varphi=\hat\varphi_{,i}(\be_i\cdot\dl\bx)
=(\hat\varphi_{,i}\be_i)\cdot\dl\bx.
\]
\emph{Step 3 (match with \eqr{1.114}).} Setting this equal to
$\grad\varphi\cdot\dl\bx$ gives
$(\grad\varphi-\hat\varphi_{,i}\be_i)\cdot\dl\bx=0$ for all $\dl\bx$.
Choose $\dl\bx=\grad\varphi-\hat\varphi_{,i}\be_i$ itself; then its squared
norm $|\grad\varphi-\hat\varphi_{,i}\be_i|^2=0$, so
\[
\grad\varphi=\hat\varphi_{,i}\be_i. \eqno(1.115)
\]
This is identical to (1.108), confirming that both routes give the same
formula. The Cartesian gradient has no extra factors because $\{\be_i\}$ is
fixed and orthonormal; when the basis varies with position (as in polar
coordinates), extra terms appear.
\end{example}

\begin{example}[cylindrical polar coordinates]\label{ex:polar}
With $\bk=\be_3$, define the rotating frame
\begin{equation}
\ber(\theta)=\cos\theta\,\be_1+\sin\theta\,\be_2,\qquad
\beth(\theta)=\bk\times\ber=-\sin\theta\,\be_1+\cos\theta\,\be_2, \tag{1.116}\label{eq:1.116}
\end{equation}
orthonormal for each $\theta$, with
\begin{equation}
\ber'(\theta)=\beth(\theta),\qquad \beth'(\theta)=-\ber(\theta). \tag{$\dagger$}
\end{equation}
A point has $\bx=r\,\ber(\theta)+z\,\bk$ (1.117). Then, using
$\dl\ber=\ber'(\theta)\dl\theta=\beth\,\dl\theta$,
\begin{equation}
\dl\bx=\dl r\,\ber+r\,\dl\ber+\dl z\,\bk
=\dl r\,\ber+r\,\dl\theta\,\beth+\dl z\,\bk, \tag{1.119}\label{eq:1.119}
\end{equation}
and dotting with the orthonormal $\{\ber,\beth,\bk\}$,
\begin{equation}
\dl r=\ber\cdot\dl\bx,\quad r\,\dl\theta=\beth\cdot\dl\bx,\quad
\dl z=\bk\cdot\dl\bx. \tag{1.120}\label{eq:1.120}
\end{equation}
With $\varphi=\tilde\varphi(r,\theta,z)$,
$\dl\tilde\varphi=\tilde\varphi_{,r}\dl r+\tilde\varphi_{,\theta}\dl\theta
+\tilde\varphi_{,z}\dl z$; substituting \eqr{1.120},
\begin{equation}
\dl\tilde\varphi=\big(\tilde\varphi_{,r}\ber
+\tfrac1r\tilde\varphi_{,\theta}\beth+\tilde\varphi_{,z}\bk\big)\cdot\dl\bx,
\tag{1.121}\label{eq:1.121}
\end{equation}
so by \eqr{1.114}, dropping the tilde,
\begin{equation}
\grad\varphi=\varphi_{,r}\,\ber+\tfrac1r\,\varphi_{,\theta}\,\beth
+\varphi_{,z}\,\bk. \tag{1.122}\label{eq:1.122}
\end{equation}
The extra $1/r$, absent from the Cartesian (1.108), arises solely because
$\{\ber,\beth,\bk\}$ varies with position while $\{\be_i\}$ is fixed.
\end{example}

\begin{example}[gradient of the radius $r=|\bx|$]\label{ex:gradr}
\emph{Index route.} Start from $r^2=\ip\bx\bx=x_jx_j$. Note that
$x_{j,i}=\partial x_j/\partial x_i=\dd_{ji}$: the Cartesian coordinate
$x_j$ depends only on itself, so its derivative with respect to $x_i$ is
$1$ when $j=i$ and $0$ otherwise. Differentiate $r^2=x_jx_j$ using the
product rule:
\begin{equation}
2r\,r_{,i}=(x_jx_j)_{,i}=x_{j,i}x_j+x_jx_{j,i}
=\dd_{ji}x_j+x_j\dd_{ji}=x_i+x_i=2x_i, \tag{1.123}\label{eq:1.123}
\end{equation}
where the Kronecker delta $\dd_{ji}$ collapsed the sum over $j$ to its
$j=i$ term: $\dd_{ji}x_j=x_i$. Divide both sides by $2r$:
$r_{,i}=x_i/r$. Since $\grad r=r_{,i}\be_i$ (by 1.108),
\begin{equation}
\grad r=\frac{x_i}{r}\be_i=\frac{1}{r}\bx. \tag{1.124}\label{eq:1.124}
\end{equation}

\emph{Coordinate free route.} The differential product rule for the
inner product reads $\dl(\ip\ba\bb)=\ip{\dl\ba}\bb+\ip\ba{\dl\bb}$
(proved by $\dl(a_ib_i)=\dl a_i\,b_i+a_i\dl b_i$). Apply it to
$r^2=\ip\bx\bx$:
\begin{equation}
2r\,\dl r=\dl(\ip\bx\bx)=\ip{\dl\bx}\bx+\ip\bx{\dl\bx}=2\,\ip\bx{\dl\bx}.
\tag{1.125}\label{eq:1.125}
\end{equation}
Cancel $2$ and divide by $r$: $\dl r=r^{-1}\ip\bx{\dl\bx}$. By \eqr{1.114},
$\dl r=\grad r\cdot\dl\bx$; matching the two gives $\grad r=r^{-1}\bx$.
Both routes arrive at the same answer, but the coordinate free route
never wrote a single index.
\end{example}

\begin{remark}[$\grad r$ points radially outward]\label{rem:gradr}
Since $\bx/r=\bx/|\bx|$ is the unit vector in the radial direction,
$\grad r$ is the unit radial vector: $|\grad r|=1$, pointing from the
origin toward the point $\bx$. The level sets of $r$ are concentric
spheres, and $\grad r$ is indeed perpendicular to them, consistent with
the general fact that the gradient is orthogonal to the level sets of
the field (Remark~\ref{rem:slope}).
\end{remark}

\subsection{Gradients of vector fields}\label{ss:gradv}

Repeat the linearisation of \S\ref{ss:grad} on a vector field. Now the
increment $\bv(\bx)-\bv(\bx_0)$ is a \emph{vector}, so its linear part is a
linear map $\Et\to\Et$, that is, a tensor. This is the double construction
one rung up: the fixed object representing the linear part is
$\grad\bv\in\Lin$.

\begin{definition}[Differentiable vector field]\label{def:diffv}
$\bv$ is differentiable at $\bx_0$ iff there is a tensor $\bL_0$ with
\begin{equation}
\bv(\bx)=\bv(\bx_0)+\bL_0(\bx-\bx_0)+\bo(\bx-\bx_0),\quad
|\bo(\bx-\bx_0)|=o(|\bx-\bx_0|), \tag{1.126--1.127}\label{eq:1.126}
\end{equation}
and one writes $\bL_0=\grad\bv\rest{\bx_0}$.
\end{definition}

\begin{theorem}[Cartesian representation]\label{thm:gradvcart}
$\grad\bv=v_{i,j}\,\be_i\otimes\be_j$, i.e. $(\grad\bv)_{ij}=v_{i,j}$.
\hfill\textup{(1.132)}
\end{theorem}

\begin{proof}
With $\bx-\bx_0=s\bu$, $|\bu|=1$,
$\bL_0\bu=\lim_{s\to0}\tfrac1s[\bv(\bx_0+s\bu)-\bv(\bx_0)]$ (1.129). Take
$\bu=\be_j$ and $\bv=\hat v_i\be_i$ with fixed $\be_i$:
$\bL_0\be_j=\be_i\lim_{s\to0}\tfrac1s[\hat v_i(\dots,x_j^0+s,\dots)
-\hat v_i(\dots)]=\hat v_{i,j}\be_i$ (1.130--1.131). Hence
$L_{ij}=\ip{\be_i}{\bL_0\be_j}=v_{i,j}$.
\end{proof}

\begin{corollary}[Chain rule for vector fields]\label{cor:dv}
$\dl\bv=(\grad\bv)\,\dl\bx$. \hfill\textup{(1.133)}
\end{corollary}

\begin{proof}
$\dl\hat v_i=\hat v_{i,j}\,\dl x_j=\hat v_{i,j}\,\be_j\cdot\dl\bx$, so
$\dl\bv=\dl\hat v_i\,\be_i=\hat v_{i,j}\,\be_i(\be_j\cdot\dl\bx)
=(\hat v_{i,j}\be_i\otimes\be_j)\dl\bx=(\grad\bv)\dl\bx$ (1.134).
\end{proof}

\begin{example}[$\bu=\mu\bx+\nu(\bx\otimes\bx)\ba$]\label{ex:gu1}
The field $\bu(\bx)=\mu\bx+\nu(\bx\otimes\bx)\ba$ arises, for instance,
in quadratic displacement fields. We compute its gradient by two routes.

\emph{Index route.} Write $u_i=\mu x_i+\nu x_ix_ja_j$ (since
$(\bx\otimes\bx)\ba=\ip\bx\ba\,\bx$, its $i$th component is
$x_i(x_ja_j)$). Differentiate with respect to $x_k$:
\begin{align}
u_{i,k}&=\mu\frac{\partial x_i}{\partial x_k}
+\nu\Big(\frac{\partial x_i}{\partial x_k}\,x_ja_j
+x_i\frac{\partial x_j}{\partial x_k}\,a_j\Big)\notag\\
&=\mu\dd_{ik}+\nu(\dd_{ik}\,x_ja_j+x_i\,\dd_{jk}\,a_j)\notag\\
&=\mu\dd_{ik}+\nu(\ip\bx\ba)\dd_{ik}+\nu x_ia_k\notag\\
&=[\mu+\nu(\ip\bx\ba)]\dd_{ik}+\nu x_ia_k. \tag{1.136}\label{eq:1.136}
\end{align}
Recognise the first term as $[\mu+\nu(\ip\bx\ba)]I_{ik}$ and the second
as $(\bx\otimes\ba)_{ik}=x_ia_k$, so
\begin{equation}
\grad\bu=[\mu+\nu(\ip\bx\ba)]\I+\nu\,\bx\otimes\ba. \tag{1.137}\label{eq:1.137}
\end{equation}

\emph{Differential route.} The product rule
$\dl(\bb\otimes\bc)=\dl\bb\otimes\bc+\bb\otimes\dl\bc$ (proved by
$\dl(b_ic_j)=\dl b_i\,c_j+b_i\,\dl c_j$) gives
$\dl(\bx\otimes\bx)=\dl\bx\otimes\bx+\bx\otimes\dl\bx$, so
\begin{align}
\dl\bu&=\mu\,\dl\bx+\nu(\dl\bx\otimes\bx+\bx\otimes\dl\bx)\ba\notag\\
&=\mu\,\dl\bx+\nu(\ip\bx\ba)\dl\bx+\nu\,\bx(\ip\ba{\dl\bx})\notag\\
&=\{[\mu+\nu(\ip\bx\ba)]\I+\nu\,\bx\otimes\ba\}\dl\bx. \tag{1.138}\label{eq:1.138}
\end{align}
In line two we used
$(\dl\bx\otimes\bx)\ba=(\ip\bx\ba)\dl\bx$ and
$(\bx\otimes\dl\bx)\ba=(\ip{\dl\bx}\ba)\bx=\bx(\ip\ba{\dl\bx})$.
Comparing with $\dl\bu=(\grad\bu)\dl\bx$ confirms \eqr{1.137}.
\end{example}

\begin{example}[affine field $\bu=\bx+(\ip\bb\bx)\ba$]\label{ex:gu2}
Here $\ba,\bb$ are fixed vectors. Write
$\bu(\bx)=\bx+\ba(\ip\bb\bx)$; its differential is
\[
\dl\bu=\dl\bx+\ba(\ip\bb{\dl\bx})=\dl\bx+\ba\otimes\bb\,\dl\bx
=(\I+\ba\otimes\bb)\dl\bx,
\]
where we recognised $\ba(\ip\bb{\dl\bx})=(\ba\otimes\bb)\dl\bx$ by
Definition~\ref{def:tp}. Thus
\begin{equation}
\grad\bu=\I+\ba\otimes\bb. \tag{1.141}\label{eq:1.141}
\end{equation}
Notice that $\grad\bu$ is \emph{uniform}: it does not depend on $\bx$.
This is no accident. Write $\bu(\bx)=(\I+\ba\otimes\bb)\bx=:\bA\bx$;
then for any $\bx_0$,
\[
\bu(\bx)=\bA\bx=\bA\bx_0+\bA(\bx-\bx_0)=\bu(\bx_0)+\bA(\bx-\bx_0),
\]
which is already in the form \eqr{1.127} with $\bL_0=\bA$ and \emph{zero
remainder}. Whenever the field is affine (linear in $\bx$), the gradient
is constant and the approximation \eqr{1.127} is exact everywhere.
\end{example}

\begin{example}[polar gradient of a vector field]\label{ex:gupolar}
For $\bu=u_r\ber+u_\theta\beth+u_z\bk$ with components functions of
$(r,\theta)$, using $(\dagger)$ and the product rule,
\[
\begin{aligned}
\dl\bu&=[u_{r,r}\dl r+(u_{r,\theta}-u_\theta)\dl\theta]\ber\\
&\quad+[u_{\theta,r}\dl r+(u_{\theta,\theta}+u_r)\dl\theta]\beth\\
&\quad+[u_{z,r}\dl r+u_{z,\theta}\dl\theta]\bk,
\end{aligned}
\]
where the $-u_\theta$ and $+u_r$ come from
$\dl\ber=\beth\dl\theta$, $\dl\beth=-\ber\dl\theta$. Inserting
$\dl r=\ber\cdot\dl\bx$, $\dl\theta=\tfrac1r\beth\cdot\dl\bx$ and
$(\be_a\cdot\dl\bx)\be_b=(\be_b\otimes\be_a)\dl\bx$,
\begin{equation}
\begin{aligned}
\grad\bu&=u_{r,r}\,\ber\otimes\ber
+\tfrac1r(u_{r,\theta}-u_\theta)\,\ber\otimes\beth
+u_{\theta,r}\,\beth\otimes\ber\\
&\quad+\tfrac1r(u_{\theta,\theta}+u_r)\,\beth\otimes\beth
+u_{z,r}\,\bk\otimes\ber+\tfrac1r u_{z,\theta}\,\bk\otimes\beth.
\end{aligned}\tag{1.145}\label{eq:1.145}
\end{equation}
\end{example}

\subsection{Divergence and curl}\label{ss:divcurl}

The tensor $\grad\bv$ carries nine numbers. Two coordinate free
combinations carry the physics: its trace measures the local rate of volume
change, the divergence, and its skew part measures local rotation, the
curl. Defining each through the trace and box product machinery keeps them
intrinsic, after which the component formulas drop out.

\begin{definition}[Divergence of a vector field]\label{def:divv}
$\dvg\bv:=\tr(\grad\bv)$.
\end{definition}

\begin{proposition}\label{prop:divv}
$\dvg\bv=v_{i,i}$.
\end{proposition}

\begin{proof}
$\dvg\bv=\tr(v_{i,j}\be_i\otimes\be_j)=v_{i,j}\dd_{ij}=v_{i,i}$ (1.147), by
$\tr(\be_i\otimes\be_j)=\dd_{ij}$.
\end{proof}

\begin{definition}[Divergence of a tensor field]\label{def:divA}
$\dvg\bA$ is the vector field with $\ip\bc{\dvg\bA}=\dvg(\bA^{t}\bc)$ for
every fixed $\bc$.
\end{definition}

\begin{proposition}\label{prop:divA}
$\dvg\bA=A_{ij,j}\,\be_i$. \hfill\textup{(1.150)}
\end{proposition}

\begin{proof}
Take $\bc=\be_i$: $\bA^{t}\be_i=A_{ij}\be_j$, so
$\ip{\be_i}{\dvg\bA}=\dvg(A_{ij}\be_j)=A_{ij,j}$ ($i$ fixed, $j$ summed).
\end{proof}

\begin{proposition}[Product rule]\label{prop:divprod}
$\ip\bu{\dvg\bA}=\dvg(\bA^{t}\bu)-\tr[\bA^{t}(\grad\bu)]$. \hfill\textup{(1.152)}
\end{proposition}

\begin{proof}
The scalar product rule $u_iA_{ij,j}=(A_{ij}u_i)_{,j}-A_{ij}u_{i,j}$
(1.151) reads, term by term, $\ip\bu{\dvg\bA}=\dvg(\bA^{t}\bu)
-\tr[\bA^{t}(\grad\bu)]$, since $(\bA^{t}\bu)_j=A_{ij}u_i$ and
$\tr[\bA^{t}\grad\bu]=A_{ij}u_{i,j}$.
\end{proof}

\begin{definition}[Curl of a vector field]\label{def:curlv}
$\crl\bv$ is the vector with $\ip\bc{\crl\bv}=\dvg(\bv\times\bc)$ for every
fixed $\bc$.
\end{definition}

\begin{proposition}\label{prop:curlv}
$\crl\bv=\eps_{kij}v_{i,k}\,\be_j$. \hfill\textup{(1.165)}
\end{proposition}

\begin{proof}
Take $\bc=\be_j$: $\bv\times\be_j=v_i\eps_{ijk}\be_k$, so
$\ip{\be_j}{\crl\bv}=\dvg(v_i\eps_{ijk}\be_k)=\eps_{ijk}v_{i,k}
=\eps_{kij}v_{i,k}$ (1.164), using $\eps_{ijk}=\eps_{kij}$.
\end{proof}

\begin{remark}[curl of a tensor field]\label{rem:curlA}
$\crl\bA$ is defined by $(\crl\bA)\bc=\crl(\bA^{t}\bc)$. Taking
$\bc=\be_m$, $\bA^{t}\be_m=A_{mi}\be_i$ has $i$th component $A_{mi}$, so
$\crl(\bA^{t}\be_m)=\eps_{kij}A_{mi,k}\be_j$; matching
$(\crl\bA)\be_m=(\crl\bA)_{jm}\be_j$ gives
$\crl\bA=\eps_{kij}A_{mi,k}\,\be_j\otimes\be_m$, equivalently
$\eps_{ilm}A_{jl,i}\,\be_m\otimes\be_j$ (Problem~28).
\end{remark}

\begin{example}[divergence of a tensor field in polar coordinates]\label{ex:divpolar}
For $\bA=A_{rr}\ber\otimes\ber+A_{r\theta}\ber\otimes\beth+\cdots$ with
components functions of $(r,\theta)$, decompose $\dvg\bA$ on
$\{\ber,\beth,\bk\}$ and use (1.152). From \eqr{1.145} applied to the frame
fields,
\begin{equation}
\grad\ber=\tfrac1r\,\beth\otimes\beth,\qquad
\grad\beth=-\tfrac1r\,\ber\otimes\beth, \tag{1.156}\label{eq:1.156}
\end{equation}
and the scalar divergence of $\bw=w_r\ber+w_\theta\beth+w_z\bk$ is
$\dvg\bw=w_{r,r}+\tfrac1r(w_{\theta,\theta}+w_r)$ (1.158). Applying (1.152)
with $\bA^{t}\ber=A_{rr}\ber+A_{r\theta}\beth+A_{rz}\bk$ and the like, and
$\tr[\bA^{t}\grad\ber]=\tfrac1rA_{\theta\theta}$,
$\tr[\bA^{t}\grad\beth]=-\tfrac1rA_{r\theta}$,
\begin{equation}
\begin{aligned}
\ber\cdot\dvg\bA&=A_{rr,r}+\tfrac1r(A_{rr}-A_{\theta\theta}+A_{r\theta,\theta}),\\
\beth\cdot\dvg\bA&=A_{\theta r,r}+\tfrac1r(A_{\theta r}+A_{r\theta}+A_{\theta\theta,\theta}),\\
\bk\cdot\dvg\bA&=A_{zr,r}+\tfrac1r(A_{zr}+A_{z\theta,\theta}).
\end{aligned}\tag{1.162}\label{eq:1.162}
\end{equation}
\end{example}

\subsection{Integral formulae}\label{ss:int}

Balance laws equate the rate of change of a quantity inside a region to
fluxes through its boundary. Passing between the two needs a bridge from
volume integrals of derivatives to surface integrals of values. That bridge
is the divergence theorem; every conservation law later rests on it.

\begin{theorem}[Divergence theorem]\label{thm:div}
For a $C^{1}$ vector field $\bv$ on a bounded region $R$ with piecewise
smooth orientable boundary $\partial R$ and exterior unit normal $\bn$,
\begin{equation}
\int_{R}\dvg\bv\,\dl v=\int_{\partial R}\ip\bv\bn\,\dl a. \tag{1.167}\label{eq:1.167}
\end{equation}
\end{theorem}

The proof is classical and is taken as given. Three consequences, used
constantly later, are immediate.

\begin{corollary}[Component, tensor and scalar forms]\label{cor:divforms}
\[\int_R v_{i,i}\,\dl v=\int_{\partial R}v_in_i\,\dl a;\]
$\int_R\dvg\bA\,\dl v=\int_{\partial R}\bA\bn\,\dl a$;
$\int_R\grad\varphi\,\dl v=\int_{\partial R}\varphi\,\bn\,\dl a$.
\end{corollary}

\begin{proof}
The first is \eqr{1.167} in components. For the tensor form, fix $\bc$ and
apply \eqr{1.167} to $\bA^{t}\bc$:
$\ip\bc{\int_R\dvg\bA\,\dl v}=\int_{\partial R}\ip{\bA^{t}\bc}\bn\,\dl a
=\ip\bc{\int_{\partial R}\bA\bn\,\dl a}$, true for all $\bc$. The scalar
form is the tensor form with $\bA=\varphi\I$, since
$\dvg(\varphi\I)=\grad\varphi$ and $\varphi\I\,\bn=\varphi\bn$.
\end{proof}

\begin{theorem}[Stokes, for reference]\label{thm:stokes}
For a differentiable $\bv$ on a piecewise smooth surface $\Omega$ bounded
by a piecewise smooth closed curve $\partial\Omega$,
\begin{equation}
\int_{\Omega}\ip{\crl\bv}\bn\,\dl a=\int_{\partial\Omega}\ip\bv{\dl\bx},
\tag{1.168}\label{eq:1.168}
\end{equation}
with $\bn$ oriented so that $\nu=\btt\times\bn$ is the outward normal to
$\partial\Omega$ in the tangent plane. \textup{(Stated for completeness;
not used in the sequel.)}
\end{theorem}

%=======================================================================
\section{Problems}\label{sec:problems}
%=======================================================================

All twenty eight problems of the chapter are solved below, in the book's
order, with every algebraic step shown. Two identities recur:
\begin{equation}
\ba\times(\bb\times\bc)=(\ip\ba\bc)\bb-(\ip\ba\bb)\bc,\qquad
\eps_{ijk}\eps_{imn}=\dd_{jm}\dd_{kn}-\dd_{jn}\dd_{km}. \tag{$\star$}
\end{equation}
Problems~1--23 use only the algebra of \S\ref{sec:vtt}; Problems~24--28 use
the field calculus of \S\ref{sec:fields}.

\begin{problem}[Cauchy--Schwarz, triangle, Pythagoras]
Prove $|\ip\ba\bb|\le|\ba||\bb|$, the triangle inequality, and the
Pythagorean theorem for orthogonal $\ba,\bb$.
\end{problem}
\begin{solution}
\emph{Cauchy--Schwarz.} If $\ba=\bze$ then $\ip\ba\bb=0$ and the inequality
reads $0\le0$; assume $\ba\neq\bze$. Following the hint, form the vector
\[
\bw=(\ip\ba\ba)\bb-(\ip\ba\bb)\ba .
\]
Positivity of the norm gives $0\le|\bw|^2=\ip\bw\bw$. Expand the dot product
by bilinearity, keeping every one of the four cross terms (write
$\alpha=\ip\ba\ba$, $\beta=\ip\ba\bb$ for brevity):
\[
\begin{aligned}
|\bw|^2&=\ip{(\alpha\bb-\beta\ba)}{(\alpha\bb-\beta\ba)}\\
&=\alpha^2\,\ip\bb\bb-\alpha\beta\,\ip\bb\ba-\beta\alpha\,\ip\ba\bb
 +\beta^2\,\ip\ba\ba\\
&=\alpha^2\,\ip\bb\bb-2\alpha\beta^2+\beta^2\alpha
 \qquad(\text{using }\ip\bb\ba=\ip\ba\bb=\beta,\ \ip\ba\ba=\alpha)\\
&=\alpha\big[\alpha\,\ip\bb\bb-\beta^2\big]
 =(\ip\ba\ba)\big[(\ip\ba\ba)(\ip\bb\bb)-(\ip\ba\bb)^2\big].
\end{aligned}
\]
Since $\ip\ba\ba=|\ba|^2>0$, dividing the inequality $|\bw|^2\ge0$ by it
leaves the bracket nonnegative:
\[
(\ip\ba\ba)(\ip\bb\bb)-(\ip\ba\bb)^2\ge0
\ \Longrightarrow\ (\ip\ba\bb)^2\le|\ba|^2|\bb|^2 .
\]
Taking the (nonnegative) square root, $|\ip\ba\bb|\le|\ba|\,|\bb|$. Equality
requires $\bw=\bze$, i.e.\ $\ba,\bb$ parallel.

\emph{Triangle inequality.} Expand the squared norm and insert
Cauchy--Schwarz at the marked step, $\ip\ba\bb\le|\ip\ba\bb|\le|\ba||\bb|$:
\[
\begin{aligned}
|\ba+\bb|^2&=\ip{(\ba+\bb)}{(\ba+\bb)}=|\ba|^2+2\,\ip\ba\bb+|\bb|^2\\
&\le|\ba|^2+2|\ba||\bb|+|\bb|^2=(|\ba|+|\bb|)^2 .
\end{aligned}
\]
Both sides are nonnegative, so $|\ba+\bb|\le|\ba|+|\bb|$.

\emph{Pythagoras.} If $\ba\perp\bb$, i.e.\ $\ip\ba\bb=0$, the middle term
above drops out and $|\ba+\bb|^2=|\ba|^2+|\bb|^2$.
\end{solution}

\begin{remark}[what the hint vector is doing]\label{rem:csuse}
The vector $\bw=(\ip\ba\ba)\bb-(\ip\ba\bb)\ba$ is (a positive multiple of)
the part of $\bb$ \emph{perpendicular} to $\ba$: subtracting the projection
$(\ip\ba\bb/\ip\ba\ba)\ba$ from $\bb$ and scaling by $\ip\ba\ba$ gives
exactly $\bw$. Cauchy--Schwarz then says nothing more than ``a squared
length is never negative'': the inequality is tight precisely when that
perpendicular part is zero, i.e.\ when $\bb$ is parallel to $\ba$. This is
why Cauchy--Schwarz underlies the very definition of angle
(Definition~\ref{def:angle}): $\cos\theta=\ip\ba\bb/(|\ba||\bb|)$ lands in
$[-1,1]$ exactly because $|\bw|^2\ge0$.
\end{remark}

\begin{problem}[Is the norm linear?]
Is $f(\bv)=|\bv|$ linear?
\end{problem}
\begin{solution}
No. Linearity requires both $f(\bu+\bv)=f(\bu)+f(\bv)$ and
$f(\alpha\bv)=\alpha f(\bv)$; each fails.

\emph{Additivity.} Take $\bv=-\bu$ with $\bu\neq\bze$:
\[
   f(\bu+\bv)=f(\bze)=0,\qquad
   f(\bu)+f(\bv)=|\bu|+|-\bu|=2|\bu|>0 .
\]

\emph{Homogeneity.} For $\alpha<0$,
\[
   f(\alpha\bv)=|\alpha\bv|=|\alpha|\,|\bv|=-\alpha|\bv|
   \neq\alpha|\bv|=\alpha f(\bv)
\]
whenever $\bv\neq\bze$.

What the norm does satisfy is $f(\alpha\bv)=|\alpha|f(\bv)$ and
$f(\bu+\bv)\le f(\bu)+f(\bv)$, the triangle inequality of Problem 1. Those
two make it a norm, not a linear function.
\end{solution}

\begin{problem}[A constant map is not a tensor]
For a fixed $\bmv$, is there a tensor $\bA$ with $\bA\bv=\bmv$ for all
$\bv$?
\end{problem}
\begin{solution}
No, unless $\bmv=\bze$. A linear map fixes the origin:
\[
   \bA\bze=\bA(0\cdot\bze)=0\cdot\bA\bze=\bze ,
\]
whereas the requirement gives $\bA\bze=\bmv$. So $\bmv=\bze$.

Second route, without using $\bze$. Homogeneity with $\alpha=2$ gives
\[
   \bmv=\bA(2\bv)=2\,\bA\bv=2\bmv
   \ \Longrightarrow\ \bmv=\bze .
\]
In the problem $\bmv=\tfrac{\sqrt2}{2}(\be_{1}+\be_{2})\neq\bze$, so no
such tensor exists. Contrast Problem 4, where $\bA\bv=(\ip\bv\bn)\bmv$
\emph{is} linear: the difference is the factor $\ip\bv\bn$, which vanishes
with $\bv$.
\end{solution}

\begin{problem}[A rank-one tensor]
With $\bmv=\tfrac{\sqrt2}{2}(\be_1+\be_2)$,
$\bn=\tfrac{\sqrt2}{2}(\be_3-\be_1)$, find $\bA$ with
$\bA\bv=(\ip\bv\bn)\bmv$.
\end{problem}
\begin{solution}
The map $\bv\mapsto(\ip\bv\bn)\bmv$ is linear in $\bv$ (the inner product
is linear in its second argument) and returns a vector, so it is a tensor.
By Definition~\ref{def:tp}, $(\bmv\otimes\bn)\bv=(\ip\bn\bv)\bmv$. Hence
$\bA=\bmv\otimes\bn$.

\emph{Components.} $A_{ij}=m_in_j$ (from Example~\ref{ex:tp}). The
components of $\bmv$ are
$m_1=\tfrac{\sqrt2}{2}$, $m_2=\tfrac{\sqrt2}{2}$, $m_3=0$;
those of $\bn$ are $n_1=-\tfrac{\sqrt2}{2}$, $n_2=0$, $n_3=\tfrac{\sqrt2}{2}$.
Each entry $A_{ij}=m_in_j$:
\[
\begin{aligned}
A_{11}&=\tfrac{\sqrt2}{2}\cdot(-\tfrac{\sqrt2}{2})=-\tfrac12,\quad
A_{12}=\tfrac{\sqrt2}{2}\cdot0=0,\quad
A_{13}=\tfrac{\sqrt2}{2}\cdot\tfrac{\sqrt2}{2}=\tfrac12,\\
A_{21}&=-\tfrac12,\quad A_{22}=0,\quad A_{23}=\tfrac12,\quad
A_{31}=0,\quad A_{32}=0,\quad A_{33}=0.
\end{aligned}
\]
Hence
\[
(A_{ij})=\tfrac12\begin{pmatrix}-1&0&1\\-1&0&1\\0&0&0\end{pmatrix}.
\]
\emph{Check.} $\bA$ has rank $1$ (all rows are multiples of $(-1,0,1)$),
confirming it is a rank one tensor: the image of $\bA$ is the line
spanned by $\bmv$.
\end{solution}

\begin{problem}[$\eps_{ijk}A_{jk}=0\iff\bA=\bA^{t}$]
\end{problem}
\begin{solution}
Write $c_i:=\eps_{ijk}A_{jk}$ (a vector, since $i$ is free while $j,k$ are
summed). We must show $c_i=0$ for all $i$ exactly when $A_{jk}=A_{kj}$.

\emph{($\Leftarrow$) Symmetric $\Rightarrow c_i=0$.} Suppose $A_{jk}=A_{kj}$.
In $c_i=\eps_{ijk}A_{jk}$ rename the two summed indices $j\leftrightarrow k$
(a dummy relabelling changes nothing):
\[
c_i=\eps_{ijk}A_{jk}=\eps_{ikj}A_{kj}.
\]
Now use symmetry $A_{kj}=A_{jk}$ and the antisymmetry of the permutation
symbol under swapping its last two slots, $\eps_{ikj}=-\eps_{ijk}$:
\[
c_i=\eps_{ikj}A_{kj}=-\eps_{ijk}A_{jk}=-c_i .
\]
Thus $c_i=-c_i$, i.e.\ $2c_i=0$, so $c_i=0$ for every $i$.

\emph{($\Rightarrow$) $c_i=0\Rightarrow$ symmetric.} Split $\bA$ into its
symmetric and skew parts, $A_{jk}=S_{jk}+W_{jk}$ with
$S_{jk}=\tfrac12(A_{jk}+A_{kj})$, $W_{jk}=\tfrac12(A_{jk}-A_{kj})$. By the
direction just proved, $\eps_{ijk}S_{jk}=0$, so
\[
0=c_i=\eps_{ijk}A_{jk}=\eps_{ijk}S_{jk}+\eps_{ijk}W_{jk}=\eps_{ijk}W_{jk}.
\]
Every skew array is $W_{jk}=-\eps_{jkl}w_l$ for a unique axial vector $\bw$
(this is the general solution of $W_{jk}=-W_{kj}$; see
Proposition~\ref{prop:axial}). Substitute and contract:
\[
\begin{aligned}
\eps_{ijk}W_{jk}&=\eps_{ijk}(-\eps_{jkl}w_l)=-\eps_{ijk}\eps_{jkl}\,w_l\\
&=-\eps_{jki}\eps_{jkl}\,w_l=-2\dd_{il}w_l=-2w_i,
\end{aligned}
\]
where we wrote $\eps_{ijk}=\eps_{jki}$ (cyclic shift) and used the contracted
identity $\eps_{jki}\eps_{jkl}=2\dd_{il}$. Hence $0=c_i=-2w_i$ gives
$w_i=0$ for all $i$, so $\bw=\bze$, so $W_{jk}=0$, so $A_{jk}=S_{jk}=A_{kj}$:
$\bA$ is symmetric.
\end{solution}

\begin{remark}[what the operation $A_{jk}\mapsto\eps_{ijk}A_{jk}$ extracts]\label{rem:p5use}
The computation shows $\eps_{ijk}A_{jk}=-2w_i$, twice the (negated) axial
vector of the skew part of $\bA$. So contracting a tensor with the
permutation symbol on two indices is a machine that \emph{throws away the
symmetric part and reads off the skew part as a vector}. This is exactly the
device behind writing the vorticity of a velocity field, or the moment of a
force, as a single axial vector; it works only in three dimensions, where a
skew tensor (three independent entries) matches a vector (three
components).
\end{remark}

\begin{problem}[Skew/symmetric and the tensor inner product]
With converses: (a) $\bA^{t}=-\bA\Rightarrow\bA\cdot(\ba\otimes\ba)=0$;
(b) $\bT\cdot\bS=0\ \forall\bS\Rightarrow\bT=\Ob$;
(c) $\bT\cdot\bS=0\ \forall\,\bS=\bS^t\Rightarrow\bT$ skew;
(d) $\bT\cdot\bS=0\ \forall\,\bS=-\bS^t\Rightarrow\bT$ symmetric.
\end{problem}
\begin{solution}
\emph{Preliminary.} Recall: $\bA\cdot\bB=\tr(\bA\bB^t)=A_{ij}B_{ij}$
(sum of all nine component products), and
$\bA\cdot(\bu\otimes\bv)=\ip\bu{\bA\bv}$.

\emph{Lemma (the key orthogonality).}
Let $\bS=\bS^t$ be symmetric and $\bW=-\bW^t$ skew. Then
$\bS\cdot\bW=S_{ij}W_{ij}$. Relabelling $i\leftrightarrow j$:
$S_{ij}W_{ij}=S_{ji}W_{ji}=S_{ij}(-W_{ij})=-S_{ij}W_{ij}$, using
$S_{ji}=S_{ij}$ and $W_{ji}=-W_{ij}$. Thus $2(\bS\cdot\bW)=0$, i.e.
$\bS\cdot\bW=0$: a symmetric tensor is orthogonal to every skew tensor.

(a) $\bA\cdot(\ba\otimes\ba)=\ip\ba{\bA\ba}$. Since $\bA^t=-\bA$,
$\ip\ba{\bA\ba}=\ip{\bA^t\ba}\ba=\ip{-\bA\ba}\ba=-\ip\ba{\bA\ba}$,
hence $\ip\ba{\bA\ba}=0$. \emph{Converse:} if
$\ip\ba{\bA\ba}=0$ for all $\ba$, put $\ba=\be_i+\be_j$ ($i\neq j$):
$\ip{(\be_i+\be_j)}{\bA(\be_i+\be_j)}=A_{ii}+A_{ij}+A_{ji}+A_{jj}=0$.
With $\ba=\be_i$: $A_{ii}=0$; with $\ba=\be_j$: $A_{jj}=0$. Hence
$A_{ij}+A_{ji}=0$ for all $i\neq j$, together with $A_{ii}=0$: $\bA$ is
skew.

(b) $\bT\cdot\bS=0$ for all $\bS$: choose $\bS=\bT$, getting
$|\bT|^2=T_{ij}T_{ij}=0$, so every $T_{ij}=0$, i.e. $\bT=\Ob$.

(c) $\bT\cdot\bS=0$ for all symmetric $\bS$. Decompose
$\bT=\bT_s+\bT_w$ (symmetric + skew). Choose $\bS=\bT_s$: by the Lemma,
$\bT_w\cdot\bT_s=0$, so $\bT\cdot\bT_s=\bT_s\cdot\bT_s=|\bT_s|^2=0$,
hence $\bT_s=\Ob$ and $\bT=\bT_w$ is skew.

(d) $\bT\cdot\bS=0$ for all skew $\bS$. Choose $\bS=\bT_w$: by the Lemma,
$\bT_s\cdot\bT_w=0$, so $\bT\cdot\bT_w=|\bT_w|^2=0$, hence $\bT_w=\Ob$
and $\bT=\bT_s$ is symmetric.
\end{solution}

\begin{problem}[Is $\bA(\bu\times\bv)=\bA\bu\times\bv+\bu\times\bA\bv$?]
\end{problem}
\begin{solution}
Write $D(\bA):=\bA(\bu\times\bv)-\bA\bu\times\bv-\bu\times\bA\bv$. The
question is when $D(\bA)=\bze$ for all $\bu,\bv$. Note $D$ is linear in
$\bA$.

\emph{No in general.} Take $\bA=\I$:
\[
   D(\I)=\bu\times\bv-\bu\times\bv-\bu\times\bv=-\,\bu\times\bv\neq\bze .
\]

\emph{Skew tensors do satisfy it.} Let $\bW\bx=\bom\times\bx$, Problem 16.
The Jacobi identity is
\[
   \bom\times(\bu\times\bv)+\bu\times(\bv\times\bom)
   +\bv\times(\bom\times\bu)=\bze .
\]
Use $\bv\times(\bom\times\bu)=-\,(\bom\times\bu)\times\bv$ and
$\bu\times(\bv\times\bom)=-\,\bu\times(\bom\times\bv)$:
\[
   \bom\times(\bu\times\bv)=(\bom\times\bu)\times\bv
   +\bu\times(\bom\times\bv),
\]
that is $D(\bW)=\bze$.

\emph{Only skew tensors do.} Suppose $D(\bA)=\bze$. Split
$\bA=\bS+\bW$ by Problem 14; $D(\bW)=\bze$, so $D(\bS)=\bze$ by
linearity. Take the principal basis of $\bS$, $\bS\be_{i}=\lambda_{i}
\be_{i}$, and put $\bu=\be_{i}$, $\bv=\be_{j}$ with $ijk$ cyclic, so
$\be_{i}\times\be_{j}=\be_{k}$:
\[
   \lambda_{k}\be_{k}
   =\lambda_{i}\be_{i}\times\be_{j}+\be_{i}\times\lambda_{j}\be_{j}
   =(\lambda_{i}+\lambda_{j})\be_{k}
   \ \Longrightarrow\
   \lambda_{k}=\lambda_{i}+\lambda_{j}.
\]
The three cyclic triples give
\[
   \lambda_{3}=\lambda_{1}+\lambda_{2},\qquad
   \lambda_{1}=\lambda_{2}+\lambda_{3},\qquad
   \lambda_{2}=\lambda_{3}+\lambda_{1}.
\]
Adding, $\sum\lambda_{i}=2\sum\lambda_{i}$, so $\sum\lambda_{i}=0$; then
the first equation gives $\lambda_{3}=-\lambda_{3}$, hence
$\lambda_{3}=0$, and likewise $\lambda_{1}=\lambda_{2}=0$. So $\bS=\Ob$
and $\bA=\bW$ is skew.
\end{solution}

\begin{remark}[what the answer says]\label{rem:P7derivation}
$D(\bA)=\bze$ is the Leibniz rule for $\bA$ acting on the product
$\bu\times\bv$. The result is therefore: \emph{a tensor differentiates the
cross product exactly when it is skew}, i.e.\ when it is an infinitesimal
rotation. This is the statement that $\Skw$ is the Lie algebra of the
rotation group, and it recurs in Chapter 4, where the spin $\bW$ is skew
for precisely this reason.
\end{remark}

\begin{problem}[Reciprocal bases]
With $\be_i=a_{ij}\be'_j$, $\be'_i=a'_{ij}\be_j$, show (a)
$a_{ik}a'_{kj}=\dd_{ij}$; (b) $(a')=(a)^{t}$.
\end{problem}
\begin{solution}
\textbf{(a)} Substitute the second relation into the first:
\[
   \be_i=a_{ij}\be'_j=a_{ij}\big(a'_{jk}\be_k\big)=a_{ij}a'_{jk}\,\be_k .
\]
But also $\be_i=\dd_{ik}\be_k$. The $\{\be_k\}$ are independent, so the
coefficients match:
\[
   a_{ij}a'_{jk}=\dd_{ik}.
\]

\textbf{(b)} Dot $\be_i=a_{ij}\be'_j$ with $\be'_k$ and use
$\ip{\be'_j}{\be'_k}=\dd_{jk}$:
\[
   \ip{\be_i}{\be'_k}=a_{ij}\dd_{jk}=a_{ik}.
\]
Dot $\be'_i=a'_{ij}\be_j$ with $\be_k$ and use $\ip{\be_j}{\be_k}=\dd_{jk}$:
\[
   \ip{\be'_i}{\be_k}=a'_{ij}\dd_{jk}=a'_{ik}.
\]
The two left sides are the same number, since the dot product is
symmetric:
\[
   a'_{ik}=\ip{\be'_i}{\be_k}=\ip{\be_k}{\be'_i}=a_{ki},
   \qquad\text{i.e.}\qquad (a')=(a)^{t}.
\]
Part (a) then reads $a_{ij}a_{kj}=\dd_{ik}$: the matrix $(a)$ is
orthogonal, $(a)(a)^{t}=(a)^{t}(a)=(\dd)$. This is the fact used in
Problem 9 and again in Problem 3 of Chapter 2.
\end{solution}

\begin{problem}[Invariance of trace and dot product; a computation]
(a) $A_{ii}=A'_{ii}$; (b) deduce invariance of $\ip\bu\bv$; (c) compute
$A'_{ij}$ for the given data.
\end{problem}
\begin{solution}
\textbf{(a)} With $a_{ij}=\ip{\be_i}{\be'_j}$, so that
$\be'_j=a_{kj}\be_k$, the components transform as
\[
   A'_{ij}=\ip{\be'_i}{\bA\be'_j}
   =\ip{a_{ki}\be_k}{\bA\,a_{lj}\be_l}
   =a_{ki}a_{lj}\ip{\be_k}{\bA\be_l}
   =a_{ki}a_{lj}A_{kl}.
\]
Orthonormality of the primed basis fixes $(a)$:
\[
   \dd_{ij}=\ip{\be'_i}{\be'_j}=a_{ki}a_{lj}\,\ip{\be_k}{\be_l}
   =a_{ki}a_{lj}\dd_{kl}=a_{ki}a_{kj},
\]
so $(a)$ is orthogonal. Set $j=i$ and sum:
\[
   A'_{ii}=a_{ki}a_{li}A_{kl}=\dd_{kl}A_{kl}=A_{kk}.
\]

\textbf{(b)} Vectors carry one factor, $u'_i=\ip{\be'_i}{\bu}
=a_{ki}u_k$. Hence
\[
   u'_iv'_i=a_{ki}a_{li}\,u_kv_l=\dd_{kl}u_kv_l=u_kv_k .
\]
Alternatively, apply (a) to $\bA=\bu\otimes\bv$, whose trace is
$\ip\bu\bv$ by Theorem~\ref{thm:trace}.

(c) The data are
\[
(A_{ij})=\begin{pmatrix}1&5&-5\\5&0&0\\-5&0&1\end{pmatrix},\qquad
\be'_1=\tfrac1{\sqrt5}(-1,0,2),\qquad \be'_2=\be_2=(0,1,0).
\]
Fix $\be'_3$ by right-handedness, $\be'_3=\be'_1\times\be'_2$. The cross
product $(-1,0,2)\times(0,1,0)$ has components
\[
\big(0\cdot0-2\cdot1,\ \ 2\cdot0-(-1)\cdot0,\ \ (-1)\cdot1-0\cdot0\big)
=(-2,0,-1),
\]
so $\be'_3=\tfrac1{\sqrt5}(-2,0,-1)$. The change-of-basis matrix
$\bQ=(a_{ij})$, $a_{ij}=\ip{\be_i}{\be'_j}$, carries the components of
$\be'_j$ in its $j$th column:
\[
\bQ=\tfrac1{\sqrt5}\begin{pmatrix}-1&0&-2\\0&\sqrt5&0\\2&0&-1\end{pmatrix}.
\]
The law $A'_{ij}=a_{ki}a_{lj}A_{kl}$ is $\bA'=\bQ^{t}\bA\,\bQ$, which we read
as $A'_{ij}=\ip{\be'_i}{\bA\be'_j}$. Compute the three columns $\bA\be'_j$
(matrix times vector); for example
\[
\bA\be'_1=\tfrac1{\sqrt5}\bA\!\begin{pmatrix}-1\\0\\2\end{pmatrix}
=\tfrac1{\sqrt5}\begin{pmatrix}1(-1)+5(0)-5(2)\\5(-1)+0+0\\-5(-1)+0+1(2)\end{pmatrix}
=\tfrac1{\sqrt5}\begin{pmatrix}-11\\-5\\7\end{pmatrix},
\]
and likewise
$\bA\be'_2=(5,0,0)^t$,
$\bA\be'_3=\tfrac1{\sqrt5}(3,-10,9)^t$. Now
$A'_{ij}=\ip{\be'_i}{\bA\be'_j}$ entry by entry:
\[
\begin{aligned}
A'_{11}&=\tfrac1{\sqrt5}(-1,0,2)\!\cdot\!\tfrac1{\sqrt5}(-11,-5,7)
=\tfrac{11+0+14}{5}=5,\\
A'_{12}&=\tfrac1{\sqrt5}(-1,0,2)\!\cdot\!(5,0,0)=-\tfrac{5}{\sqrt5}=-\sqrt5,\\
A'_{13}&=\tfrac1{\sqrt5}(-1,0,2)\!\cdot\!\tfrac1{\sqrt5}(3,-10,9)
=\tfrac{-3+0+18}{5}=3,\\
A'_{22}&=(0,1,0)\!\cdot\!(5,0,0)=0,\\
A'_{23}&=(0,1,0)\!\cdot\!\tfrac1{\sqrt5}(3,-10,9)=-\tfrac{10}{\sqrt5}=-2\sqrt5,\\
A'_{33}&=\tfrac1{\sqrt5}(-2,0,-1)\!\cdot\!\tfrac1{\sqrt5}(3,-10,9)
=\tfrac{-6+0-9}{5}=-3,
\end{aligned}
\]
with $A'_{21}=A'_{12}$, $A'_{31}=A'_{13}$, $A'_{32}=A'_{23}$ since $\bA'$ is
symmetric ($\bA$ is). Therefore
\[
\bA'=\begin{pmatrix}5&-\sqrt5&3\\-\sqrt5&0&-2\sqrt5\\3&-2\sqrt5&-3\end{pmatrix}.
\]
Check: $\tr\bA'=5+0-3=2=1+0+1=\tr\bA$, confirming the invariance proved in
part (a).
\end{solution}

\begin{problem}[Projection and reflection tensors]
\textup{(a)} Consider a plane $P$ with unit normal $\bn$. Let $\bv\in\Et$ be
arbitrary and $\bv_p$ its projection onto $P$. Verify
$\bv=\bv_p+(\ip\bv\bn)\bn$ and find $\Proj{\bn}$ in
$\bv_p=\Proj{\bn}\bv$.
\textup{(b)} A tensor $\bA$ is a \emph{projection} iff $\bA^{t}=\bA$ and
$\bA^{2}=\bA$. Show $\Proj{\bn}$ is a projection, and that
$\Ob$, $\I$, $\bn\otimes\bn$ are projections.
\textup{(c)} Let $\Refl{\bn}$ be the \emph{reflection} tensor for $P$. Find
$\Refl{\bn}$ and its components relative to $\{\be_i\otimes\be_j\}$ for
\textup{(i)} $\bn=\be_1$, \textup{(ii)}
$\bn=\tfrac{\sqrt2}{2}(\be_1+\be_2)$.
\end{problem}

\begin{solution}
The normal is a unit vector and the plane through the origin with that
normal consists of the vectors orthogonal to it, so throughout
\begin{equation}
   \ip\bn\bn=1,\qquad
   P=\{\bw\in\Et:\ \ip\bw\bn=0\}=\{\bn\}^{\perp},\qquad
   \mathcal{N}=\operatorname{span}\{\bn\}, \tag{$\star_1$}
\end{equation}
$P$ being two dimensional and $\mathcal{N}$ one dimensional.

\emph{(a)} The shadow of $\bv$ on the direction $\bn$ has signed length
$\ip\bv\bn$, so the part of $\bv$ along the normal, and what is left after
removing it, are
\begin{equation}
   \bv_n:=(\ip\bv\bn)\bn,\qquad \bv_p:=\bv-(\ip\bv\bn)\bn . \tag{$\star_2$}
\end{equation}
That $\bv_p$ lies in $P$ is not an assumption but a computation; dot it with
$\bn$ and use bilinearity together with $\ip\bn\bn=1$:
\[
   \ip{\bv_p}{\bn}=\ip{\big[\bv-(\ip\bv\bn)\bn\big]}{\bn}
   =\ip\bv\bn-(\ip\bv\bn)\,\ip\bn\bn
   \overset{(\star_1)}{=}\ip\bv\bn-\ip\bv\bn=0 ,
\]
hence $\bv_p\in\{\bn\}^{\perp}=P$, and by construction $\bv=\bv_p+\bv_n$,
which is the required identity.

The splitting is unique. Suppose $\bv=\bw_p+\alpha\bn$ with $\bw_p\in P$ and
$\alpha\in\Rn$. Subtracting the two expressions for $\bv$,
\[
   \bv_p-\bw_p=(\alpha-\ip\bv\bn)\bn ,
\]
whose left side lies in $P$ and whose right side lies in $\mathcal{N}$. A
vector of $P\cap\mathcal{N}$ satisfies both $\bw=\lambda\bn$ and
$\ip\bw\bn=0$, so $\lambda=\lambda\,\ip\bn\bn=\ip\bw\bn=0$ and
$P\cap\mathcal{N}=\{\bze\}$. Therefore $\bw_p=\bv_p$ and
$\alpha=\ip\bv\bn$, and $\Et=P\oplus\mathcal{N}$.

Before writing $\bv_p=\Proj{\bn}\bv$ we must know the map $\bv\mapsto\bv_p$
is linear, since only linear maps are tensors:
\[
\begin{aligned}
   (\alpha\bu+\beta\bv)_p
   &=\alpha\bu+\beta\bv-\big(\ip{(\alpha\bu+\beta\bv)}{\bn}\big)\bn\\
   &=\alpha\bu+\beta\bv-\big(\alpha\,\ip\bu\bn+\beta\,\ip\bv\bn\big)\bn\\
   &=\alpha\big[\bu-(\ip\bu\bn)\bn\big]+\beta\big[\bv-(\ip\bv\bn)\bn\big]
   =\alpha\bu_p+\beta\bv_p .
\end{aligned}
\]
Now recognise the subtracted term as a tensor acting on $\bv$. By the
definition of the tensor product, $(\bn\otimes\bn)\bv=(\ip\bn\bv)\bn
=(\ip\bv\bn)\bn$, so $(\star_2)$ becomes
$\bv_p=\bv-(\bn\otimes\bn)\bv=[\I-\bn\otimes\bn]\bv$ for every $\bv$, and
equality of tensors is equality of their action. Hence
\begin{equation}
   \Proj{\bn}=\I-\bn\otimes\bn,\qquad
   (\Proj{\bn})_{ij}=\dd_{ij}-n_in_j , \tag{$\star_3$}
\end{equation}
the components following from $\I=\dd_{ij}\be_i\otimes\be_j$ and
$(\bn\otimes\bn)_{ij}=n_in_j$.


\begin{figure}[htbp]\centering
\begin{tikzpicture}[scale=1.05]
  \fill[gray!12] (-2.5,-0.55) -- (2.5,-0.55) -- (3.3,0.5) -- (-1.7,0.5) -- cycle;
  \draw[gray!65,thin] (-2.5,-0.55) -- (2.5,-0.55) -- (3.3,0.5) -- (-1.7,0.5) -- cycle;
  \node[gray!75] at (-2.15,-0.28) {$P$};
  \coordinate (O) at (0,0); \fill (O) circle (1.4pt);
  \draw[vec] (O) -- (0,0.62); \node[left=2pt] at (0,0.38) {$\bn$};
  \draw[gray!70,thin] (-0.17,0) -- (-0.17,0.17) -- (0,0.17);
  \coordinate (Q) at (1.95,0.30);
  \coordinate (V) at (1.95,1.45);
  \draw[vec] (O) -- (V); \node[above left=-1pt] at (V) {$\bv$};
  \draw[vec,black] (O) -- (Q);
  \node[below=3pt] at (1.35,0.24) {$\bv_p=\Proj{\bn}\bv$};
  \draw[vec,black] (Q) -- (V);
  \node[right=3pt] at (1.98,0.88) {$(\ip\bv\bn)\bn=(\bn\otimes\bn)\bv$};
  \draw[gray!70,thin] (Q) ++(-0.19,0.02) -- ++(0.02,0.19) -- ++(0.19,-0.02);
\end{tikzpicture}
\\[1.2ex]
{\footnotesize\textsc{Figure P10.1.} $\bv=\bv_p+(\ip\bv\bn)\bn$ with
$\ip{\bv_p}\bn=0$; $\Proj{\bn}$ returns the leg in $P$, $\bn\otimes\bn$ the
leg along $\bn$, and $\Proj{\bn}+\bn\otimes\bn=\I$.}
\end{figure}


\emph{(b)} Everything rests on one property of the dyad, obtained by
applying it twice to an arbitrary $\bv$ and unfolding the definition at each
step:
\begin{equation}
   (\bn\otimes\bn)^{2}\bv=(\bn\otimes\bn)\big[(\ip\bn\bv)\bn\big]
   =(\ip\bn\bv)(\ip\bn\bn)\bn
   \overset{(\star_1)}{=}(\ip\bn\bv)\bn=(\bn\otimes\bn)\bv , \tag{$\star_4$}
\end{equation}
true for every $\bv$, so $(\bn\otimes\bn)^{2}=\bn\otimes\bn$. Note where the
hypothesis enters: the middle step used $\ip\bn\bn=1$, and for a normal of
length other than one the dyad would not be idempotent. Its symmetry follows
from the definition of the transpose, since for all $\bu,\bv$
\[
   \ip\bu{(\bn\otimes\bn)\bv}=\ip{\bu}{(\ip\bn\bv)\bn}
   =(\ip\bn\bv)(\ip\bu\bn)=\ip\bv{(\ip\bn\bu)\bn}
   =\ip\bv{(\bn\otimes\bn)\bu},
\]
which says $(\bn\otimes\bn)^{t}=\bn\otimes\bn$. Transposition being linear,
and $\I^{t}=\I$,
\[
   \Proj{\bn}^{t}=\I^{t}-(\bn\otimes\bn)^{t}=\I-\bn\otimes\bn=\Proj{\bn},
\]
and expanding the square, the two equal cross terms combining into one,
\[
   \Proj{\bn}^{2}=\I-\bn\otimes\bn-\bn\otimes\bn+(\bn\otimes\bn)^{2}
   =\I-2\,\bn\otimes\bn+(\bn\otimes\bn)^{2}
   \overset{(\star_4)}{=}\I-\bn\otimes\bn=\Proj{\bn}.
\]
The same two facts in components, which is the form used in calculations;
recall $(\bA\bB)_{ij}=A_{ik}B_{kj}$, and that a Kronecker delta collapses
the sum it is contracted with, $\dd_{ik}n_k=n_i$ and
$\dd_{ik}\dd_{kj}=\dd_{ij}$:
\[
\begin{aligned}
   \big[(\bn\otimes\bn)^{2}\big]_{ij}&=(n_in_k)(n_kn_j)
   =n_i(n_kn_k)n_j=n_in_j,\\
   \big[\Proj{\bn}^{2}\big]_{ij}&=(\dd_{ik}-n_in_k)(\dd_{kj}-n_kn_j)\\
   &=\dd_{ik}\dd_{kj}-\dd_{ik}n_kn_j-n_in_k\dd_{kj}+n_in_k n_kn_j\\
   &=\dd_{ij}-n_in_j-n_in_j+n_in_j\\
   &=\dd_{ij}-n_in_j=(\Proj{\bn})_{ij},\\
   (\Proj{\bn})_{ji}&=\dd_{ji}-n_jn_i=\dd_{ij}-n_in_j=(\Proj{\bn})_{ij}.
\end{aligned}
\]
The remaining three are immediate from the definitions. For the zero tensor,
$\ip\bu{\Ob\bv}=\ip\bu\bze=0=\ip\bv\bze=\ip\bv{\Ob\bu}$ gives $\Ob^{t}=\Ob$,
and $\Ob^{2}\bv=\Ob(\Ob\bv)=\Ob\bze=\bze=\Ob\bv$ gives $\Ob^{2}=\Ob$; it
projects onto the trivial subspace $\{\bze\}$. For the identity,
$\ip\bu{\I\bv}=\ip\bu\bv=\ip\bv\bu=\ip\bv{\I\bu}$ gives $\I^{t}=\I$, and
$\I^{2}\bv=\I(\I\bv)=\I\bv$ gives $\I^{2}=\I$; it projects onto all of
$\Et$. For $\bn\otimes\bn$, symmetry and idempotency are exactly what was
proved above, and it projects onto the normal line, since
$(\bn\otimes\bn)\bv=\bv_n$ by $(\star_2)$.


\begin{figure}[htbp]\centering
\begin{tikzpicture}[scale=0.95]
  \fill[gray!12] (-2.2,-0.5) -- (2.6,-0.5) -- (3.35,0.45) -- (-1.45,0.45) -- cycle;
  \draw[gray!65,thin] (-2.2,-0.5) -- (2.6,-0.5) -- (3.35,0.45) -- (-1.45,0.45) -- cycle;
  \node[gray!75] at (-1.9,-0.28) {$P$};
  \coordinate (O) at (0,0); \fill (O) circle (1.3pt);
  \coordinate (Q) at (1.75,0.22);
  \coordinate (V) at (1.75,1.45);
  \draw[vec] (O) -- (V); \node[above left=-1pt] at (V) {$\bv$};
  \draw[vec] (O) -- (Q); \node[below=3pt] at (1.35,0.18) {$\Proj{\bn}\bv$};
  \draw[guide] (V) -- (Q);
  \draw[-{Stealth[length=2mm]},thick,gray!55] (Q) ++(0.12,0.30) arc (110:-110:0.34);
  \node[gray!70,right=12pt] at (2.05,0.22) {$\Proj{\bn}$ again};
\end{tikzpicture}
\\[1.2ex]
{\footnotesize\textsc{Figure P10.2.} $\Proj{\bn}^{2}=\Proj{\bn}$: the second
application meets a vector already in $P$, since
$\ip{(\Proj{\bn}\bv)}{\bn}=0$.}
\end{figure}


The two projections $\Proj{\bn}$ and $\bn\otimes\bn$ are not independent of
each other: they are \emph{complementary}, and this is the algebraic face of
$\Et=P\oplus\mathcal{N}$. Directly from $(\star_3)$ and $(\star_4)$,
\[
\begin{aligned}
   \Proj{\bn}&+\bn\otimes\bn=\I,\\
   \Proj{\bn}(\bn\otimes\bn)&=(\I-\bn\otimes\bn)(\bn\otimes\bn)
   =\bn\otimes\bn-(\bn\otimes\bn)^{2}\\
   &\overset{(\star_4)}{=}\Ob=(\bn\otimes\bn)\Proj{\bn},
\end{aligned}
\]
so each annihilates the range of the other. Their traces count the
dimensions of what they keep; using linearity of the trace with $\tr\I=3$
and $\tr(\bu\otimes\bv)=\ip\bu\bv$,
\[
   \tr\Proj{\bn}=\tr\I-\tr(\bn\otimes\bn)=3-\ip\bn\bn=3-1=2=\dim P,
\]
\[
   \tr(\bn\otimes\bn)=\ip\bn\bn=1=\dim\mathcal{N}.
\]
Finally $\Proj{\bn}$ is singular, since it kills the nonzero vector $\bn$:
\[
   \Proj{\bn}\bn=\bn-(\bn\otimes\bn)\bn=\bn-(\ip\bn\bn)\bn=\bn-\bn=\bze,
\]
\[
   \bn\neq\bze
   \ \Longrightarrow\ \det\Proj{\bn}=0 ,
\]
so no projection other than $\I$ is invertible.

\emph{(c)} The mirror image must leave the part lying in the plane alone and
send the normal part to its negative. In the notation of $(\star_2)$ this
says $\Refl{\bn}\bv=\bv_p-\bv_n$, and subtracting the normal part once to
reach the plane and once more to pass through it produces the factor $2$:
\begin{equation}
   \Refl{\bn}\bv=\bv_p-\bv_n
   =\big[\bv-(\ip\bv\bn)\bn\big]-(\ip\bv\bn)\bn
   =\bv-2(\ip\bv\bn)\bn
   =\big[\I-2\,\bn\otimes\bn\big]\bv,
\end{equation}
valid for every $\bv$, whence
\begin{equation}
   \Refl{\bn}=\I-2\,\bn\otimes\bn=\Proj{\bn}-\bn\otimes\bn,\qquad
   (\Refl{\bn})_{ij}=\dd_{ij}-2n_in_j . \tag{$\star_5$}
\end{equation}


\begin{figure}[htbp]\centering
\begin{tikzpicture}[scale=1.05]
  \fill[gray!12] (-2.6,-0.5) -- (2.6,-0.5) -- (3.35,0.45) -- (-1.85,0.45) -- cycle;
  \draw[gray!65,thin] (-2.6,-0.5) -- (2.6,-0.5) -- (3.35,0.45) -- (-1.85,0.45) -- cycle;
  \node[gray!75] at (-2.2,-0.28) {$P$};
  \coordinate (O) at (0,0); \fill (O) circle (1.4pt);
  \coordinate (Q) at (1.8,-0.02);
  \coordinate (V) at (1.8,1.2);
  \coordinate (W) at (1.8,-1.24);
  \draw[vec] (O) -- (0,0.62); \node[left=2pt] at (0,0.38) {$\bn$};
  \draw[vec] (O) -- (V); \node[above left=-1pt] at (V) {$\bv$};
  \draw[vec] (O) -- (W); \node[below left=-1pt] at (W) {$\Refl{\bn}\bv$};
  \draw[vec,black] (O) -- (Q); \node[above left=-1pt and 6pt] at (Q) {$\bv_p$};
  \draw[thinvec,black] (Q) -- (V); \node[right=3pt] at (1.83,0.62) {$\bv_n$};
  \draw[thinvec,black] (Q) -- (W); \node[right=3pt] at (1.83,-0.64) {$-\bv_n$};
  \draw[gray!70,thin] (Q) ++(-0.19,0.02) -- ++(0.02,0.19) -- ++(0.19,-0.02);
\end{tikzpicture}
\\[1.2ex]
{\footnotesize\textsc{Figure P10.3.} $\Refl{\bn}\bv=\bv_p-\bv_n$; the normal
part is removed twice, giving the factor $2$ in
$\Refl{\bn}=\I-2\,\bn\otimes\bn$.}
\end{figure}


Three properties can be established before any component is computed, and
they will serve as checks on the two matrices below. Symmetry is inherited
from $\I$ and $\bn\otimes\bn$, and the square collapses because the dyad is
idempotent:
\[
   \Refl{\bn}^{t}=\I^{t}-2(\bn\otimes\bn)^{t}=\I-2\,\bn\otimes\bn
   =\Refl{\bn},
\]
\[
   \Refl{\bn}^{2}=\I-2\,\bn\otimes\bn-2\,\bn\otimes\bn
   +4(\bn\otimes\bn)^{2}
   =\I-4\,\bn\otimes\bn+4(\bn\otimes\bn)^{2}
   \overset{(\star_4)}{=}\I,
\]
and the same in components,
\[
\begin{aligned}
   \big[\Refl{\bn}^{2}\big]_{ij}&=(\dd_{ik}-2n_in_k)(\dd_{kj}-2n_kn_j)\\
   &=\dd_{ij}-2n_in_j-2n_in_j+4n_i(n_kn_k)n_j=\dd_{ij},
\end{aligned}
\]
Thus $\Refl{\bn}$ is its own inverse, $\Refl{\bn}^{-1}=\Refl{\bn}$:
reflecting twice restores the original, as it must. Combining this with
symmetry gives $\Refl{\bn}^{t}\Refl{\bn}=\Refl{\bn}^{2}=\I$, so
$\Refl{\bn}$ is orthogonal and therefore preserves every inner product,
hence every length and angle:
\[
   \ip{\Refl{\bn}\bu}{\Refl{\bn}\bv}
   =\ip\bu{\Refl{\bn}^{t}\Refl{\bn}\bv}=\ip\bu{\I\bv}=\ip\bu\bv .
\]
For the determinant, write $\Refl{\bn}=\I+\ba\otimes\bb$ with $\ba=-2\bn$
and $\bb=\bn$; Problem~21(f) states $\det(\I+\ba\otimes\bb)=1+\ip\ba\bb$, so
\[
   \det\Refl{\bn}=1+\ip{(-2\bn)}{\bn}=1-2\,\ip\bn\bn=1-2=-1 ,
\]
the negative sign saying that a reflection reverses orientation, in contrast
to a rotation. The eigenvalues follow from the geometry: the normal is
reversed and the plane is fixed,
\[
\begin{aligned}
   \Refl{\bn}\bn&=\bn-2(\ip\bn\bn)\bn=\bn-2\bn=-\bn,\\
   \Refl{\bn}\bw&=\bw-2(\ip\bw\bn)\bn=\bw \quad (\bw\in P),
\end{aligned}
\]
so $\lambda=-1$ once and $\lambda=+1$ twice, since $\dim P=2$. This is
confirmed by the characteristic polynomial: with
\[
   I_1=\tr\Refl{\bn}=\tr\I-2\,\tr(\bn\otimes\bn)=3-2\,\ip\bn\bn=1,
\]
\[
   I_2=\tfrac12\big[I_1^{2}-\tr(\Refl{\bn}^{2})\big]
   =\tfrac12\big[1-\tr\I\big]=\tfrac12[1-3]=-1,
   \qquad I_3=\det\Refl{\bn}=-1,
\]
Problem~23 gives
\[
   \det(\Refl{\bn}-\lambda\I)=-\lambda^{3}+I_1\lambda^{2}-I_2\lambda+I_3
   =-\lambda^{3}+\lambda^{2}+\lambda-1
   =-(\lambda-1)^{2}(\lambda+1),
\]
whose roots are $+1,+1,-1$ as expected.

\emph{(i) $\bn=\be_1$.} Here $n_i=\ip{\be_1}{\be_i}=\dd_{i1}$, so
$n_1=1$ and $n_2=n_3=0$, and the outer product has entries
$n_in_j=\dd_{i1}\dd_{j1}$, which vanish unless $i=j=1$:
\[
   (n_in_j)=\begin{pmatrix}1&0&0\\0&0&0\\0&0&0\end{pmatrix}.
\]
Substituting into $(\star_5)$ and subtracting entry by entry,
\[
   (R_{ij})=\dd_{ij}-2n_in_j
   =\begin{pmatrix}1&0&0\\0&1&0\\0&0&1\end{pmatrix}
   -2\begin{pmatrix}1&0&0\\0&0&0\\0&0&0\end{pmatrix}
   =\begin{pmatrix}-1&0&0\\0&1&0\\0&0&1\end{pmatrix}.
\]
Reassembling the tensor from its components,
$\Refl{\be_1}=R_{ij}\,\be_i\otimes\be_j$, the only surviving terms being the
diagonal ones,
\[
   \Refl{\be_1}=-\be_1\otimes\be_1+\be_2\otimes\be_2+\be_3\otimes\be_3 .
\]
Its trace is $R_{11}+R_{22}+R_{33}=-1+1+1=1$ and, the matrix being
diagonal, its determinant is the product of the diagonal entries,
$(-1)(1)(1)=-1$, matching the general results.

\emph{(ii) $\bn=\tfrac{\sqrt2}{2}(\be_1+\be_2)$.} First confirm the
hypothesis $\ip\bn\bn=1$:
\[
   |\bn|^{2}=\Big(\tfrac{\sqrt2}{2}\Big)^{2}(1^{2}+1^{2}+0^{2})
   =\tfrac12\cdot2=1 .
\]
The components are $n_1=n_2=\tfrac{\sqrt2}{2}=\tfrac{1}{\sqrt2}$ and
$n_3=0$, so the nonzero products are
\[
   n_1n_1=n_1n_2=n_2n_1=n_2n_2=\tfrac{1}{\sqrt2}\cdot\tfrac{1}{\sqrt2}
   =\tfrac12 ,
\]
every product involving the index $3$ vanishing because $n_3=0$. Hence
\[
   (n_in_j)=\tfrac12\begin{pmatrix}1&1&0\\1&1&0\\0&0&0\end{pmatrix},
   \qquad 2(n_in_j)=\begin{pmatrix}1&1&0\\1&1&0\\0&0&0\end{pmatrix},
\]
and $(\star_5)$ gives, subtracting entry by entry,
\[
   (R_{ij})=\begin{pmatrix}1&0&0\\0&1&0\\0&0&1\end{pmatrix}
   -\begin{pmatrix}1&1&0\\1&1&0\\0&0&0\end{pmatrix}
   =\begin{pmatrix}0&-1&0\\-1&0&0\\0&0&1\end{pmatrix}.
\]
Reassembling $\Refl{\bn}=R_{ij}\be_i\otimes\be_j$ with the three nonzero
entries $R_{12}=R_{21}=-1$ and $R_{33}=1$,
\[
   \Refl{\bn}=-\be_1\otimes\be_2-\be_2\otimes\be_1+\be_3\otimes\be_3 .
\]


\begin{figure}[htbp]\centering
\begin{tikzpicture}[scale=1.45]
  \draw[guide] (-1.75,0) -- (1.55,0); \draw[guide] (0,-1.6) -- (0,1.35);
  \node[gray!70] at (1.68,0.05) {$x_1$}; \node[gray!70] at (0.15,1.3) {$x_2$};
  \draw[very thick] (-1.15,1.15) -- (1.15,-1.15);
  \node[gray!85,font=\footnotesize,align=center] at (-1.45,1.45)
     {$P:\ x_1+x_2=0$ (edge on),\\[-2pt]fixed pointwise};
  \draw[vec] (0,0) -- (0.707,0.707);
  \node[above right=-1pt] at (0.70,0.70) {$\bn$};
  \draw (0.34,0) arc (0:45:0.34); \node at (0.53,0.16) {\footnotesize$45^{\circ}$};
  \draw[vec] (0,0) -- (1,0); \node[above right=-1pt] at (1.0,0.02) {$\be_1$};
  \draw[vec] (0,0) -- (0,1); \node[above left=-1pt] at (0,1.0) {$\be_2$};
  \draw[vec,gray!62] (0,0) -- (0,-1);
  \draw[vec,gray!62] (0,0) -- (-1,0);
  \node[gray!72,font=\small] at (-0.78,-0.30) {$\Refl{\bn}\be_2=-\be_1$};
  \node[gray!72,font=\small] at (-0.30,-1.18) {$\Refl{\bn}\be_1=-\be_2$};
  \draw[guide,-{Stealth[length=1.5mm]}] (1.02,0.13) arc (7:-83:1.03);
  \draw[guide,-{Stealth[length=1.5mm]}] (0.13,1.02) arc (83:173:1.03);
\end{tikzpicture}
\\[1.2ex]
{\footnotesize\textsc{Figure P10.4.} Case (ii): the mirror $x_1+x_2=0$ seen
edge on, normal at $45^{\circ}$. The axes are swapped and negated, which is
why $(R_{ij})$ has zero diagonal and $-1$ off it.}
\end{figure}


It remains to verify that this matrix does what the general theory
predicted. Compute the action on each basis vector from
$\Refl{\bn}\be_j=R_{ij}\be_i$, that is by reading the $j$th column and
multiplying it against $\{\be_i\}$:
\[
\begin{aligned}
   \Refl{\bn}\be_1&=R_{11}\be_1+R_{21}\be_2+R_{31}\be_3
    =0\,\be_1+(-1)\be_2+0\,\be_3=-\be_2,\\
   \Refl{\bn}\be_2&=R_{12}\be_1+R_{22}\be_2+R_{32}\be_3
    =(-1)\be_1+0\,\be_2+0\,\be_3=-\be_1,\\
   \Refl{\bn}\be_3&=R_{13}\be_1+R_{23}\be_2+R_{33}\be_3
    =0\,\be_1+0\,\be_2+1\,\be_3=\be_3 .
\end{aligned}
\]
\emph{The normal is reversed.} Using linearity and the three images just
computed,
\[
   \Refl{\bn}\bn=\tfrac{\sqrt2}{2}\big(\Refl{\bn}\be_1+\Refl{\bn}\be_2\big)
   =\tfrac{\sqrt2}{2}(-\be_2-\be_1)
   =-\tfrac{\sqrt2}{2}(\be_1+\be_2)=-\bn .
\]
\emph{The plane is fixed pointwise.} A basis of $P$ is
$\{\be_1-\be_2,\ \be_3\}$; both really lie in $P$, since
\[
   \ip{(\be_1-\be_2)}{\bn}=\tfrac{\sqrt2}{2}(1-1)=0,\qquad
   \ip{\be_3}{\bn}=\tfrac{\sqrt2}{2}(0+0)=0,
\]
and both are fixed,
\[
   \Refl{\bn}(\be_1-\be_2)=\Refl{\bn}\be_1-\Refl{\bn}\be_2
   =-\be_2-(-\be_1)=\be_1-\be_2,\qquad
   \Refl{\bn}\be_3=\be_3 .
\]
So the eigenvalues really are $-1$ once and $+1$ twice.
\emph{The involution, checked on the matrix.} Multiplying $R$ by itself and
writing out all nine entries,
\[
\begin{aligned}
   R^{2}&=\begin{pmatrix}0&-1&0\\-1&0&0\\0&0&1\end{pmatrix}
         \begin{pmatrix}0&-1&0\\-1&0&0\\0&0&1\end{pmatrix}\\[1mm]
   &={\footnotesize\begin{pmatrix}
     (0)(0)+(-1)(-1)+0 & (0)(-1)+(-1)(0)+0 & 0\\
     (-1)(0)+(0)(-1)+0 & (-1)(-1)+(0)(0)+0 & 0\\
     0 & 0 & (1)(1)\end{pmatrix}}
   =\begin{pmatrix}1&0&0\\0&1&0\\0&0&1\end{pmatrix},
\end{aligned}
\]
so $R^{2}=I$, and since $R^{t}=R$ this also confirms $R^{t}R=I$ entry by
entry. \emph{Trace and determinant.} Summing the diagonal,
$\tr R=R_{11}+R_{22}+R_{33}=0+0+1=1$; expanding the determinant along the
third row, whose only nonzero entry is $R_{33}=1$,
\[
   \det R=1\cdot\det\begin{pmatrix}0&-1\\-1&0\end{pmatrix}
   =1\cdot\big[(0)(0)-(-1)(-1)\big]=-1 .
\]
Both agree with $I_1=1$ and $I_3=-1$ found from $(\star_5)$, and with the
spectrum $\{-1,+1,+1\}$.
\end{solution}


\begin{problem}[Triple products; two induced tensors]
(a) Prove $\ba\times(\bb\times\bc)=(\ip\ba\bc)\bb-(\ip\ba\bb)\bc$ and
$\bv=(\ip\bv\be)\be+\be\times(\bv\times\be)$ ($\be$ unit); (b) tensor for
$f(\bv)=\be\times(\bv\times\be)$; (c) tensor for $f(\bv)=\bv+\ba\times\bv$.
\end{problem}
\begin{solution}
\textbf{(a)} Write the inner cross product first, then the outer one:
\[
   (\bb\times\bc)_k=\eps_{klm}b_lc_m,
   \qquad
   [\ba\times(\bb\times\bc)]_i=\eps_{ijk}a_j\,(\bb\times\bc)_k
   =\eps_{ijk}\eps_{klm}\,a_jb_lc_m .
\]
Contract the two symbols. Theorem~\ref{thm:epsdelta} contracts on the last
index, and $\eps_{ijk}=\eps_{kij}$ by cyclic shift, so
\[
   \eps_{ijk}\eps_{klm}=\eps_{kij}\eps_{klm}
   =\dd_{il}\dd_{jm}-\dd_{im}\dd_{jl}.
\]
Substituting and letting the deltas collapse the sums,
\[
   [\ba\times(\bb\times\bc)]_i
   =(\dd_{il}\dd_{jm}-\dd_{im}\dd_{jl})a_jb_lc_m
   =a_jc_j\,b_i-a_jb_j\,c_i
   =(\ip\ba\bc)b_i-(\ip\ba\bb)c_i .
\]
Now put $\ba=\bc=\be$ with $|\be|=1$, so $\ip\be\be=1$:
\[
   \be\times(\bv\times\be)=(\ip\be\be)\bv-(\ip\be\bv)\be
   =\bv-(\ip\bv\be)\be ,
\]
which rearranges to $\bv=(\ip\bv\be)\be+\be\times(\bv\times\be)$.

\textbf{(b)} From (a), $f(\bv)=\bv-(\ip\be\bv)\be$. Both terms are linear
in $\bv$, and $(\ip\be\bv)\be=(\be\otimes\be)\bv$, so
\[
   f(\bv)=\big(\I-\be\otimes\be\big)\bv,
   \qquad \bA=\I-\be\otimes\be=\Proj\be ,
\]
the projection of Problem 10 onto the plane perpendicular to $\be$. Check:
$f(\be)=\be-\be=\bze$ and $f(\bw)=\bw$ for $\bw\perp\be$.

\textbf{(c)} $f(\bv)=\bv+\ba\times\bv$. In components,
$(\ba\times\bv)_i=\eps_{ijk}a_jv_k$, so the second term is
$\bOm\bv$ with $\Omega_{ik}=\eps_{ijk}a_j=-\eps_{ikj}a_j$, skew by
Problem 16. Hence
\[
   \bA=\I+\bOm,\qquad A_{ij}=\dd_{ij}-\eps_{ijk}a_k .
\]
Linearity of $f$ is then visible, and $\det\bA=1+|\ba|^{2}\neq0$ by
Problem 21(g), so $\bA$ is invertible for every $\ba$.
\end{solution}

\begin{problem}[Two-point tensors]
For orthonormal bases $\{\be_i\}$ of $\Et$ and $\{\hbE_A\}$ of $\hEt$,
show a linear $\bA:\hEt\to\Et$ is $\bA=A_{iB}\be_i\otimes\hbE_B$ with
$A_{iB}=\ip{\be_i}{\bA\hbE_B}$.
\end{problem}
\begin{solution}
\emph{Step 1.} $\bA\hbE_B\in\Et$, so resolve it on the spatial basis and
name the coefficients:
\[
   \bA\hbE_B=\big(\ip{\be_i}{\bA\hbE_B}\big)\be_i=:A_{iB}\,\be_i .
\]
The dot product is legal, both arguments being in $\Et$.

\emph{Step 2.} For $\bV=V_B\hbE_B\in\hEt$, by linearity,
\[
   \bA\bV=V_B\,\bA\hbE_B=A_{iB}V_B\,\be_i .
\]

\emph{Step 3.} The mixed dyad acts by $(\ba\otimes\bb)\bc=(\ip\bb\bc)\ba$:
\[
   (\be_i\otimes\hbE_B)\bV=\big(\ip{\hbE_B}{\bV}\big)\be_i=V_B\,\be_i ,
\]
legal since $\hbE_B,\bV\in\hEt$ while $\be_i\in\Et$. Hence
$\big(A_{iB}\be_i\otimes\hbE_B\big)\bV=A_{iB}V_B\be_i=\bA\bV$ for all
$\bV$, and
\[
   \boxed{\ \bA=A_{iB}\,\be_i\otimes\hbE_B,\qquad
   A_{iB}=\ip{\be_i}{\bA\hbE_B}.\ }
\]

\emph{Step 4, consistency.} Feed $\hbE_C$ to the dyadic form, then dot
with $\be_j$:
\[
   \big(A_{iB}\be_i\otimes\hbE_B\big)\hbE_C=A_{iB}\dd_{BC}\be_i
   =A_{iC}\be_i,
   \qquad
   \ip{\be_j}{\bA\hbE_C}=A_{iC}\dd_{ji}=A_{jC}.
\]
\end{solution}

\begin{remark}[which dots are legal]\label{rem:twopointlegal}
\[
\begin{aligned}
   \ip{\be_i}{\bA\hbE_B}&\quad\text{defined, both in }\Et,\\
   \ip{\hbE_B}{\bV}&\quad\text{defined, both in }\hEt,\\
   \ip{\be_i}{\hbE_B}&\quad\text{\emph{not} defined},\\
   \be_i\otimes\hbE_B&\quad\text{defined, a map }\hEt\to\Et .
\end{aligned}
\]
Also $\bA^{t}=A_{iB}\hbE_B\otimes\be_i$ maps $\Et\to\hEt$, so
$\bA+\bA^{t}$ is meaningless: a two-point tensor is neither symmetric nor
skew. The object built here is $\bF=F_{iA}\be_i\otimes\hbE_A$ of \eqr{2.39}.
\end{remark}

\begin{problem}[Riesz representation for tensors]
Show a linear $g:\Lin\to\Rn$ is $g(\bV)=\bA\cdot\bV$ for a unique $\bA$.
\end{problem}
\begin{solution}
\emph{Existence.} Define the nine numbers
\[
   A_{ij}:=g(\be_i\otimes\be_j),\qquad
   \bA:=A_{ij}\,\be_i\otimes\be_j .
\]
Resolve an arbitrary $\bV$ on the same basis,
$\bV=V_{ij}\be_i\otimes\be_j$ with $V_{ij}=\ip{\be_i}{\bV\be_j}$. Since
$g$ is linear and the $V_{ij}$ are scalars,
\[
   g(\bV)=g\big(V_{ij}\be_i\otimes\be_j\big)
        =V_{ij}\,g(\be_i\otimes\be_j)
        =V_{ij}A_{ij}=\bA\cdot\bV ,
\]
the last step being the definition $\bA\cdot\bV=A_{ij}V_{ij}$ of the
tensor inner product.

\emph{Uniqueness.} If $\bA\cdot\bV=\bA'\cdot\bV$ for every $\bV$, then
$(\bA-\bA')\cdot\bV=0$ for every $\bV$. Choose $\bV=\bA-\bA'$:
\[
   0=(\bA-\bA')\cdot(\bA-\bA')=|\bA-\bA'|^{2}
   \ \Longrightarrow\ \bA=\bA' .
\]
This is the tensor version of Problem 13's vector counterpart, and the
proof is the same one twice: linearity gives existence through a basis,
positive definiteness of the inner product gives uniqueness.
\end{solution}

\begin{problem}[Symmetric--skew decomposition]
$\bA=\bS+\bW$, $\bS=\tfrac12(\bA+\bA^t)$, $\bW=\tfrac12(\bA-\bA^t)$. (a)
uniqueness; (b) $\Sym,\Skw$ are orthogonal subspaces,
$\mathbb{E}^9=\Sym\oplus\Skw$.
\end{problem}
\begin{solution}
\emph{Existence.} With $\bS=\tfrac12(\bA+\bA^{t})$ and
$\bW=\tfrac12(\bA-\bA^{t})$,
\[
   \bS+\bW=\tfrac12(\bA+\bA^{t})+\tfrac12(\bA-\bA^{t})
   =\tfrac12\big(\bA+\bA^{t}+\bA-\bA^{t}\big)=\tfrac12(2\bA)=\bA ,
\]
and, by $(\bA+\bB)^{t}=\bA^{t}+\bB^{t}$ and $(\bA^{t})^{t}=\bA$,
\[
\begin{aligned}
   \bS^{t}&=\tfrac12\big(\bA^{t}+(\bA^{t})^{t}\big)
         =\tfrac12(\bA^{t}+\bA)=\bS,\\
   \bW^{t}&=\tfrac12\big(\bA^{t}-\bA\big)
         =-\tfrac12(\bA-\bA^{t})=-\bW .
\end{aligned}
\]

\textbf{(a) Uniqueness.} Let $\bA=\bS'+\bW'$ with $(\bS')^{t}=\bS'$,
$(\bW')^{t}=-\bW'$. Transpose:
\[
   \bA^{t}=(\bS'+\bW')^{t}=(\bS')^{t}+(\bW')^{t}=\bS'-\bW' .
\]
Now add and subtract this from $\bA=\bS'+\bW'$:
\[
\begin{aligned}
   \bA+\bA^{t}&=2\bS'\ \Longrightarrow\ \bS'=\tfrac12(\bA+\bA^{t})=\bS,\\
   \bA-\bA^{t}&=2\bW'\ \Longrightarrow\ \bW'=\tfrac12(\bA-\bA^{t})=\bW .
\end{aligned}
\]
\textbf{(b) Orthogonal subspaces.} Each is a subspace:
$(\alpha\bS_{1}+\beta\bS_{2})^{t}=\alpha\bS_{1}+\beta\bS_{2}$, and the
same with a minus sign for $\Skw$. For orthogonality, with
$\bS\cdot\bW=S_{ij}W_{ij}$,
\[
   \bS\cdot\bW=S_{ij}W_{ij}
   \ \overset{(1)}{=}\ S_{ji}W_{ij}
   \ \overset{(2)}{=}\ S_{ij}W_{ji}
   \ \overset{(3)}{=}\ -S_{ij}W_{ij}=-\bS\cdot\bW ,
\]
with (1) $S_{ij}=S_{ji}$, (2) a relabelling, (3) $W_{ji}=-W_{ij}$. Hence
$\bS\cdot\bW=0$.

\emph{Dimensions.} $\bS$ is fixed by
$S_{11},S_{22},S_{33},S_{12},S_{13},S_{23}$, so $\dim\Sym=6$. For $\bW$,
$W_{ii}=-W_{ii}$ gives a zero diagonal, leaving
$W_{12},W_{13},W_{23}$, so $\dim\Skw=3$. Their sum is $9$, and they meet
only in $\Ob$, since $\bA=\bA^{t}=-\bA$ forces $2\bA=\Ob$. Therefore
\[
   \mathbb{E}^{9}=\Sym\oplus\Skw .
\]
\end{solution}

\begin{remark}[three uses]\label{rem:symskwuse}
(i) $\ip{\ba}{\bA\ba}=\ip{\ba}{\bS\ba}$ by Problem 6(a), so positive
definiteness \eqr{3.13} sees only $\bS$.
(ii) Chapter 4 splits $\bL=\grad\bv=\bD+\bW$ into stretching and spin;
part (a) is what makes them unique.
(iii) By Problem 16 every skew $\bW$ is $\bW\bu=\bom\times\bu$, which is
why $\dim\Skw=3$.

Expanding $\bA\cdot\bA=(\bS+\bW)\cdot(\bS+\bW)$ and killing the cross
terms by (b) gives $|\bA|^{2}=|\bS|^{2}+|\bW|^{2}$: the splitting is an
orthogonal projection, not merely a decomposition.
\end{remark}

\begin{problem}[Spherical--deviatoric decomposition]
$\bT=\bA+\alpha\I$, $\tr\bA=0$, $3\alpha=\tr\bT$. (a) uniqueness; (b)
$\Sph,\Dev$ orthogonal, $\mathbb{E}^9=\Sph\oplus\Dev$.
\end{problem}
\begin{solution}
\emph{Existence.} Put $\alpha=\tfrac13\tr\bT$ and $\bA=\bT-\alpha\I$.
Then $\bT=\bA+\alpha\I$ and, using $\tr\I=\dd_{ii}=3$,
\[
   \tr\bA=\tr\bT-\alpha\,\tr\I=\tr\bT-\tfrac13(\tr\bT)(3)=0 .
\]

\textbf{(a) Uniqueness.} Let $\bT=\bB+\beta\I$ with $\tr\bB=0$. Take the
trace:
\[
   \tr\bT=\tr\bB+\beta\,\tr\I=0+3\beta
   \ \Longrightarrow\ \beta=\tfrac13\tr\bT=\alpha ,
\]
and then $\bB=\bT-\beta\I=\bT-\alpha\I=\bA$.

\textbf{(b) Orthogonal subspaces.} Both are subspaces: $\Sph$ is the span
of $\I$, and $\tr$ is linear so $\Dev=\{\bD:\tr\bD=0\}$ is its kernel.
For $\lambda\I\in\Sph$ and $\bD\in\Dev$, using
$\I\cdot\bD=\dd_{ij}D_{ij}=D_{ii}$,
\[
   (\lambda\I)\cdot\bD=\lambda\,\tr\bD=\lambda\cdot0=0 .
\]
Dimensions: $\dim\Sph=1$; $\tr:\mathbb{E}^{9}\to\Rn$ is onto, so its
kernel has $\dim\Dev=9-1=8$. Since $1+8=9$ and
$\Sph\cap\Dev=\{\Ob\}$, because $\lambda\I\in\Dev$ forces
$3\lambda=0$,
\[
   \mathbb{E}^{9}=\Sph\oplus\Dev .
\]
\end{solution}

\begin{remark}[the two splittings are independent]\label{rem:twosplits}
Problems 14 and 15 split $\mathbb{E}^{9}$ two different ways,
$\Sym\oplus\Skw$ and $\Sph\oplus\Dev$, and the two are compatible rather
than competing: $\Sph\subset\Sym$, so a single tensor decomposes as
\[
   \bT=\underbrace{\alpha\I}_{1}
      +\underbrace{(\Sym\bT-\alpha\I)}_{5}
      +\underbrace{\Skw\bT}_{3},
   \qquad \alpha=\tfrac13\tr\bT,
\]
with dimensions $1+5+3=9$. In Chapter 3 the middle piece is what carries
distortion at constant volume and the first is what carries volume change:
there $\tr(\ln\vek{U})=\ln J$ exactly, so the spherical part of the Hencky
strain is $\tfrac13(\ln J)\I$.
\end{remark}

\begin{problem}[Axial vector of a skew tensor]
$f(\bv)=\bom\times\bv=\bOm\bv$. (a) $\Omega_{ij}=\eps_{jik}\omega_k$, so
$\bOm\in\Skw$; (b) $\omega_k=\tfrac12\Omega_{ij}\eps_{jik}$.
\end{problem}
\begin{solution}
\textbf{(a)} Write both sides of $\bOm\bv=\bom\times\bv$ in components:
\[
   (\bOm\bv)_i=\Omega_{ij}v_j,
   \qquad
   (\bom\times\bv)_i=\eps_{ikj}\omega_kv_j ,
\]
the second by Proposition~\ref{prop:crosscomp} with the summation indices
named so that $v$ carries $j$ in both. Equating for every $\bv$,
\[
   \Omega_{ij}=\eps_{ikj}\omega_k .
\]
Now $\eps_{ikj}=-\eps_{ijk}=\eps_{jik}$: one swap gives a minus, a second
swap gives it back. So $\Omega_{ij}=\eps_{jik}\omega_k$, and
interchanging $i,j$,
\[
   \Omega_{ji}=\eps_{ijk}\omega_k=-\eps_{jik}\omega_k=-\Omega_{ij},
\]
i.e.\ $\bOm\in\Skw$.

\textbf{(b)} Substitute $\Omega_{ij}=\eps_{jil}\omega_l$ and contract:
\[
   \tfrac12\,\Omega_{ij}\eps_{jik}
   =\tfrac12\,\eps_{jil}\eps_{jik}\,\omega_l
   =\tfrac12\,(2\dd_{lk})\,\omega_l=\omega_k ,
\]
the middle step by Theorem~\ref{thm:epsdelta}: $\eps_{jil}\eps_{jik}
=\eps_{ilj}\eps_{ikj}=2\dd_{lk}$ after a cyclic shift brings the repeated
pair to the front.

So $\bom\mapsto\bOm$ and $\bOm\mapsto\bom$ invert one another: $\Skw$ and
$\Et$ are in one-to-one correspondence, which is why
$\dim\Skw=3$ in Problem 14.
\end{solution}

\begin{problem}[Identities for the tensor product]
Verify (a) $(\ba\otimes\bb)^t=\bb\otimes\ba$; (b)
$\bA(\ba\otimes\bb)=(\bA\ba)\otimes\bb$; (c)
$(\ba\otimes\bb)\bA=\ba\otimes(\bA^{t}\bb)$; (d)
$\bA=(\bA\be_i)\otimes\be_i=\be_i\otimes(\bA^{t}\be_i)$.
\end{problem}
\begin{solution}
Throughout, the definition is $(\ba\otimes\bb)\bc=(\ip\bb\bc)\ba$: the
second slot dots, the first survives.

\textbf{(a)} Test against arbitrary $\bu,\bv$ and use the symmetry of the
dot product:
\[
   \ip{\bu}{(\ba\otimes\bb)\bv}=(\ip\bb\bv)(\ip\ba\bu)
   =(\ip\ba\bu)(\ip\bb\bv)
   =\ip{\bv}{(\bb\otimes\ba)\bu} .
\]
By the definition of the transpose this says
$(\ba\otimes\bb)^{t}=\bb\otimes\ba$.

\textbf{(b)} Apply both sides to arbitrary $\bv$ and use linearity of
$\bA$:
\[
   \big[\bA(\ba\otimes\bb)\big]\bv=\bA\big[(\ip\bb\bv)\ba\big]
   =(\ip\bb\bv)\,\bA\ba
   =\big[(\bA\ba)\otimes\bb\big]\bv .
\]

\textbf{(c)} Same test, now moving $\bA$ across the dot product:
\[
   \big[(\ba\otimes\bb)\bA\big]\bv=(\ip{\bb}{\bA\bv})\,\ba
   =(\ip{\bA^{t}\bb}{\bv})\,\ba
   =\big[\ba\otimes(\bA^{t}\bb)\big]\bv .
\]
So post-multiplication touches only the second slot, pre-multiplication
only the first.

\textbf{(d)} With $\bA=A_{ij}\be_i\otimes\be_j$ one has
$\bA\be_i=A_{ki}\be_k$ and $\bA^{t}\be_i=A_{ik}\be_k$. Hence, summing on
$i$,
\[
   (\bA\be_i)\otimes\be_i=A_{ki}\,\be_k\otimes\be_i=\bA,
   \qquad
   \be_i\otimes(\bA^{t}\be_i)=A_{ik}\,\be_i\otimes\be_k=\bA ,
\]
the second after renaming the two dummy indices. Both are \eqr{3.11} with
$\bA$ inserted on one side or the other.
\end{solution}

\begin{problem}[Transpose of the inverse]
Show $(\bA^{t})^{-1}=(\bA^{-1})^{t}=:\bA^{-t}$.
\end{problem}
\begin{solution}
Transpose the two defining relations of $\bA^{-1}$, using
$(\bB\bC)^{t}=\bC^{t}\bB^{t}$ from Problem 20 and $\I^{t}=\I$:
\[
   \bA\bA^{-1}=\I\ \Longrightarrow\ (\bA^{-1})^{t}\bA^{t}=\I,
   \qquad
   \bA^{-1}\bA=\I\ \Longrightarrow\ \bA^{t}(\bA^{-1})^{t}=\I .
\]
So $(\bA^{-1})^{t}$ is a two-sided inverse of $\bA^{t}$, and inverses are
unique: if $\bP\bA^{t}=\I=\bA^{t}\bQ$ then
$\bP=\bP(\bA^{t}\bQ)=(\bP\bA^{t})\bQ=\bQ$. Hence
$(\bA^{t})^{-1}=(\bA^{-1})^{t}$, and the common value may be written
$\bA^{-t}$ without ambiguity.
\end{solution}

\begin{problem}[Inverse of a product]
$(\bB\bA)^{-1}=\bA^{-1}\bB^{-1}$.
\end{problem}
\begin{solution}
Multiply both ways, cancelling the inner pair first:
\[
   (\bA^{-1}\bB^{-1})(\bB\bA)=\bA^{-1}(\bB^{-1}\bB)\bA=\bA^{-1}\I\bA
   =\bA^{-1}\bA=\I,
\]
\[
   (\bB\bA)(\bA^{-1}\bB^{-1})=\bB(\bA\bA^{-1})\bB^{-1}=\bB\I\bB^{-1}
   =\bB\bB^{-1}=\I .
\]
By uniqueness of inverses, $(\bB\bA)^{-1}=\bA^{-1}\bB^{-1}$. The order
reverses, as it must: the last map applied is the first to be undone.
\end{solution}

\begin{problem}[Transpose of a product]
$(\bB\bA)^{t}=\bA^{t}\bB^{t}$.
\end{problem}
\begin{solution}
Move the two factors across the dot product one at a time, each by the
definition $\ip{\bu}{\bC\bv}=\ip{\bC^{t}\bu}{\bv}$:
\[
   \ip{\bu}{(\bB\bA)\bv}=\ip{\bu}{\bB(\bA\bv)}
   =\ip{\bB^{t}\bu}{\bA\bv}
   =\ip{\bA^{t}\bB^{t}\bu}{\bv}
   =\ip{\bu}{\big(\bA^{t}\bB^{t}\big)^{t}\bv}.
\]
Comparing the first and last for all $\bu,\bv$, and using that a tensor is
determined by its bilinear form, $(\bB\bA)^{t}=\bA^{t}\bB^{t}$. Again the
order reverses.
\end{solution}

\begin{problem}[Determinant identities via box products]
Verify (a) $\det(\alpha\bA)=\alpha^3\det\bA$; (b)
$\det(\bA\bB)=\det\bA\det\bB$; (c) $\det\I=1$; (d)
$\det(\bA^{-1})=(\det\bA)^{-1}$; (e) $\det(\ba\otimes\bb)=0$; (f)
$\det(\I+\ba\otimes\bb)=1+\ip\ba\bb$; (g)
$\det(\bA+\ba\otimes\bb)=\det\bA\,(1+\ip\ba{\bA^{-t}\bb})$.
\end{problem}
\begin{solution}
The master identity is $\det\bA\,\tbox\ba\bb\bc=\tbox{\bA\ba}{\bA\bb}{\bA\bc}$
for all $\ba,\bb,\bc$ (Definition~\ref{def:inv3}$_3$). Every part below
follows from it.

(a)
\[
\begin{aligned}
\det(\alpha\bA)\,\tbox\ba\bb\bc
&=\tbox{(\alpha\bA)\ba}{(\alpha\bA)\bb}{(\alpha\bA)\bc}\\
&=\alpha^3\tbox{\bA\ba}{\bA\bb}{\bA\bc}
=\alpha^3\det\bA\,\tbox\ba\bb\bc;
\end{aligned}
\]
cancel $\tbox\ba\bb\bc\neq0$.

(b)
\[
\begin{aligned}
\tbox{(\bA\bB)\ba}{(\bA\bB)\bb}{(\bA\bB)\bc}
&=\tbox{\bA(\bB\ba)}{\bA(\bB\bb)}{\bA(\bB\bc)}\\
&=\det\bA\,\tbox{\bB\ba}{\bB\bb}{\bB\bc}\\
&=\det\bA\,\det\bB\,\tbox\ba\bb\bc.
\end{aligned}
\]

(c) $\det\I\,\tbox\ba\bb\bc=\tbox{\I\ba}{\I\bb}{\I\bc}
=\tbox\ba\bb\bc$, so $\det\I=1$.

(d) From (b) and (c), $\det\bA\det(\bA^{-1})=\det(\bA\bA^{-1})=\det\I=1$.

(e) $(\ba\otimes\bb)\bu=(\ip\bb\bu)\ba$: each image is a scalar multiple
of $\ba$. Hence $\tbox{(\ba\otimes\bb)\be_1}{(\ba\otimes\bb)\be_2}
{(\ba\otimes\bb)\be_3}=\tbox{b_1\ba}{b_2\ba}{b_3\ba}$. A box product
with two parallel entries vanishes (alternation), so $\det(\ba\otimes\bb)=0$.

(f) Set $\bT=\I+\ba\otimes\bb$. Then $\bT\be_i=\be_i+b_i\ba$, and
\[
\det\bT=\tbox{(\be_1+b_1\ba)}{(\be_2+b_2\ba)}{(\be_3+b_3\ba)}.
\]
Expand multilinearly: the box product is trilinear, so there are $2^3=8$
terms, classified by how many of the three factors contribute $b_i\ba$:
\[
\begin{aligned}
&\underbrace{\tbox{\be_1}{\be_2}{\be_3}}_{=1}
+b_1\tbox\ba{\be_2}{\be_3}
+b_2\tbox{\be_1}\ba{\be_3}
+b_3\tbox{\be_1}{\be_2}\ba\\
&\quad+b_1b_2\underbrace{\tbox\ba\ba{\be_3}}_{=0}
+b_1b_3\underbrace{\tbox\ba{\be_2}\ba}_{=0}
+b_2b_3\underbrace{\tbox{\be_1}\ba\ba}_{=0}
+b_1b_2b_3\underbrace{\tbox\ba\ba\ba}_{=0}.
\end{aligned}
\]
Every term with $\ba$ appearing twice or thrice vanishes by alternation.
The surviving terms are: $1+b_1a_1+b_2a_2+b_3a_3=1+b_ia_i=1+\ip\ba\bb$.

(g) For invertible $\bA$, factor:
$\bA+\ba\otimes\bb=\bA(\I+\bA^{-1}(\ba\otimes\bb))
=\bA(\I+(\bA^{-1}\ba)\otimes\bb)$,
where Lemma~\ref{lem:Aab} gives
$\bA^{-1}(\ba\otimes\bb)=(\bA^{-1}\ba)\otimes\bb$. Apply (b), then (f):
\[
\det(\bA+\ba\otimes\bb)=\det\bA\,\det(\I+(\bA^{-1}\ba)\otimes\bb)
=\det\bA\,(1+\ip{\bA^{-1}\ba}\bb).
\]
Finally, $\ip{\bA^{-1}\ba}\bb=\ip\ba{(\bA^{-1})^t\bb}=\ip\ba{\bA^{-t}\bb}$,
using the definition of the transpose.
\end{solution}

\begin{problem}[Cofactor identities]
Let $\bA^{*}$ satisfy $\bA^{*}(\ba\times\bb)=\bA\ba\times\bA\bb$. Show:
(a) $I_2=\tr\bA^{*}$;\ (b) $I_2=\tfrac12[(\tr\bA)^2-\tr(\bA^2)]$;\
(c) $\bA^{t}\bA^{*}=(\det\bA)\I$, hence $\bA^{*}=(\det\bA)\bA^{-t}$;\
(d) recover $\det\bA=\tfrac16\eps_{ijk}\eps_{lmn}A_{il}A_{jm}A_{kn}$.
\end{problem}
\begin{solution}
\emph{Step 0: is $\bA^{*}$ even well defined?}
Definition~\ref{def:cof} prescribes the value of $\bA^{*}$ only on vectors
that are presented to us in the form $\ba\times\bb$. Before the symbol may
be used at all, two questions must be answered. Does the prescription reach
every vector, so that at most one tensor can satisfy it? And does a tensor
satisfying it exist, given that the same vector can be written as a cross
product in infinitely many ways, so that the prescription might contradict
itself?

\emph{Uniqueness.} Every vector is a cross product. If $\bv=\bze$, then
$\bv=\ba\times\ba$ for any $\ba$, since $\ba\times\ba=\bze$. If
$\bv\neq\bze$, set $\bu_{3}=\bv/|\bv|$, a unit vector, and complete it to a
right-handed orthonormal triad $\{\bu_{1},\bu_{2},\bu_{3}\}$, so that
$\bu_{1}\times\bu_{2}=\bu_{3}$. Then
\[
   \bu_{1}\times(|\bv|\,\bu_{2})=|\bv|\,(\bu_{1}\times\bu_{2})
   =|\bv|\,\bu_{3}=\bv .
\]
So the defining relation does prescribe $\bA^{*}\bv$ for every $\bv$, and
two tensors obeying it take the same value on every vector, hence are the
same tensor.

\emph{Existence.} We produce $\bA^{*}$ explicitly and verify the definition.
Put
\begin{equation}
   A^{*}_{mi}:=\tfrac12\,\eps_{mjk}\eps_{ipq}A_{jp}A_{kq},
   \qquad \bA^{*}:=A^{*}_{mi}\,\be_{m}\otimes\be_{i}. \tag{$\ast_1$}
\end{equation}
This is a tensor by construction, being a linear combination of dyads. Now
compute its action on $\ba\times\bb$, whose components are
$(\ba\times\bb)_{i}=\eps_{irs}a_{r}b_{s}$ by
Proposition~\ref{prop:crosscomp}:
\[
\begin{aligned}
   \big[\bA^{*}(\ba\times\bb)\big]_{m}
   &=A^{*}_{mi}\,(\ba\times\bb)_{i}\\
   &=\tfrac12\,\eps_{mjk}\eps_{ipq}A_{jp}A_{kq}\,\eps_{irs}a_{r}b_{s}\\
   &=\tfrac12\,\eps_{mjk}A_{jp}A_{kq}\,
     \big(\eps_{ipq}\eps_{irs}\big)\,a_{r}b_{s}.
\end{aligned}
\]
The bracket is contracted on its \emph{first} index, whereas
Theorem~\ref{thm:epsdelta} contracts on the last; the two are the same
statement because $\eps$ is unchanged by a cyclic shift of its indices,
$\eps_{ipq}=\eps_{pqi}$ and $\eps_{irs}=\eps_{rsi}$, so
\[
   \eps_{ipq}\eps_{irs}=\eps_{pqi}\eps_{rsi}
   =\dd_{pr}\dd_{qs}-\dd_{ps}\dd_{qr}.
\]
Substituting,
\[
\begin{aligned}
   \big[\bA^{*}(\ba\times\bb)\big]_{m}
   &=\tfrac12\,\eps_{mjk}A_{jp}A_{kq}
     \big(\dd_{pr}\dd_{qs}-\dd_{ps}\dd_{qr}\big)a_{r}b_{s}\\
   &=\tfrac12\,\eps_{mjk}A_{jp}A_{kq}a_{p}b_{q}
    -\tfrac12\,\eps_{mjk}A_{jp}A_{kq}a_{q}b_{p},
\end{aligned}
\]
the deltas having collapsed the sums by forcing $r=p,\ s=q$ in the first
term and $s=p,\ r=q$ in the second. The two terms are in fact equal. In the
second, first interchange the \emph{names} of the summation indices $p$ and
$q$, which changes nothing since both are summed over $1,2,3$:
\[
   -\tfrac12\,\eps_{mjk}A_{jp}A_{kq}a_{q}b_{p}
   =-\tfrac12\,\eps_{mjk}A_{jq}A_{kp}a_{p}b_{q};
\]
then interchange the names of the summation indices $j$ and $k$:
\[
   -\tfrac12\,\eps_{mjk}A_{jq}A_{kp}a_{p}b_{q}
   =-\tfrac12\,\eps_{mkj}A_{kq}A_{jp}a_{p}b_{q}
   =+\tfrac12\,\eps_{mjk}A_{jp}A_{kq}a_{p}b_{q},
\]
the last equality using $\eps_{mkj}=-\eps_{mjk}$. Hence the two halves add
rather than cancel, and
\[
   \big[\bA^{*}(\ba\times\bb)\big]_{m}
   =\eps_{mjk}\,(A_{jp}a_{p})(A_{kq}b_{q})
   =\eps_{mjk}\,(\bA\ba)_{j}(\bA\bb)_{k}
   =(\bA\ba\times\bA\bb)_{m}.
\]
So $(\ast_1)$ does satisfy Definition~\ref{def:cof}, and by the uniqueness
just proved it is \emph{the} cofactor. Both the object and its components
are now in hand.

\medskip
\emph{(a) $I_{2}(\bA)=I_{1}(\bA^{*})$.}
By Theorem~\ref{thm:trace}, $I_{1}(\bA^{*})=\tr\bA^{*}=A^{*}_{11}+A^{*}_{22}
+A^{*}_{33}$ with $A^{*}_{ii}=\ip{\be_{i}}{\bA^{*}\be_{i}}$ (no sum). For a
right-handed orthonormal basis,
$\be_{1}=\be_{2}\times\be_{3}$, $\be_{2}=\be_{3}\times\be_{1}$,
$\be_{3}=\be_{1}\times\be_{2}$, so the definition of $\bA^{*}$ evaluates on
each basis vector separately. The three computations, written out in full
rather than left to a cyclic instruction:
\[
\begin{aligned}
   \bA^{*}\be_{1}&=\bA^{*}(\be_{2}\times\be_{3})=\bA\be_{2}\times\bA\be_{3},
   &A^{*}_{11}&=\ip{\be_{1}}{\bA\be_{2}\times\bA\be_{3}}
               =\tbox{\be_{1}}{\bA\be_{2}}{\bA\be_{3}},\\
   \bA^{*}\be_{2}&=\bA^{*}(\be_{3}\times\be_{1})=\bA\be_{3}\times\bA\be_{1},
   &A^{*}_{22}&=\ip{\be_{2}}{\bA\be_{3}\times\bA\be_{1}}
               =\tbox{\be_{2}}{\bA\be_{3}}{\bA\be_{1}},\\
   \bA^{*}\be_{3}&=\bA^{*}(\be_{1}\times\be_{2})=\bA\be_{1}\times\bA\be_{2},
   &A^{*}_{33}&=\ip{\be_{3}}{\bA\be_{1}\times\bA\be_{2}}
               =\tbox{\be_{3}}{\bA\be_{1}}{\bA\be_{2}} .
\end{aligned}
\]
The box product is unchanged by a cyclic permutation of its three slots,
$\tbox\bu\bv\bw=\tbox\bw\bu\bv$, because a cyclic permutation is two
transpositions and each transposition reverses the sign. Applying it to the
second and third lines so that all three entries are ordered as in
Definition~\ref{def:inv3}$_{2}$:
\[
   \tbox{\be_{2}}{\bA\be_{3}}{\bA\be_{1}}
   =\tbox{\bA\be_{1}}{\be_{2}}{\bA\be_{3}},\qquad
   \tbox{\be_{3}}{\bA\be_{1}}{\bA\be_{2}}
   =\tbox{\bA\be_{1}}{\bA\be_{2}}{\be_{3}} .
\]
Adding the three diagonal entries,
\[
   \tr\bA^{*}=\tbox{\bA\be_{1}}{\bA\be_{2}}{\be_{3}}
             +\tbox{\bA\be_{1}}{\be_{2}}{\bA\be_{3}}
             +\tbox{\be_{1}}{\bA\be_{2}}{\bA\be_{3}} .
\]
The right side is exactly Definition~\ref{def:inv3}$_{2}$ evaluated at
$(\ba,\bb,\bc)=(\be_{1},\be_{2},\be_{3})$, for which
$\tbox{\be_{1}}{\be_{2}}{\be_{3}}=1$. Hence $\tr\bA^{*}=I_{2}(\bA)$, which
is the assertion.

\medskip
\emph{(b) $I_{2}(\bA)=\tfrac12[(\tr\bA)^{2}-\tr(\bA^{2})]$.}
By (a) it suffices to contract $(\ast_1)$ on $m=i$:
\[
   I_{2}(\bA)=\tr\bA^{*}=A^{*}_{ii}
   =\tfrac12\,\eps_{ijk}\eps_{ipq}A_{jp}A_{kq}.
\]
Again shifting cyclically to expose the contraction pattern of
Theorem~\ref{thm:epsdelta},
$\eps_{ijk}\eps_{ipq}=\eps_{jki}\eps_{pqi}=\dd_{jp}\dd_{kq}
-\dd_{jq}\dd_{kp}$, so
\[
   I_{2}(\bA)=\tfrac12\big(\dd_{jp}\dd_{kq}A_{jp}A_{kq}
                          -\dd_{jq}\dd_{kp}A_{jp}A_{kq}\big).
\]
Take the two pieces one at a time. In the first, $\dd_{jp}$ forces $p=j$ and
$\dd_{kq}$ forces $q=k$, leaving
\[
   \dd_{jp}\dd_{kq}A_{jp}A_{kq}=A_{jj}A_{kk}=(\tr\bA)(\tr\bA)=(\tr\bA)^{2}.
\]
In the second, $\dd_{jq}$ forces $q=j$ and $\dd_{kp}$ forces $p=k$, leaving
\[
   \dd_{jq}\dd_{kp}A_{jp}A_{kq}=A_{jk}A_{kj}.
\]
This last expression is a trace: the components of $\bA^{2}=\bA\bA$ are
$(\bA^{2})_{jl}=A_{jk}A_{kl}$, so setting $l=j$ and summing gives
$(\bA^{2})_{jj}=A_{jk}A_{kj}=\tr(\bA^{2})$. Therefore
\[
   I_{2}(\bA)=\tfrac12\big[(\tr\bA)^{2}-\tr(\bA^{2})\big].
\]

\medskip
\emph{(c) $\bA^{t}\bA^{*}=(\det\bA)\I$.}
Two independent proofs are given, the first free of components.

\emph{Without coordinates.} Let $\ba,\bb,\bc$ be arbitrary and set
$\bv=\ba\times\bb$. Using the definition of the transpose,
$\ip{\bc}{\bA^{t}\bw}=\ip{\bA\bc}{\bw}$, and then the definition of the
cofactor,
\[
\begin{aligned}
   \ip{\bc}{\bA^{t}\bA^{*}\bv}
   &=\ip{\bA\bc}{\bA^{*}(\ba\times\bb)}\\
   &=\ip{\bA\bc}{\bA\ba\times\bA\bb}\\
   &=\tbox{\bA\bc}{\bA\ba}{\bA\bb}\\
   &=\tbox{\bA\ba}{\bA\bb}{\bA\bc},
\end{aligned}
\]
the last step being the cyclic shift again. By
Definition~\ref{def:inv3}$_{3}$ together with Theorem~\ref{thm:det}, which
identify $I_{3}$ with $\det$,
\[
   \tbox{\bA\ba}{\bA\bb}{\bA\bc}=(\det\bA)\tbox\ba\bb\bc
   =(\det\bA)\,\ip{\bc}{\ba\times\bb}=(\det\bA)\,\ip\bc\bv ,
\]
where the middle equality is one more cyclic shift,
$\tbox\ba\bb\bc=\tbox\bc\ba\bb=\ip{\bc}{\ba\times\bb}$. Comparing the two
displays,
\[
   \ip{\bc}{\big[\bA^{t}\bA^{*}-(\det\bA)\I\big]\bv}=0
   \qquad\text{for every }\bc .
\]
A vector orthogonal to every vector is the zero vector, so
$[\bA^{t}\bA^{*}-(\det\bA)\I]\bv=\bze$ for every $\bv$ of the form
$\ba\times\bb$, that is, by Step 0, for \emph{every} $\bv$. A tensor
sending every vector to $\bze$ is $\Ob$, so $\bA^{t}\bA^{*}=(\det\bA)\I$.

\emph{In components.} Start from $(\bA^{t})_{ik}=A_{ki}$ and $(\ast_1)$:
\[
   (\bA^{t}\bA^{*})_{ij}=(\bA^{t})_{ik}A^{*}_{kj}=A_{ki}A^{*}_{kj}
   =\tfrac12\,\eps_{kmn}\eps_{jpq}\,A_{ki}A_{mp}A_{nq}.
\]
The triple sum $\eps_{kmn}A_{ki}A_{mp}A_{nq}$ is precisely
Proposition~\ref{prop:detforms} read with the matrix $(A_{ij})$ and columns
$i,p,q$:
\[
   \eps_{kmn}A_{ki}A_{mp}A_{nq}=\eps_{ipq}\det\bA .
\]
(For completeness, here is why. The left side changes sign whenever two of
$i,p,q$ are interchanged: swapping $i$ and $p$ gives
$\eps_{kmn}A_{kp}A_{mi}A_{nq}$, and renaming the dummies $k\leftrightarrow
m$ turns this into $\eps_{mkn}A_{mp}A_{ki}A_{nq}
=-\eps_{kmn}A_{ki}A_{mp}A_{nq}$; the other two swaps are identical in form.
Being trilinear and alternating in the columns, it equals a fixed multiple
of $\eps_{ipq}$ by Lemma~\ref{lem:altform}, and evaluating at
$(i,p,q)=(1,2,3)$ identifies that multiple as
$\eps_{kmn}A_{k1}A_{m2}A_{n3}=\tbox{\bA\be_{1}}{\bA\be_{2}}{\bA\be_{3}}
=\det\bA$.)
Substituting and using $\eps_{jpq}\eps_{ipq}=2\dd_{ij}$ from
Theorem~\ref{thm:epsdelta},
\[
   (\bA^{t}\bA^{*})_{ij}=\tfrac12\,(\det\bA)\,\eps_{jpq}\eps_{ipq}
   =\tfrac12\,(\det\bA)\,2\dd_{ij}=(\det\bA)\,\dd_{ij},
\]
which is the component statement of $\bA^{t}\bA^{*}=(\det\bA)\I$.

\emph{The consequence.} Suppose now $\bA$ is invertible, equivalently
$\det\bA\neq0$ by Theorem~\ref{thm:det}. Then $\bA^{t}$ is invertible too,
and multiplying $\bA^{t}\bA^{*}=(\det\bA)\I$ on the left by $(\bA^{t})^{-1}$
gives
\[
   \bA^{*}=(\det\bA)\,(\bA^{t})^{-1}=(\det\bA)\,\bA^{-t},
\]
the last equality being $(\bA^{t})^{-1}=(\bA^{-1})^{t}$, proved in
Problem~18, which is what licenses the single symbol $\bA^{-t}$.

\medskip
\emph{(d) The contracted determinant formula.}
Take the trace of (c). On the right, $\tr[(\det\bA)\I]=(\det\bA)\dd_{ii}
=3\det\bA$, since $\dd_{ii}=1+1+1=3$. On the left, setting $j=i$ in the
component computation just performed,
\[
   \tr(\bA^{t}\bA^{*})=(\bA^{t}\bA^{*})_{ii}=A_{ki}A^{*}_{ki}
   =\tfrac12\,\eps_{kmn}\eps_{ipq}A_{ki}A_{mp}A_{nq}.
\]
Equating,
\[
   3\det\bA=\tfrac12\,\eps_{kmn}\eps_{ipq}A_{ki}A_{mp}A_{nq},
   \qquad\text{so}\qquad
   \det\bA=\tfrac16\,\eps_{kmn}\eps_{ipq}A_{ki}A_{mp}A_{nq}.
\]
Renaming the summation indices, $k\to i$, $m\to j$, $n\to k$ in the first
$\eps$ and $i\to l$, $p\to m$, $q\to n$ in the second, so that the three
factors read $A_{il},A_{jm},A_{kn}$,
\[
   \det\bA=\tfrac16\,\eps_{ijk}\eps_{lmn}A_{il}A_{jm}A_{kn},
\]
in agreement with Corollary~\ref{cor:det6}, which reached the same formula
by contracting Proposition~\ref{prop:detforms} directly. The two routes are
genuinely different: there the $\tfrac16$ came from
$\eps_{mnp}\eps_{mnp}=6$, here from $\tfrac12\cdot 3$, the $\tfrac12$ being
built into the cofactor and the $3$ being $\tr\I$.

\medskip
\emph{Checks.} Take $\bA=\I$ first. Then $\I\ba\times\I\bb=\ba\times\bb$, so
$\I^{*}=\I$ and $\tr\I^{*}=3$; part (b) gives
$\tfrac12[3^{2}-3]=\tfrac12\cdot 6=3$, and part (c) gives
$\I^{t}\I^{*}=\I=(\det\I)\I$ since $\det\I=1$. Part (d) gives
$\tfrac16\eps_{ijk}\eps_{lmn}\dd_{il}\dd_{jm}\dd_{kn}
=\tfrac16\eps_{ijk}\eps_{ijk}=\tfrac16\cdot 6=1$.

Take next $\bA=\alpha\,\be_{1}\otimes\be_{1}+\beta\,\be_{2}\otimes\be_{2}
+\gamma\,\be_{3}\otimes\be_{3}$, so $\bA\be_{1}=\alpha\be_{1}$,
$\bA\be_{2}=\beta\be_{2}$, $\bA\be_{3}=\gamma\be_{3}$. From the definition,
\[
\begin{aligned}
   \bA^{*}\be_{1}&=\bA\be_{2}\times\bA\be_{3}
    =\beta\gamma\,(\be_{2}\times\be_{3})=\beta\gamma\,\be_{1},\\
   \bA^{*}\be_{2}&=\bA\be_{3}\times\bA\be_{1}
    =\gamma\alpha\,(\be_{3}\times\be_{1})=\gamma\alpha\,\be_{2},\\
   \bA^{*}\be_{3}&=\bA\be_{1}\times\bA\be_{2}
    =\alpha\beta\,(\be_{1}\times\be_{2})=\alpha\beta\,\be_{3},
\end{aligned}
\]
so $\tr\bA^{*}=\beta\gamma+\gamma\alpha+\alpha\beta$. Part (b) predicts
$\tfrac12[(\alpha+\beta+\gamma)^{2}-(\alpha^{2}+\beta^{2}+\gamma^{2})]$, and
expanding the square,
$(\alpha+\beta+\gamma)^{2}=\alpha^{2}+\beta^{2}+\gamma^{2}
+2(\alpha\beta+\beta\gamma+\gamma\alpha)$, the squares cancel and the
$\tfrac12$ removes the $2$, leaving
$\alpha\beta+\beta\gamma+\gamma\alpha$, which agrees. Part (c) predicts
$\bA^{t}\bA^{*}=(\alpha\beta\gamma)\I$, and indeed $\bA^{t}=\bA$ here while
$\bA\bA^{*}\be_{1}=\beta\gamma\,\bA\be_{1}=\alpha\beta\gamma\,\be_{1}$, and
likewise $\bA\bA^{*}\be_{2}=\alpha\beta\gamma\,\be_{2}$ and
$\bA\bA^{*}\be_{3}=\alpha\beta\gamma\,\be_{3}$, while
$\det\bA=\alpha\beta\gamma$.
\end{solution}

\begin{remark}[what the cofactor is for]\label{rem:cofuse}
Part (c) is not a curiosity; it is the identity that makes surface
deformation computable. If $\ba\times\bb$ is the vector area of the
parallelogram on $\ba$ and $\bb$, then Definition~\ref{def:cof} says that
$\bA^{*}$ carries vector areas exactly as $\bA$ carries vectors, and (c)
converts that geometric statement into the working formula
$\bA^{*}=(\det\bA)\bA^{-t}$. Applied in Chapter~2 to the deformation
gradient $\bF$ with $J=\det\bF$, the same two lines become
$\bF^{*}=J\bF^{-t}$ and the Piola--Nanson formula
$\bn\,\dl a=\bF^{*}\bN\,\dl A$. Part (b) is the identity that makes
$I_{2}$ computable from $\bA$ alone, without ever forming $\bA^{*}$, and it
is what turns the characteristic polynomial of Problem~23 into something
one can evaluate.
\end{remark}

\begin{problem}[Characteristic polynomial]
Show $\det(\bA-\lambda\I)=-\lambda^3+I_1\lambda^2-I_2\lambda+I_3$.
\end{problem}
\begin{solution}
By Definition~\ref{def:inv3}$_3$ applied to the tensor $\bA-\lambda\I$,
\[
\begin{aligned}
   \det(\bA-\lambda\I)\,\tbox\ba\bb\bc
   &=\tbox{(\bA-\lambda\I)\ba}{(\bA-\lambda\I)\bb}{(\bA-\lambda\I)\bc}\\
   &=\tbox{\bA\ba-\lambda\ba}{\bA\bb-\lambda\bb}{\bA\bc-\lambda\bc}.
\end{aligned}
\]
Expand by linearity in each slot separately. Each slot contributes either
$\bA(\cdot)$ or $-\lambda(\cdot)$, so there are $2^{3}=8$ terms. Group
them by how many slots took the $-\lambda$:
\[
\begin{aligned}
   0\ \text{slots}:&\quad \tbox{\bA\ba}{\bA\bb}{\bA\bc}
     =I_{3}\tbox\ba\bb\bc,\\
   1\ \text{slot}:&\quad -\lambda\Big[
     \tbox{\ba}{\bA\bb}{\bA\bc}+\tbox{\bA\ba}{\bb}{\bA\bc}
     +\tbox{\bA\ba}{\bA\bb}{\bc}\Big]=-\lambda I_{2}\tbox\ba\bb\bc,\\
   2\ \text{slots}:&\quad \lambda^{2}\Big[
     \tbox{\bA\ba}{\bb}{\bc}+\tbox{\ba}{\bA\bb}{\bc}
     +\tbox{\ba}{\bb}{\bA\bc}\Big]=\lambda^{2}I_{1}\tbox\ba\bb\bc,\\
   3\ \text{slots}:&\quad (-\lambda)^{3}\tbox\ba\bb\bc
     =-\lambda^{3}\tbox\ba\bb\bc,
\end{aligned}
\]
the three groupings being exactly Definition~\ref{def:inv3}$_{3}$,
$_{2}$, $_{1}$ in turn. Adding and cancelling the non-zero
$\tbox\ba\bb\bc$,
\[
   \det(\bA-\lambda\I)=-\lambda^{3}+I_{1}(\bA)\lambda^{2}
   -I_{2}(\bA)\lambda+I_{3}(\bA).
\]
The alternating signs are simply the $(-\lambda)^{n}$ of the grouping, and
this is why the invariants enter \eqr{3.2} in that order: $I_{k}$ is the
coefficient produced by putting $\bA$ into exactly $k$ of the three slots.
\end{solution}

\begin{problem}[Field product rules]
For scalar fields $\varphi,\psi$ and vector fields $\bu,\bv$, show
(a) $\grad(\varphi+\psi)=\grad\varphi+\grad\psi$;
(b) $\dvg(\bu+\bv)=\dvg\bu+\dvg\bv$;
(c) $\dvg(\varphi\bv)=\ip{\grad\varphi}\bv+\varphi\dvg\bv$;
(d) $\crl(\grad\varphi)=\bze$;
(e) $\dvg(\crl\bv)=0$;
(f) $\grad(\varphi\bu)=\bu\otimes\grad\varphi+\varphi\grad\bu$;
(g) $\dvg(\bu\otimes\bv)=(\grad\bu)\bv+(\dvg\bv)\bu$.
\end{problem}
\begin{solution}
All follow from the Cartesian representations and the ordinary product and
equality-of-mixed-partials rules; write everything in components.

(a) $[\grad(\varphi+\psi)]_i=(\varphi+\psi)_{,i}=\varphi_{,i}+\psi_{,i}$.
(b) $\dvg(\bu+\bv)=(u_i+v_i)_{,i}=u_{i,i}+v_{i,i}$.
(c) $\dvg(\varphi\bv)=(\varphi v_i)_{,i}=\varphi_{,i}v_i+\varphi v_{i,i}
=\ip{\grad\varphi}\bv+\varphi\dvg\bv$.
(d) $[\crl(\grad\varphi)]_j=\eps_{kij}(\varphi_{,i})_{,k}
=\eps_{kij}\varphi_{,ik}$; since $\varphi_{,ik}=\varphi_{,ki}$ is symmetric
in $(i,k)$ while $\eps_{kij}$ is antisymmetric under $i\leftrightarrow k$,
the contraction vanishes.
(e) $\dvg(\crl\bv)=(\eps_{kij}v_{i,k})_{,j}=\eps_{kij}v_{i,kj}$; $v_{i,kj}$
is symmetric in $(k,j)$ and $\eps_{kij}$ antisymmetric under
$k\leftrightarrow j$, so the sum is $0$.
(f) $[\grad(\varphi\bu)]_{ij}=(\varphi u_i)_{,j}=u_i\varphi_{,j}
+\varphi u_{i,j}=(\bu\otimes\grad\varphi)_{ij}+\varphi(\grad\bu)_{ij}$.
(g) $[\dvg(\bu\otimes\bv)]_i=(u_iv_j)_{,j}=u_{i,j}v_j+u_iv_{j,j}
=[(\grad\bu)\bv]_i+(\dvg\bv)u_i$.
(In (d),(e) twice continuous differentiability is assumed, so mixed
partials commute.)

Parts (d) and (e) are the same statement twice: a symmetric array
contracted with an antisymmetric one vanishes, since
$A_{pq}S_{pq}=-A_{qp}S_{pq}=-A_{pq}S_{qp}=-A_{pq}S_{pq}$ after relabelling,
so the quantity equals minus itself. The same argument proves the Piola
identity in Problem 2 of Chapter 2.
\end{solution}

\begin{problem}[Gradient in spherical coordinates]
With $\bx=r\berho(\theta,\varphi)$, $\berho=\cos\varphi\,\ber(\theta)
+\sin\varphi\,\be_3$, $\beph=\berho\times\beth$, find the components of
$\grad f$ in the variable basis $\{\berho,\beth,\beph\}$.
\end{problem}
\begin{solution}
\emph{Step 1: build the frame from the Cartesian basis.}
Nothing is assumed about spherical coordinates; everything is constructed.
The angle $\theta$ is the azimuth, measured in the equatorial plane from
$\be_{1}$, and $\varphi$ is the elevation above that plane, so $\varphi$ is
latitude and not the colatitude used in much of the physics literature.
Start with the unit vector pointing along the projection of $\bx$ onto the
equatorial plane,
\[
   \ber(\theta)=\cos\theta\,\be_{1}+\sin\theta\,\be_{2},
\]
and differentiate it with respect to $\theta$, the $\be_{i}$ being fixed:
\[
\begin{aligned}
   \beth(\theta)&:=\ber'(\theta)=-\sin\theta\,\be_{1}+\cos\theta\,\be_{2},\\
   \beth'(\theta)&=-\cos\theta\,\be_{1}-\sin\theta\,\be_{2}=-\ber(\theta).
\end{aligned}
\]
Both are unit vectors, since
$\cos^{2}\theta+\sin^{2}\theta=1$ in each case, and they are orthogonal:
\[
   \ip\ber\beth=\cos\theta(-\sin\theta)+\sin\theta(\cos\theta)=0 .
\]
Also $\ip\ber{\be_{3}}=0$ and $\ip\beth{\be_{3}}=0$, both vectors having no
$\be_{3}$ component. Tilting $\ber$ up out of the equatorial plane by the
angle $\varphi$ produces the radial direction given in the statement,
\[
   \berho(\theta,\varphi)=\cos\varphi\,\ber(\theta)+\sin\varphi\,\be_{3}
   =\cos\varphi\cos\theta\,\be_{1}+\cos\varphi\sin\theta\,\be_{2}
    +\sin\varphi\,\be_{3},
\]
so that $\bx=r\berho$ has Cartesian coordinates
$x_{1}=r\cos\varphi\cos\theta$, $x_{2}=r\cos\varphi\sin\theta$,
$x_{3}=r\sin\varphi$, and $|\bx|=r$ as the notation intends.

\emph{Step 2: the third leg of the frame, computed rather than quoted.}
The statement defines $\beph=\berho\times\beth$. Using
$(\bu\times\bv)_{1}=u_{2}v_{3}-u_{3}v_{2}$ and its two companions from
Proposition~\ref{prop:crosscomp}, with
\[
   \berho=(\cos\varphi\cos\theta,\ \cos\varphi\sin\theta,\ \sin\varphi),
   \qquad
   \beth=(-\sin\theta,\ \cos\theta,\ 0),
\]
one obtains
\[
\begin{aligned}
   (\berho\times\beth)_{1}
   &=(\cos\varphi\sin\theta)(0)-(\sin\varphi)(\cos\theta)\\
   &=-\sin\varphi\cos\theta,\\
   (\berho\times\beth)_{2}
   &=(\sin\varphi)(-\sin\theta)-(\cos\varphi\cos\theta)(0)\\
   &=-\sin\varphi\sin\theta,\\
   (\berho\times\beth)_{3}
   &=(\cos\varphi\cos\theta)(\cos\theta)-(\cos\varphi\sin\theta)(-\sin\theta)\\
   &=\cos\varphi(\cos^{2}\theta+\sin^{2}\theta)=\cos\varphi .
\end{aligned}
\]
Collecting, and recognising $\cos\theta\,\be_{1}+\sin\theta\,\be_{2}=\ber$,
\[
   \beph=-\sin\varphi\cos\theta\,\be_{1}-\sin\varphi\sin\theta\,\be_{2}
         +\cos\varphi\,\be_{3}
        =-\sin\varphi\,\ber(\theta)+\cos\varphi\,\be_{3}.
\]

\emph{Step 3: the frame is orthonormal and right-handed.}
Each of the six conditions is verified, using
$\ip\ber\ber=1$, $\ip\beth\beth=1$, $\ip\ber\beth=0$,
$\ip\ber{\be_{3}}=\ip\beth{\be_{3}}=0$, $\ip{\be_{3}}{\be_{3}}=1$:
\[
\begin{aligned}
   \ip\berho\berho&=\cos^{2}\varphi\,\ip\ber\ber
     +2\cos\varphi\sin\varphi\,\ip\ber{\be_{3}}
     +\sin^{2}\varphi\,\ip{\be_{3}}{\be_{3}}\\
     &=\cos^{2}\varphi+\sin^{2}\varphi=1,\\
   \ip\beph\beph&=\sin^{2}\varphi\,\ip\ber\ber
     -2\sin\varphi\cos\varphi\,\ip\ber{\be_{3}}
     +\cos^{2}\varphi\,\ip{\be_{3}}{\be_{3}}\\
     &=\sin^{2}\varphi+\cos^{2}\varphi=1,\\
   \ip\beth\beth&=1,\\
   \ip\berho\beth&=\cos\varphi\,\ip\ber\beth
     +\sin\varphi\,\ip{\be_{3}}\beth=0+0=0,\\
   \ip\beph\beth&=-\sin\varphi\,\ip\ber\beth
     +\cos\varphi\,\ip{\be_{3}}\beth=0+0=0,\\
   \ip\berho\beph&=(\cos\varphi)(-\sin\varphi)\ip\ber\ber
     +(\sin\varphi)(\cos\varphi)\ip{\be_{3}}{\be_{3}}\\
     &=-\cos\varphi\sin\varphi+\sin\varphi\cos\varphi=0 .
\end{aligned}
\]
For the handedness, use
$\bu\times(\bv\times\bw)=\ip\bu\bw\,\bv-\ip\bu\bv\,\bw$ from
Problem~11(a) with $\bu=\bw=\beth$ and $\bv=\berho$:
\[
   \beth\times\beph=\beth\times(\berho\times\beth)
   =\ip\beth\beth\,\berho-\ip\beth\berho\,\beth
   =1\cdot\berho-0\cdot\beth=\berho,
\]
so $\tbox\berho\beth\beph=\ip\berho{\beth\times\beph}
=\ip\berho\berho=1>0$. The triad is a right-handed orthonormal basis at
every point where it is defined.

\emph{Step 4: how the frame turns.}
Differentiate $\berho$ with respect to each angle, remembering that
$\be_{3}$ is constant and $\ber'=\beth$:
\[
   \frac{\partial\berho}{\partial\theta}
   =\cos\varphi\,\ber'(\theta)+0=\cos\varphi\,\beth,
   \qquad
   \frac{\partial\berho}{\partial\varphi}
   =-\sin\varphi\,\ber+\cos\varphi\,\be_{3}=\beph .
\]
The second is exactly the vector computed in Step 2, which is why the
problem defines $\beph$ by a cross product: the cross product and the
$\varphi$ derivative deliver the same direction, so the frame that is
natural geometrically is also the frame that is natural for differentiation.
Hence, treating $r,\theta,\varphi$ as independent,
\[
   \dl\berho=\frac{\partial\berho}{\partial\theta}\,\dl\theta
            +\frac{\partial\berho}{\partial\varphi}\,\dl\varphi
            =\cos\varphi\,\dl\theta\ \beth+\dl\varphi\ \beph .
\]

\emph{Step 5: the line element.}
Differentiate $\bx=r\berho(\theta,\varphi)$ by the product rule:
\begin{equation}
   \dl\bx=\dl r\,\berho+r\,\dl\berho
         =\underbrace{1}_{h_{r}}\,\dl r\ \berho
          +\underbrace{r\cos\varphi}_{h_{\theta}}\,\dl\theta\ \beth
          +\underbrace{r}_{h_{\varphi}}\,\dl\varphi\ \beph .
   \tag{$\ast_2$}
\end{equation}
The three coefficients $h_{r}=1$, $h_{\theta}=r\cos\varphi$, $h_{\varphi}=r$
are the only content of the problem. They are lengths per unit coordinate:
increasing $\theta$ by $\dl\theta$ slides the point along the circle of
latitude, whose radius is not $r$ but $r\cos\varphi$ (Figure~\ref{fig:sph}),
so the distance travelled is $r\cos\varphi\,\dl\theta$; increasing $\varphi$
by $\dl\varphi$ slides it along a meridian of radius $r$, a distance
$r\,\dl\varphi$; increasing $r$ moves it radially a distance $\dl r$,
whence $h_{r}=1$.

\begin{figure}[htbp]\centering
\begin{tikzpicture}[x={(-0.62cm,-0.42cm)},y={(1cm,-0.28cm)},z={(0cm,1cm)},
  scale=2.55,
  vec/.style={-{Stealth[length=2.4mm]},thick},
  lbl/.style={fill=white,inner sep=1pt,font=\small}]
\def\LO{112}   % longitude of the point
\def\LA{30}    % latitude  of the point
% polar axis and the reference direction for the azimuth
\draw[gray!55,thin] (0,0,-1.25) -- (0,0,1.40);
\node[lbl,text=gray!80] at (0,0,1.52) {$\be_{3}$};
\draw[gray!45,thin] (0,0,0) -- (1.22,0,0);
\node[lbl,text=gray!80] at (1.34,0,0) {$\be_{1}$};
% equator and the meridian through the point
\draw[gray!45,thin] plot[domain=0:360,samples=120] ({cos(\x)},{sin(\x)},0);
\draw[gray!45,thin] plot[domain=0:360,samples=120]
      ({cos(\LO)*cos(\x)},{sin(\LO)*cos(\x)},{sin(\x)});
% the parallel of latitude, drawn heavy: this is the circle whose radius
% is r cos(phi), not r
\draw[black,thick] plot[domain=0:360,samples=120]
      ({cos(\LA)*cos(\x)},{cos(\LA)*sin(\x)},{sin(\LA)});
% point, position vector, and the distance from the polar axis
\coordinate (O) at (0,0,0);
\coordinate (P) at ({cos(\LA)*cos(\LO)},{cos(\LA)*sin(\LO)},{sin(\LA)});
\coordinate (Q) at (0,0,{sin(\LA)});
\draw[vec] (O) -- (P);
\draw[gray!75,densely dashed] (Q) -- (P);
\fill (P) circle (0.4pt);  \fill (O) circle (0.4pt);
\node[lbl] at ({0.32*cos(\LA)*cos(\LO)},{0.32*cos(\LA)*sin(\LO)},
               {sin(\LA)+0.16}) {$r\cos\varphi$};
\node[lbl] at ({0.45*cos(\LA)*cos(\LO)},{0.45*cos(\LA)*sin(\LO)},
               {0.45*sin(\LA)-0.13}) {$r$};
% the local orthonormal triad
\draw[vec] (P) -- ++({0.46*cos(\LA)*cos(\LO)},{0.46*cos(\LA)*sin(\LO)},
                     {0.46*sin(\LA)}) node[lbl,right] {$\berho$};
\draw[vec] (P) -- ++({-0.46*sin(\LO)},{0.46*cos(\LO)},0)
                     node[lbl,above right=-1pt] {$\beth$};
\draw[vec] (P) -- ++({-0.46*sin(\LA)*cos(\LO)},{-0.46*sin(\LA)*sin(\LO)},
                     {0.46*cos(\LA)}) node[lbl,above left=-1pt] {$\beph$};
% the two angles
\draw[gray!80,thin] plot[domain=0:\LO,samples=40]
      ({0.42*cos(\x)},{0.42*sin(\x)},0);
\node[lbl,text=gray!90] at ({0.52*cos(0.5*\LO)},{0.52*sin(0.5*\LO)},0)
     {$\theta$};
\draw[gray!80,thin] plot[domain=0:\LA,samples=25]
      ({0.62*cos(\LO)*cos(\x)},{0.62*sin(\LO)*cos(\x)},{0.62*sin(\x)});
\node[lbl,text=gray!90] at ({0.76*cos(\LO)*cos(0.5*\LA)},
      {0.76*sin(\LO)*cos(0.5*\LA)},{0.76*sin(0.5*\LA)}) {$\varphi$};
\end{tikzpicture}
\hspace{0.35cm}
\raisebox{1.7cm}{\begin{tikzpicture}[scale=0.92]
\draw[thick] (0,0) rectangle (2.5,1.55);
\draw[thick] (2.5,1.55) -- (3.2,2.10) -- (3.2,0.55) -- (2.5,0);
\draw[thick] (0,1.55) -- (0.7,2.10) -- (3.2,2.10);
\draw[gray!55,densely dashed] (0,0) -- (0.7,0.55) -- (3.2,0.55);
\draw[gray!55,densely dashed] (0.7,0.55) -- (0.7,2.10);
\node[below,font=\small] at (1.25,0) {$r\cos\varphi\,\dl\theta$};
\node[right,font=\small] at (3.2,1.30) {$r\,\dl\varphi$};
\node[above left,font=\small] at (0.42,1.87) {$\dl r$};
\node[font=\small] at (1.55,0.85) {$\dl v$};
\end{tikzpicture}}
\caption{Left: the spherical frame at $\bx=r\berho$. The heavy circle is the
parallel of latitude $\varphi$, and its radius is the dashed distance
$r\cos\varphi$ from the polar axis, not $r$; this single fact is the whole
reason $h_{\theta}=r\cos\varphi$. The triad $\{\berho,\beth,\beph\}$ points
outward, east and north, and $\beth$ appears short only because it is
foreshortened, pointing towards the reader. Right: the box built on the three
coordinate increments, with edge lengths $\dl r$, $r\cos\varphi\,\dl\theta$
and $r\,\dl\varphi$; each edge is a coordinate increment multiplied by its
scale factor.}
\label{fig:sph}
\end{figure}

\begin{remark}[the three scale factors, read off the picture]
\label{rem:sphscale}
Each coefficient in $(\ast_{2})$ is the distance moved per unit
coordinate:
\[
\begin{aligned}
   \dl r\ \text{at fixed }\theta,\varphi&:\quad
     \text{outward along }\berho,\ \text{distance }\dl r,
     &&h_{r}=1,\\
   \dl\theta\ \text{at fixed }r,\varphi&:\quad
     \text{along the parallel, radius }\varrho=r\cos\varphi,
     &&h_{\theta}=r\cos\varphi,\\
   \dl\varphi\ \text{at fixed }r,\theta&:\quad
     \text{along the meridian, radius }r,
     &&h_{\varphi}=r .
\end{aligned}
\]
The middle one is the trap: the parallel has radius $\varrho=r\cos\varphi$,
the distance to the \emph{axis}, not $r$. Two checks. At the equator
$\cos\varphi=1$ and $h_{\theta}=r$. At a pole $\cos\varphi\to0$ and
$h_{\theta}\to0$, the degeneracy that excludes the poles below.
\end{remark}

\begin{figure}[htbp]\centering
\begin{tikzpicture}[scale=1.78,
  vec/.style={-{Stealth[length=2.4mm]},thick},
  lbl/.style={fill=white,inner sep=1.2pt,font=\small}]
% meridional section; radius 2.2, latitude phi = 38 deg, so
% P = 2.2(cos38, sin38) = (1.734, 1.355)
\draw[gray!45,thin] (-2.75,0)--(2.75,0);
\draw[gray!45,thin] (0,-2.60)--(0,2.45);
\node[lbl,text=gray!85,font=\small] at (0,2.62) {polar axis};
\node[lbl,text=gray!85,font=\small] at (2.90,0.22) {equator};
\draw[thick] (0,0) circle (2.2);
\draw[thick,black] (33:2.2) arc (33:43:2.2);
\fill (0,0) circle (1.0pt);
\coordinate (P) at (1.734,1.355);
\fill (P) circle (1.2pt);
% the radius r
\draw[vec] (0,0)--(P);
\node[lbl] at (1.47,0.80) {$r$};
% the latitude angle, measured up from the equator
\draw[gray!80,thin] (0:0.75) arc (0:38:0.75);
\node[lbl,font=\small] at (0.87,0.30) {$\varphi$};
% the distance from P to the POLAR AXIS
\draw[guide] (P)--(0,1.355);
\node[lbl] at (0.78,1.355) {$\varrho=r\cos\varphi$};
% the outward and northward directions at P
\draw[vec] (P)--++(0.473,0.369);
\node[lbl,right] at (2.26,1.76) {$\berho$};
\draw[vec] (P)--++(-0.369,0.473);
\node[lbl] at (1.24,2.02) {$\beph$};
% the azimuthal direction comes out of the page at P
\draw[thick] (P) circle (0.10);
\fill (P) circle (0.028);
\node[lbl] at (2.18,1.10) {$\beth$};
% the two arc lengths, as a legend clear of the circle
\node[lbl,align=center,font=\small] at (-1.55,-2.80)
  {along the meridian\\ (drawn heavy): $r\,\dl\varphi$};
\node[lbl,align=center,font=\small] at (1.55,-2.80)
  {around the parallel\\ (out of the page): $r\cos\varphi\,\dl\theta$};
\end{tikzpicture}
\caption{A meridional section, showing why $h_{\theta}=r\cos\varphi$ while
$h_{\varphi}=r$. Since $\varphi$ is measured from the equator, the point
is at distance $r$ from the origin but only $\varrho=r\cos\varphi$ from
the polar axis. A change of longitude sweeps it around that smaller
circle, a change of latitude along a meridian of radius $r$. The
azimuthal direction $\beth$ comes out of the page, marked by the circled
dot.}
\label{fig:sphscale}
\end{figure}

\emph{Step 6: invert the line element.}
Because $\{\berho,\beth,\beph\}$ is orthonormal, the three coefficients in
$(\ast_2)$ are recovered by dotting with the corresponding leg, every cross
term vanishing:
\[
\begin{aligned}
   \ip\berho{\dl\bx}&=\dl r\,\ip\berho\berho
     +r\cos\varphi\,\dl\theta\,\ip\berho\beth
     +r\,\dl\varphi\,\ip\berho\beph=\dl r,\\
   \ip\beth{\dl\bx}&=\dl r\,\ip\beth\berho
     +r\cos\varphi\,\dl\theta\,\ip\beth\beth
     +r\,\dl\varphi\,\ip\beth\beph=r\cos\varphi\,\dl\theta,\\
   \ip\beph{\dl\bx}&=\dl r\,\ip\beph\berho
     +r\cos\varphi\,\dl\theta\,\ip\beph\beth
     +r\,\dl\varphi\,\ip\beph\beph=r\,\dl\varphi .
\end{aligned}
\]
Solving for the coordinate increments, which requires $r\neq0$ and
$\cos\varphi\neq0$,
\[
   \dl r=\ip\berho{\dl\bx},\qquad
   \dl\theta=\frac{1}{r\cos\varphi}\,\ip\beth{\dl\bx},\qquad
   \dl\varphi=\frac{1}{r}\,\ip\beph{\dl\bx}.
\]

\emph{Step 7: the chain rule and the identification.}
Write $f=\hat f(r,\theta,\varphi)$ for the composite that assigns to the
coordinates the value of the field at the corresponding point. The ordinary
chain rule in three variables gives
\[
   \dl f=f_{,r}\,\dl r+f_{,\theta}\,\dl\theta+f_{,\varphi}\,\dl\varphi,
   \qquad f_{,r}:=\frac{\partial\hat f}{\partial r},\ \text{etc.}
\]
Substituting the three expressions just obtained and collecting the common
factor $\dl\bx$ on the right of every dot product,
\[
\begin{aligned}
   \dl f&=f_{,r}\,\ip\berho{\dl\bx}
        +f_{,\theta}\,\frac{1}{r\cos\varphi}\,\ip\beth{\dl\bx}
        +f_{,\varphi}\,\frac{1}{r}\,\ip\beph{\dl\bx}\\
   &=\Big(f_{,r}\,\berho+\frac{1}{r\cos\varphi}\,f_{,\theta}\,\beth
        +\frac{1}{r}\,f_{,\varphi}\,\beph\Big)\cdot\dl\bx .
\end{aligned}
\]
Compare with \eqr{1.114}, $\dl f=\ip{\grad f}{\dl\bx}$. Subtracting,
$\ip{\bw}{\dl\bx}=0$ where $\bw$ is the difference of $\grad f$ and the
bracket. As $(\dl r,\dl\theta,\dl\varphi)$ ranges over $\Rn^{3}$ the vector
$\dl\bx$ in $(\ast_2)$ ranges over all of $\Vt$, precisely because the three
scale factors are non-zero and the triad is a basis; so $\bw$ is orthogonal
to every vector and therefore $\bw=\bze$. Hence
\begin{equation}
   \boxed{\ \grad f=f_{,r}\,\berho
     +\frac{1}{r\cos\varphi}\,f_{,\theta}\,\beth
     +\frac{1}{r}\,f_{,\varphi}\,\beph\ }
   \tag{$\ast_3$}
\end{equation}
and the components sought are
$(\grad f)_{\rho}=f_{,r}$,
$(\grad f)_{\theta}=f_{,\theta}/(r\cos\varphi)$,
$(\grad f)_{\varphi}=f_{,\varphi}/r$.

\emph{Where it fails.} The derivation used $r\neq0$ and $\cos\varphi\neq0$.
The first excludes the origin, where no radial direction exists; the second
excludes the two poles $\varphi=\pm\pi/2$, where the circle of latitude
degenerates to a point, $\beth$ is undefined, and longitude carries no
information. These are failures of the coordinate system, not of the field:
$\grad f$ itself is perfectly well defined there, as \eqr{1.114} makes plain,
since it never mentions coordinates.

\emph{Checks.} Three independent ones, each computed both ways.

Take $f=r=|\bx|$. Then $f_{,r}=1$ and $f_{,\theta}=f_{,\varphi}=0$, so
$(\ast_3)$ gives $\grad r=\berho$. Independently, from
$r^{2}=\ip\bx\bx$ one has $2r\,\dl r=2\ip\bx{\dl\bx}$, hence
$\dl r=\ip{(\bx/r)}{\dl\bx}$ and $\grad r=\bx/r=\berho$, in agreement.

Take $f=x_{3}=r\sin\varphi$. Then $f_{,r}=\sin\varphi$, $f_{,\theta}=0$,
$f_{,\varphi}=r\cos\varphi$, so
\[
   \grad f=\sin\varphi\,\berho+\frac{1}{r}\,r\cos\varphi\,\beph
          =\sin\varphi\,\berho+\cos\varphi\,\beph .
\]
Substituting the Cartesian forms of $\berho$ and $\beph$,
\[
\begin{aligned}
   \grad f&=\sin\varphi\big(\cos\varphi\,\ber+\sin\varphi\,\be_{3}\big)
           +\cos\varphi\big(-\sin\varphi\,\ber+\cos\varphi\,\be_{3}\big)\\
   &=(\sin\varphi\cos\varphi-\cos\varphi\sin\varphi)\,\ber
     +(\sin^{2}\varphi+\cos^{2}\varphi)\,\be_{3}=\be_{3},
\end{aligned}
\]
which is right, since $\grad x_{3}=\be_{3}$ directly from
$\dl x_{3}=\ip{\be_{3}}{\dl\bx}$.

Take $f=\theta$. Then $(\ast_3)$ gives
$\grad\theta=\beth/(r\cos\varphi)$. This is the familiar polar result
$\grad\theta=\beth/\varrho$ with $\varrho=r\cos\varphi$ the distance from
the polar axis, which is exactly what the dashed segment in
Figure~\ref{fig:sph} measures, and it correctly blows up on the axis, where
$\theta$ is discontinuous.
\end{solution}

\begin{remark}[the scale-factor recipe, and the colatitude convention]%
\label{rem:lame}
Nothing in Steps 5 to 7 used the sphere. Whenever coordinates
$(q_{1},q_{2},q_{3})$ produce a line element of the orthogonal form
$\dl\bx=\sum_{k}h_{k}\,\dl q_{k}\,\bff_{k}$ with $\{\bff_{k}\}$ orthonormal,
the same three steps give
\[
   \grad f=\sum_{k}\frac{1}{h_{k}}\frac{\partial f}{\partial q_{k}}\,\bff_{k},
\]
the $h_{k}$ being called scale factors. Cartesian coordinates have
$h_{1}=h_{2}=h_{3}=1$ and the formula returns $f_{,i}\be_{i}$; the polar
coordinates of \S1.2.5 have $h_{r}=1$, $h_{\theta}=r$, $h_{z}=1$ and it
returns $f_{,r}\ber+r^{-1}f_{,\theta}\beth+f_{,z}\be_{3}$. The division by
$h_{k}$ is not a convention but a units statement: $\partial f/\partial
q_{k}$ is a change per unit \emph{coordinate}, while a gradient component
must be a change per unit \emph{length}, and $h_{k}$ is the number of
lengths per coordinate.

One warning on conventions. Most physics texts measure the polar angle
$\vartheta$ down from the north pole, the colatitude, and obtain
$h_{\vartheta}=r$, $h_{\theta}=r\sin\vartheta$. Here $\varphi$ is measured
up from the equator, so $\varphi=\pi/2-\vartheta$ and
$\sin\vartheta=\cos\varphi$, which is why $\cos\varphi$ appears in
$(\ast_3)$ where the familiar formula shows a sine. The two agree; only the
angle being named has moved.
\end{remark}

\begin{problem}[Laplace's equation between two spheres]
For $f=f(r)$, $r=|\bx|$, solve $\dvg(\grad f)=0$ on $a\le r\le b$ with
$f(a)=f_a$, $f(b)=f_b$.
\end{problem}
\begin{solution}
By Problem~25 with $f_{,\theta}=f_{,\varphi}=0$, $\grad f=f'(r)\berho$. Put
$\bv=\grad f=f'(r)\berho$; then, using $\dl\berho$ as above,
\[
\dl\bv=f''(r)\,\dl r\,\berho+f'(r)\big(\cos\varphi\,\beth\,\dl\theta
+\beph\,\dl\varphi\big),
\]
and inserting $\dl r=\berho\cdot\dl\bx$,
$\dl\theta=\tfrac{1}{r\cos\varphi}\beth\cdot\dl\bx$,
$\dl\varphi=\tfrac1r\beph\cdot\dl\bx$ and
$(\be_a\cdot\dl\bx)\be_b=(\be_b\otimes\be_a)\dl\bx$,
\[
\grad\bv=f''\,\berho\otimes\berho+\frac{f'}{r}\,\beth\otimes\beth
+\frac{f'}{r}\,\beph\otimes\beph,\qquad
\dvg\bv=\tr(\grad\bv)=f''+\frac{2f'}{r}.
\]
So Laplace's equation is $f''+2f'/r=0$, i.e. $(r^2f')'=0$, giving
$r^2f'=$ const, $f'=\alpha/r^2$, $f(r)=\beta-\alpha/r$. Rewriting the
constants through the boundary conditions,
\[
\boxed{\,f(r)=f_a+(f_b-f_a)\,\frac{\tfrac1a-\tfrac1r}{\tfrac1a-\tfrac1b}\,}
\]
which indeed gives $f(a)=f_a$, $f(b)=f_b$.
\end{solution}

\begin{problem}[A curl identity and the curl of a polar field]
Establish $\ip\bu{\crl\bv}=\dvg(\bv\times\bu)+\ip\bv{\crl\bu}$ and derive
the curl of the field \eqr{1.145}'s underlying
$\bu=u_r\ber+u_\theta\beth+u_z\bk$.
\end{problem}
\begin{solution}
\emph{Identity.} Write $(\crl\bv)_i=\eps_{ijk}v_{k,j}$. Then
\[
\dvg(\bv\times\bu)=\eps_{ijk}(v_ju_k)_{,i}
=\eps_{ijk}v_{j,i}u_k+\eps_{ijk}v_ju_{k,i}.
\]
In the first term $\eps_{ijk}v_{j,i}=\eps_{kij}v_{j,i}=(\crl\bv)_k$, so it
is $\ip\bu{\crl\bv}$; in the second $\eps_{ijk}u_{k,i}
=-\eps_{jik}u_{k,i}=-(\crl\bu)_j$, so it is $-\ip\bv{\crl\bu}$. Hence
$\dvg(\bv\times\bu)=\ip\bu{\crl\bv}-\ip\bv{\crl\bu}$, which rearranges to
the stated identity.

\emph{Curl of the polar field.} For any field, $\crl\bu$ is the axial
vector of $\grad\bu-(\grad\bu)^t$ (that is,
$(\grad\bu-(\grad\bu)^t)\ba=(\crl\bu)\times\ba$). From \eqr{1.145},
\[
\begin{aligned}
\grad\bu-(\grad\bu)^t
&=c\,(\ber\otimes\beth-\beth\otimes\ber)
+u_{z,r}(\bk\otimes\ber-\ber\otimes\bk)\\
&\quad+\tfrac1r u_{z,\theta}(\bk\otimes\beth-\beth\otimes\bk),
\end{aligned}
\]
with $c=\tfrac1r(u_{r,\theta}-u_\theta)-u_{\theta,r}$. The axial vector of
$\be_a\otimes\be_b-\be_b\otimes\be_a$ is $\be_b\times\be_a$, so
\[
\begin{aligned}
\crl\bu&=c\,(\beth\times\ber)+u_{z,r}(\ber\times\bk)
+\tfrac1r u_{z,\theta}(\beth\times\bk)\\
&=\tfrac1r u_{z,\theta}\,\ber-u_{z,r}\,\beth
+\Big[u_{\theta,r}+\tfrac1r(u_\theta-u_{r,\theta})\Big]\bk,
\end{aligned}
\]
using $\beth\times\ber=-\bk$, $\ber\times\bk=-\beth$,
$\beth\times\bk=\ber$.
\end{solution}

\begin{problem}[Cartesian formula for the curl of a tensor]
Establish $\crl\bA=\eps_{ilm}A_{jl,i}\,\be_m\otimes\be_j$.
\end{problem}
\begin{solution}
By definition $(\crl\bA)\bc=\crl(\bA^{t}\bc)$ for fixed $\bc$. Take
$\bc=\be_j$; then $\bA^{t}\be_j=A_{jl}\be_l$ is a vector field with $l$th
component $A_{jl}$, and with $(\crl\bw)_m=\eps_{miq}w_{q,i}$,
\[
\crl(\bA^{t}\be_j)=\eps_{miq}A_{jq,i}\,\be_m
=\eps_{ilm}A_{jl,i}\,\be_m
\]
(renaming $q\to l$ and using $\eps_{miq}=\eps_{qmi}$, i.e.\ writing the
antisymmetric symbol in the cyclic order $ilm$). Writing
$\crl\bA=C_{mj}\be_m\otimes\be_j$ and matching
$(\crl\bA)\be_j=C_{mj}\be_m$ gives $C_{mj}=\eps_{ilm}A_{jl,i}$, i.e.
$\crl\bA=\eps_{ilm}A_{jl,i}\,\be_m\otimes\be_j$.
\end{solution}

%=======================================================================
\appendix
\chapter{Dictionary to the book}\label{app:dict}
%=======================================================================

\renewcommand{\arraystretch}{1.2}
\begin{center}\footnotesize
\begin{tabular}{@{}ll@{\hspace{1.3em}}ll@{}}
\hline
Book & Here & Book & Here\\
\hline
(1.1)--(1.3) & Def.~\ref{def:lindep}, Thm.~\ref{thm:coord}
 & (1.79)--(1.84) & Thm.~\ref{thm:trace}\\
ip axioms & Def.~\ref{def:ip}
 & (1.85), (1.86) & Def.~\ref{def:cof}, Prop.~\ref{prop:I2cof}\\
(1.4)--(1.6) & Thm.~\ref{thm:cs}--Cor.~\ref{cor:tri}
 & (1.87)--(1.93) & Def.~\ref{def:tip}, Thm.~\ref{thm:tip}\\
(1.7)--(1.16) & Def.~\ref{def:onb}--Prop.~\ref{prop:comp}
 & (1.95) & Prop.~\ref{prop:posvec}\\
(1.17)--(1.19) & Thm.~\ref{thm:riesz}
 & \eqr{1.96}, (1.97) & Def.~\ref{def:diff}, Def.~\ref{def:landau}\\
(1.20)--(1.22) & Thm.~\ref{thm:cob}
 & (1.98), \eqr{1.99} & Thm.~\ref{thm:grad}\\
(1.23)--(1.30) & Def.~\ref{def:cross}--Prop.~\ref{prop:crosscomp}
 & \eqr{1.100}--(1.108) & Thm.~\ref{thm:dir}--Thm.~\ref{thm:cart}\\
(1.31)--(1.39) & Def.~\ref{def:det}--Cor.~\ref{cor:det6}
 & \eqr{1.109}--(1.115) & Thm.~\ref{thm:chain}--Ex.~\ref{ex:cartrecover}\\
(1.40)--(1.50) & Def.~\ref{def:tensor}--Ex.~\ref{ex:tp}
 & \eqr{1.116}--\eqr{1.125} & Ex.~\ref{ex:polar}, Ex.~\ref{ex:gradr}\\
(1.55)--(1.58) & Thm.~\ref{thm:tenstrans}
 & \eqr{1.126}--(1.134) & Def.~\ref{def:diffv}--Cor.~\ref{cor:dv}\\
(1.59)--(1.62) & Prop.~\ref{prop:transpose}, Prop.~\ref{prop:axial}
 & (1.135)--\eqr{1.145} & Ex.~\ref{ex:gu1}--Ex.~\ref{ex:gupolar}\\
(1.63)--(1.72) & Prop.~\ref{prop:prodinv}
 & (1.146)--(1.152) & Def.~\ref{def:divv}--Prop.~\ref{prop:divprod}\\
(1.73)--(1.77) & Prop.~\ref{prop:orth}
 & (1.153)--\eqr{1.162} & Ex.~\ref{ex:divpolar}\\
(1.78) & Def.~\ref{def:inv3}
 & (1.163)--(1.166) & Def.~\ref{def:curlv}--Rem.~\ref{rem:curlA}\\
sym/skw, sph/dev & Def.~\ref{def:splits}
 & \eqr{1.167}, \eqr{1.168} & Thm.~\ref{thm:div}, Thm.~\ref{thm:stokes}\\
\hline
\end{tabular}
\end{center}

\medskip\noindent
Lemma~\ref{lem:altform} (one dimensionality of alternating $3$ forms)
replaces the book's Appendix~A computations in
Propositions~\ref{prop:detforms},~\ref{prop:invwd} and
Theorem~\ref{thm:det}. All twenty eight chapter problems are solved in
\S\ref{sec:problems}.

\end{document}
